%%%%%%%%%%%%%%%%%%%%%%%%%%%%%%%%%%%%%%%%%%%%%%%%%%%%%%%%%%%%%%%%%%%%%%
% Lớp 12 | Nguyên hàm, tích phân | Ứng dụng thực tế và hình học của tích phân
% Dạng 2D4?3-3: Thể tích giới hạn bởi các đồ thị (tròn xoay)
% Số câu: 102  (sinh tự động bởi tach_cau.py; từ nay thầy sửa trực tiếp file này)
%%%%%%%%%%%%%%%%%%%%%%%%%%%%%%%%%%%%%%%%%%%%%%%%%%%%%%%%%%%%%%%%%%%%%%

%%%%%%%%%%%%%%%% TRẮC NGHIỆM 4 PHƯƠNG ÁN %%%%%%%%%%%%%%%%

% nguồn: Phan1_Lop_12_Dai_so_Chuong_4
\begin{ex}%[2D4N3-3]%[ma:aaf36382]
	Thể tích $V$ của khối tròn xoay được tạo  ra khi quay hình phẳng giới hạn bởi đồ thị hàm số $y=f(x)$, trục $Ox$ và hai đường thẳng $x=a$, $x=b$ $(a<b)$, xung quanh trục $Ox$ là
	\choice 
		{$V=\displaystyle\int_{a}^{b} |f(x)|\mathrm{\,d}{x}$} 
		 {\True $V=\pi\displaystyle\int_{a}^{b} f^2(x)\mathrm{\,d}{x}$}
		{$V=\displaystyle\int_{a}^{b} f^2(x)\mathrm{\,d}{x}$} 
		 {$V=\pi\displaystyle\int_{a}^{b} f(x)\mathrm{\,d}{x}$}
	 \loigiai{
		$V=\pi\displaystyle\int_{a}^{b}  f^2(x)\mathrm{\,d}{x}$.
	 }
\end{ex}

% nguồn: Phan1_Lop_12_Dai_so_Chuong_4
\begin{ex}%[2D4N3-3]%[ma:8593edbe]
	Cho hình phẳng $(H)$ giới hạn bởi các đường $y=x^2+6$, $y=0$, $x=0$,  $x=2$. Gọi $V$ là thể tích khối tròn xoay được tạo thành khi quay $(H)$ xung quanh trục $Ox$. Mệnh đề nào sau đây là đúng?
	\choice 
		 {\True $V=\pi\displaystyle\int_{0}^{2}\left(x^2+6\right)^2\mathrm{\,d}x$}
		 {$V=\displaystyle\int_{0}^{2}\left(x^2+6\right)\mathrm{\,d}x$}
		 {$V=\displaystyle\int_{0}^{2}\left(x^2+6\right)^2\mathrm{\,d}x$}
		{$V=\pi\displaystyle\int_{0}^{2}\left(x^2+6\right)\mathrm{\,d}x$} 
	\loigiai{ 
		 $V=\pi\displaystyle\int_{0}^{2}\left(x^2+6\right)^2\mathrm{\,d}x$.
	 }
\end{ex}

% nguồn: Phan1_Lop_12_Dai_so_Chuong_4
\begin{ex}%[2D4N3-3]%[ma:82ad427b]
	Cho hàm số $y=f(x)$ liên tục trên $[a;b]$. Gọi $D$ là hình phẳng giới hạn bởi đồ thị hàm số $y=f(x)$, trục hoành và hai đường thẳng $x=a$, $x=b$. Thể tích $V$ của vật thể  tròn xoay được tạo ra khi cho $D$ quay xung quanh trục $Ox$ là
	 \choice
		{$V=\pi^2 \displaystyle\int_a^b f^2(x)\mathrm{\,d}x$} 
		{$V=\displaystyle\int_a^b  f^2(x)\mathrm{\,d}x$}
		{$V=\pi\displaystyle\int_a^b  f(x)\mathrm{\,d}x$}
		 {\True $V=\pi\displaystyle\int_a^b f^2(x)\mathrm{\,d}x$}
	\loigiai{ 
		Ta có  $V=\pi\displaystyle\int_a^b f^2(x)\mathrm{\,d}x$.
	 }
\end{ex}

% nguồn: Phan1_Lop_12_Dai_so_Chuong_4
\begin{ex}%[2D4N3-3]%[ma:8baacad5]
	Cho hàm số $y = f(x)$ liên tục trên đoạn $[a;b]$. Gọi $D$ là hình phẳng giới hạn bởi đồ thị của hàm số $y = f(x)$, trục hoành và hai đường thẳng $x = a$, $x = b$ ($a < b$). Thể tích  của khối tròn xoay tạo thành khi quay $D$ quanh trục hoành được tính theo công thức nào?
	\choice 
		{\True $V =  \pi \displaystyle\int_a^b f^2(x) \mathrm{\,d}x$}
		 {$V = \pi \displaystyle\int_a^b f^2(x) \mathrm{\,d}x$}
		{$V =  \pi^2 \displaystyle\int_a^b f(x) \mathrm{\,d}x$}
		{$V = 2\pi  \displaystyle\int_a^b f^2(x) \mathrm{\,d}x$}
	 \loigiai{
		Theo  lý thuyết, ta có công thức $V = \pi \displaystyle\int_a^b f^2(x) \mathrm{\,d}x$.
	} 
\end{ex}

% nguồn: Phan1_Lop_12_Dai_so_Chuong_4
\begin{ex}%[2D4N3-3]%[ma:aa6f08d3]
	Gọi $H$ là hình phẳng giới hạn bởi đồ thị hàm số $y=\sqrt{x}$, trục hoành và hai đường thẳng $x=1$, $x=2$. Thể tích $V$ của khối tròn xoay tạo thành khi cho hình phẳng $H$ quay xung quanh trục  $Ox$ là
	 \choice
		{$V=\pi  \displaystyle\int_1^2 \sqrt{x} \mathrm{\,d} x$}
		{$V=\pi^2\displaystyle\int_1^2  x \mathrm{\,d} x$}
		{$V=\pi^2\displaystyle\int_1^2 \sqrt{x}  \mathrm{\,d} x$}
		{\True $V=\pi\displaystyle\int_1^2  x \mathrm{\,d} x$}
	\loigiai{ 
		Áp  dụng công thức, ta có \[V=\pi\displaystyle\int_1^2 \left(\sqrt{x}\right)^2 \mathrm{\,d} x=\pi\displaystyle\int_1^2 x \mathrm{\,d} x.\]
	} 
\end{ex}

% nguồn: Phan1_Lop_12_Dai_so_Chuong_4
\begin{ex}%[2D4N3-3]%[ma:c9e7a016]
	Hình phẳng giới hạn bởi đồ thị hàm số $y=2^{x}$, trục hoành và hai đường thẳng $x=0, x=7$ quay quanh trục $Ox$ tạo  thành khối tròn xoay có thể tích bằng
	\choice 
		 {$\pi\displaystyle\int_{0}^{7}2^{x}~\mathrm{\,d}x$}
		{$\displaystyle\int_{0}^{7}\left|2^{x}\right|\mathrm{\,d}x$} 
		{\True $\pi\displaystyle\int_{0}^{7}2^{2x}~\mathrm{\,d}x$} 
		 {$\displaystyle\int_{0}^{7}2^{2x}~\mathrm{\,d}x$}
	 \loigiai{
		Diện tích cần tính là  $S=\pi\displaystyle\int_{0}^{7}2^{2x}~\mathrm{\,d}x$.
	 }
\end{ex}

% nguồn: Phan1_Lop_12_Dai_so_Chuong_4
\begin{ex}%[2D4N3-3]%[ma:899eb38d]
	Cho hình phẳng $\left(S\right)$ giới hạn bởi đồ  thị hàm số $y=\sqrt{1-x^2}$, trục hoành và hai đường thẳng $x=-1$, $ x=1$. Thể tích của khối tròn xoay khi quay $\left(S\right)$ quanh $Ox$ là
	\choice 
		 {$\dfrac{3 \pi}{4}$}
		{$\dfrac{3 \pi}{2}$} 
		 {$\dfrac{2 \pi}{3}$}
		{\True  $\dfrac{4 \pi}{3}$}
	 \loigiai{
		 Ta có \[V=\pi \displaystyle\int_{-1}^{1} \left(\sqrt{1-x^2}\right)^2 \mathrm{\,d}x =\pi \displaystyle\int_{-1}^{1}\left(1-x^2\right) \mathrm{\,d}x=\dfrac{4 \pi}{3}.\]
	} 
\end{ex}

% nguồn: Phan1_Lop_12_Dai_so_Chuong_4
\begin{ex}%[2D4N3-3]%[ma:cee7b1ca]
	Cho hình phẳng $(H)$ giới hạn bởi đồ thị $y=2x-x^2$, trục hoành và hai đường thẳng $x=0$, $x=2$. Tính thể tích $V$ của khối tròn xoay sinh ra  khi cho $(H)$ quay quanh trục $Ox$.
	 \choice
		{$V=\pi  \displaystyle\int_0^2\left(2x-x^2\right)\mathrm{\,d}x$}
		{$V=\pi^2\displaystyle\int_0^2\left(2x-x^2\right)^2\mathrm{\,d}x$} 
		 {\True $V=\pi \displaystyle\int_0^2\left(2x-x^2\right)^2\mathrm{\,d}x$}
		{$V=\displaystyle\int_0^2\left(2x-x^2\right)^2 \mathrm{\,d}x$} 
	\loigiai{Thể  tích $V$ của khối tròn xoay sinh ra khi cho $(H)$ quay quanh trục $Ox$ là $V=\pi \displaystyle\int_0^2\left(2x-x^2\right)^2 \mathrm{\,d}x$.}
\end{ex}

% nguồn: Phan1_Lop_12_Dai_so_Chuong_4
\begin{ex}%[2D4N3-3]%[ma:7a6c649f]
	Thể tích khối tròn xoay tạo thành  khi cho hình phẳng giới hạn bởi đồ thị hàm số $f(x)=\sqrt{x}$, trục hoành và hai đường thẳng $x=0$, $x=2$ quay quanh trục $O x$ là
	 \choice
		{$\displaystyle\int_{0}^{2} x  \mathrm{\,d}x$}
		 {$\displaystyle\int_{0}^{2} \sqrt{x} \mathrm{\,d} x$}
		 {$\pi \displaystyle\int_{0}^{2} \sqrt{x} \mathrm{\,d} x$}
		{\True $\pi  \displaystyle\int_{0}^{2} x \mathrm{\,d} x$}
	\loigiai{ 
		Áp dụng lý thuyết ta có thể  tích cần tìm là $V=\pi \displaystyle\int_{0}^{2} (f(x))^2 \mathrm{\,d} x=\pi \displaystyle\int_{0}^{2} x \mathrm{\,d} x$.
	} 
\end{ex}

% nguồn: Phan1_Lop_12_Dai_so_Chuong_4
\begin{ex}%[2D4N3-3]%[ma:0e3682a9]
	Thể tích của vật thể  tròn xoay tạo thành khi quay hình phẳng giới hạn bởi các đường $y=x^2-2x$, $y=0$, $x=0$, $x=1$ quanh trục $Ox$ là
	 \choice
		 {$\dfrac{7\pi}{8}$}
		 {$\dfrac{8\pi}{15}$}
		{$\dfrac{15\pi}{8}$} 
		 {$\dfrac{8\pi}{7}$}
	\loigiai{ 
		Ta có  $V=\pi\displaystyle\int_{0}^{1}\left(x^2-2x\right)^2\mathrm{\,d}x=\dfrac{8\pi}{15}$.
	} 
\end{ex}

% nguồn: Phan1_Lop_12_Dai_so_Chuong_4
\begin{ex}%[2D4N3-3]%[ma:b1431e93]
	Tính thể tích $V$ của khối tròn xoay tạo thành khi quay quanh trục hoành hình phẳng giới hạn bởi các đường $y=x^2-2x$, $y=0$, $x=0$ và $x=1$. 
	\choice 
		{\True $V=\dfrac{8\pi}{15}$} 
		 {$V=\dfrac{7\pi}{8}$}
		 {$V=\dfrac{8\pi}{7}$}
		{$V=\dfrac{15\pi}{8}$} 
	\loigiai{ 
		 $V=\pi \displaystyle\int_0^1\left(x^2-2x\right)^2\mathrm{\,d}x=\pi \displaystyle\int_0^1\left(x^4-4x^3+4x^2\right) \mathrm{\,d}x=\pi\left(\dfrac{x^5}{5}-x^4+\dfrac{4}{3} x^3\right)\Bigg|_0^1=\pi\left(\dfrac{1}{5}-1+\dfrac{4}{3}\right)=\dfrac{8\pi}{15}$.
	 }
\end{ex}

% nguồn: Phan1_Lop_12_Dai_so_Chuong_4
\begin{ex}%[2D4N3-3]%[ma:61099b43]
	Cho hàm số $y=f(x)$ liên tục trên đoạn $[a;b]$. Gọi $D$ là hình phẳng  giới hạn bởi đồ thị hàm số $y=f(x)$, trục hoành và hai đường thẳng $x=a$, $x=b$ $(a<b)$. Thể tích của khối tròn xoay tạo thành khi quay $D$ quanh trục hoành được tính theo công thức
	\choice 
		{$V=\pi^2\displaystyle\int_a^b [f(x)]^2\mathrm{\,d}x$} 
		{$V=2\pi\displaystyle\int_a^b  [f(x)]^2\mathrm{\,d}x$}
		{\True  $V=\pi\displaystyle\int_a^b [f(x)]^2\mathrm{\,d}x$}
		 {$V=\pi^2\displaystyle\int_a^b f(x)\mathrm{\,d}x$}
	 \loigiai{
	Thể tích  $V=\pi\displaystyle\int_a^b [f(x)]^2\mathrm{\,d}x$.
	} 
\end{ex}

% nguồn: Phan1_Lop_12_Dai_so_Chuong_4
\begin{ex}%[2D4H3-3]%[ma:8d874cc1]
	Gọi $(H)$ là hình phẳng giới hạn bởi đồ thị hàm số $y=x^2-5x+4$ và trục $Ox$. Tính thể tích của khối tròn xoay sinh ra khi quay hình  $(H)$ quanh trục $Ox$.
	\choice 
		{$\dfrac{9}{2}$} 
		{$\dfrac{81}{10}$} 
		{\True $\dfrac{81\pi}{10}$} 
		 {$\dfrac{9\pi}{2}$}
	 \loigiai{
		Ta có: $x^2-5x+4=0\Leftrightarrow \hoac{&x= 1\\  &x= 4.}$\\
		Thể tích của khối tròn xoay sinh ra khi quay hình $(H)$  quanh trục $Ox$ là
		\[V=\pi \displaystyle\int_1^4\left(x^2-5x+4\right)^2 \mathrm{d}x=\dfrac{81\pi}{10}.\] 
	 }
\end{ex}

% nguồn: Phan1_Lop_12_Dai_so_Chuong_4
\begin{ex}%[2D4H3-3]%[ma:69d2b8f7]
	Cho hình phẳng $D$ giới hạn bởi đồ thị hàm số $y = x + 5$, $x = 0$, $x = 2$. Thể tích khối tròn xoay tạo thành khi quay $D$ quanh trục  $Ox$ là bao nhiêu?
	\choice 
		{\True  $\dfrac{218\pi}{3}$}
		 {$\dfrac{218}{3}$}
		{$218\pi$} 
		{$\dfrac{2\pi}{6}$} 
	 \loigiai{
		Thể tích khối tròn xoay khi quay  $D$ quanh $Ox$ là $V = \pi \displaystyle\int_{0}^{2} (x+5)^2 \mathrm{\,d}x = \pi \displaystyle\int_{0}^{2} (x^2 + 10x + 25) \mathrm{\,d}x=\dfrac{218\pi}{3}$.
	} 
\end{ex}

% nguồn: Phan1_Lop_12_Dai_so_Chuong_4
\begin{ex}%[2D4H3-3]%[ma:43e0412e]
	Thể tích khối tròn xoay giới hạn  bởi các đường $y=2x-x^2$, $y=0$ khi quay quanh trục $Ox$ là:
	 \choice
		{$\dfrac{4\pi}{3}$} 
		{$\dfrac{13\pi}{15}$} 
		{$\dfrac{14\pi}{15}$} 
		{\True $\dfrac{16\pi}{15}$} 
		\begin{center}
			\begin{tikzpicture}[line cap=butt,line join=miter,>=stealth]
					 \tikzset{declare function={f (\x)=2*\x-(\x)^2; a=f (1);}}
					\begin{scope}
						\clip (-1,-2) rectangle  (4,2);
						\fill[gray!35] plot[domain=0:2]  (\x,{f (\x)});
						\draw[samples=100]  plot[domain=-1:4] (\x,{f (\x)});
						\draw[samples=100,dashed]  plot[domain=-1:4] (\x,{-f (\x)});
					\end{scope}
					\draw[->] (-2,0)--(3.5,0) 
						node[shift={(-100:7pt)},font=\normalsize]{$x$}; 
					 \draw[->] (0,-2)--(0,2.3)
						 node[shift={(180:7pt)},font=\normalsize]{$y$};
					 \draw (5pt,0) |- (0,5pt)
						 (0,0) node[shift={(215:9pt)},font=\normalsize]{$O$};
					\draw (1,a) arc (90:270:{0.25*a}  and a);
					\draw[dashed] (1,a) arc  (90:-90:{0.25*a} and a);
					\path  (1.2,a) node[above]{$y=2x-x^2$};
				\end{tikzpicture}
	\end{center}
	\loigiai{ 
		Phương trình hoành độ giao điểm là: $-x^2+2x=0\Leftrightarrow \hoac{&x= 0\\  &x= 2.}$\\
		Thể  tích khối tròn xoay là: $V=\pi \displaystyle\int_0^2\left(2x-x^2\right)^2\mathrm{d}x=\dfrac{16\pi}{15}$
	} 
\end{ex}

% nguồn: Phan1_Lop_12_Dai_so_Chuong_4
\begin{ex}%[2D4H3-3]%[ma:4bf87b81]
	Thể tích khối tròn xoay thu được khi quay  hình phẳng giới hạn bởi các đường $y=x^2-2x$, $y=0$ và $x=3$ quay quanh trục $Ox$ bằng
	\choice 
		{\True $\dfrac{18\pi}{5}$} 
		 {$\dfrac{8\pi}{3}$}
		{$\dfrac{18}{5}$} 
		{$\dfrac{8}{3}$} 
	\loigiai{ 
		Xét phương  trình $x^2-2x=0\Leftrightarrow \hoac{&x= 0\\ &x= 2.}$\\
		Khi đó $V=\pi \displaystyle\int_0^3\left(x^2-2x\right)^2\mathrm{d}x=\pi \displaystyle\int_0^3\left(x^4-4x^3+4x^2\right)\mathrm{d}x=\pi \left(\dfrac{x^5}{5}-x^4+\dfrac{4x^3}{3} \right)\left|\begin{array}{l} 3\\ 0
	\end{array}\right.=\dfrac{18\pi}{5}\cdot$
 }
\end{ex}

% nguồn: Phan1_Lop_12_Dai_so_Chuong_4
\begin{ex}%[2D4H3-3]%[ma:c08b0a96]
	Tính thể tích  khối tròn xoay thu được khi quay hình phẳng giới hạn bởi hai đường $y=-x^2+3x$ và $y=0$ xung quanh trục $Ox$.
	 \choice
		{$\dfrac{5\pi}{2}$} 
		{$\dfrac{27\pi}{10}$} 
		 {\True $\dfrac{81\pi}{10}$}
		 {$\dfrac{9\pi}{2}$}
	 \loigiai{
		Phương trình hoành độ giao điểm giữa hai  đường $y=-x^2+3x$ và $y=0$ là
		$-x^2+3x=0\Leftrightarrow \hoac{&x= 0\\ &x= 3.}$\\ 
		Do đó thể tích khối tròn xoay cần tính là  $V=\pi \displaystyle\int_0^3\left(-x^2+3x\right)^2\mathrm{d}x=\dfrac{81\pi}{10}$.
	 }
\end{ex}

% nguồn: Phan1_Lop_12_Dai_so_Chuong_4
\begin{ex}%[2D4H3-3]%[ma:e0a17e19]
	Cho hình phẳng $(H)$ giới hạn bởi đồ thị hàm số $y=4-x^2$ và trục hoành. Thể tích khối tròn xoay được tạo thành khi quay $(H)$ xung quanh  trục hoành bằng
	 \choice
		{$\dfrac{32\pi}{3}$} 
		 {\True $\dfrac{512\pi}{15}$}
		 {$\dfrac{32}{3}$}
		{$\dfrac{512}{15}$} 
	 \loigiai{
		Phương trình  hoành độ giao điểm: $4-x^2=0\Leftrightarrow x=\pm 2$.\\
		Khi đó: $V=\pi \displaystyle\int_{-2}^2\left(4-x^2\right)^2\mathrm{d}x=\pi \displaystyle\int_{-2}^2\left(16-8x^2+x^4\right)\mathrm{d}x=\pi \left.  \left(16x-\dfrac{8}{3} x^3+\dfrac{x^5}{5} \right)\right|_{-2}^2=\dfrac{512}{15} \pi$.
	 }
\end{ex}

% nguồn: Phan1_Lop_12_Dai_so_Chuong_4
\begin{ex}%[2D4H3-3]%[ma:c2bc253a]
	Thể tích khối  tròn xoay thu được khi quay hình phẳng giới hạn bởi hai đường $y=x^2-4x+3$ và $y=0$ quanh trục $Ox$ bằng
	\choice 
		 {$\dfrac{16}{15}$}
		{\True  $\dfrac{16\pi}{15}$}
		{$\dfrac{31\pi}{30}$} 
		 {$\dfrac{31}{30}$}
	 \loigiai{
		Phương trình hoành độ giao điểm  của đường $y=x^2-4x+3$ và đường $y=0$ là
		\[x^2-4x+3=0\Leftrightarrow \hoac{&x=  1\\ &x= 3.} \]
		Thể tích khối tròn xoay thu được khi quay hình phẳng giới hạn bởi hai đường $y=x^2-4x+3$ và  $y=0$ quanh trục $Ox$ là $V=\pi \displaystyle\int_1^3\left(x^2-4x+3\right)^2\mathrm{d}x=\dfrac{16\pi}{15}$.
	 }
\end{ex}

% nguồn: Phan1_Lop_12_Dai_so_Chuong_4
\begin{ex}%[2D4H3-3]%[ma:3eefc464]
	Thể tích của khối tròn xoay thu được khi quay hình phẳng giới hạn bởi hai đường $y=4-x^2$ và $y=0$  quanh trục $Ox$ bằng
	 \choice
		{\True $V=\dfrac{512  \pi}{15}$}
		{$V=\dfrac{32}{3}$} 
		 {$V=\dfrac{32 \pi}{3}$}
		 {$V=\dfrac{512}{15}$}
	\loigiai{ 
		Giao điểm của đồ thị  $y=4-x^2$ với trục $ Ox $ (tức là $ y=0 $).\\
		 $4-x^2=0 \Rightarrow x^2=4 \Rightarrow x=\pm 2$.\\
		Thể tích khối tròn xoay khi quay quanh trục $Ox$ là\\ 
		 $$V=\pi\displaystyle\int_{-2}^{2} (4-x^2)^2\,\mathrm{d}x=\dfrac{512 \pi}{15}.$$
	} 
\end{ex}

% nguồn: Phan1_Lop_12_Dai_so_Chuong_4
\begin{ex}%[2D4H3-3]%[ma:bf6aba32]
	Cho hình phẳng giới hạn bởi đồ thị hàm số $y=\sin x$,  trục $Ox$ và các đường thẳng $x=0$, $x=\pi$ quay quanh trục $Ox$. Tính thể tích khối tròn xoay tạo thành bằng
	\choice 
		{$\displaystyle\int_0^{\pi}\sin  x \mathrm{d}x$}
		 {$\displaystyle\int_0^{\pi}\sin ^2x \mathrm{d}x$}
		{\True $\pi  \displaystyle\int_0^{\pi}\sin ^2x \mathrm{d}x$}
		{$\pi  \displaystyle\int_0^{\pi}\cos ^2x \mathrm{d}x$}
	 \loigiai{
		Công thức tính thể tích khối  tròn xoay $V=\pi \displaystyle\int_a^bf^2(x)\mathrm{d}x$.
	} 
\end{ex}

% nguồn: Phan1_Lop_12_Dai_so_Chuong_4
\begin{ex}%[2D4H3-3]%[ma:20c76af1]
	Cho hàm số $y=f(x)$. Đồ thị của đạo hàm $f'(x)$ là đường cong trong Hình 2. Biết rằng diện tích của các phần hình phẳng $A$ và $B$ lần lượt là $S_A=2$ và  $S_B=3$. Nếu $f(0)=4$ thì giá trị của $f(5)$ bằng
	 \choice
		 {\True $3$}
		 {$5$}
		{$9$} 
		 {$-1$}
		\begin{center}
			\begin{tikzpicture}
					\def  \k{3/10}
					 \fill[orange!40](0,0)--plot[domain=0:2](\x,{\x*(\x-2)*(\x-5)*\k})--(2,0)--(0,0);
					\fill[blue!30](2,0)--plot[domain=2:5](\x,{\x*(\x-2)*(\x-5)*\k})--(5,0)--(2,0); 
					 \draw plot[samples=100,domain=-0.4:5.2](\x,{\x*(\x-2)*(\x-5)*\k}) node[above left]{$y=f'(x)$};
					\draw [->](-0.5,0)  node[below]{$O$}--(6,0) node[below]{$x$};
					\draw[->](0,-2.7)--(0,2.5)  node[left]{$y$};
					\foreach \i  in {1,2,3,4,5}{\draw ($(\i,0)+(90:0.1)$)--($(\i,0)+(-90:0.1)$);}
					\foreach \i in {-1,1,2}{\draw ($(0,\i)+(180:0.1)$)--($(0,\i)+(0:0.1)$);} 
					\foreach \i/\t/\g in {1/A/100,2/2/-90,4/B/-120,5/5/-80}{\path  ($(\i,0)+(\g:0.5)$) node{$\t$};}
				\end{tikzpicture}
	\end{center}
	\loigiai{ 
		 Ta có $\displaystyle\int_{0}^{5}f'(x)\mathrm{\,d}x=f(x)\big|_{0}^{5}=f(5)-f(0)$.\qquad(1)\\
		Mặt  khác, $\displaystyle\int_{0}^{5}f'(x)\mathrm{\,d}x=\displaystyle\int_{0}^{2}f'(x)\mathrm{\,d}x+\displaystyle\int_{2}^{5}f'(x)\mathrm{\,d}x$.\\
		Mà $f'(x)\ge0,\forall x\in[0;2]$ và $f'(x)\le  0,\forall x\in[2;5]$.\\
		 Do đó $\displaystyle\int_{0}^{5}f'(x)\mathrm{\,d}x=\displaystyle\int_{0}^{2}f'(x)\mathrm{\,d}x+\displaystyle\int_{2}^{5}f'(x)\mathrm{\,d}x=S_A-S_B=2-3=-1$.\qquad (2)\\
		Từ $(1)$ và $(2)$ suy  ra $f(5)-f(0)=-1\Leftrightarrow f(5)=3$.
	} 
\end{ex}

% nguồn: Phan1_Lop_12_Dai_so_Chuong_4
\begin{ex}%[2D4H3-3]%[ma:8123a9c6]
	Cho hình $(H)$ giới hạn bởi các  đường $y=x^2-2x$, trục hoành. Quay hình phẳng $(H)$ quanh trục $Ox$ ta được khối tròn xoay có thể tích là
	\choice 
		 {$\dfrac{496\pi}{15}$}
		{$\dfrac{32\pi}{15}$} 
		{$\dfrac{4\pi}{15}$} 
		 {\True $\dfrac{16\pi}{15}$}
	 \loigiai{
		 Hoành đồ giao điểm của đồ thị hàm số $y=x^2-2x$ và trục hoành $x^2-2x=0\Leftrightarrow \hoac{x=0\\x=2.}$\\
		Thể tích khối  tròn xoay cần tìm là
		$$V=\pi  \displaystyle\int_0^2 (x^2-2x)^2\mathrm{\,d}x=\pi \displaystyle\int_0^2 \left(x^4-4x^3+4x^2\right)\mathrm{\,d}x=\pi\left(\dfrac{x^5}{5}-x^4+\dfrac{4}{3}x^3\right)\Bigg|_0^2=\dfrac{16\pi}{15}.$$
	} 
\end{ex}

% nguồn: Phan1_Lop_12_Dai_so_Chuong_4
\begin{ex}%[2D4H3-3]%[ma:cc46d219]
	Thể  tích của khối tròn xoay do hình phẳng giới hạn bởi đồ thị các hàm số $y=-x^2+3x$ và $y=0$ khi quay quanh trục $Ox$ bằng
	 \choice
		{$\dfrac{81}{10}$} 
		 {$\dfrac{9}{2}$}
		 {$\dfrac{9\pi}{2}$}
		{\True  $\dfrac{81\pi}{10}$}
	 \loigiai{
		Phương trình hoành  độ giao điểm $-x^2+3x=0\Leftrightarrow\hoac{&x=0\\
		 &x=3.}$\\
		Ta có  $V=\pi \displaystyle\int_0^3\left(-x^2+3x\right)^2\mathrm{\,d} x=\dfrac{81\pi}{10}$ (đvtt).
	} 
\end{ex}

% nguồn: Phan1_Lop_12_Dai_so_Chuong_4
\begin{ex}%[2D4H3-3]%[ma:ea230f15]
	Kí hiệu $(H)$ là hình phẳng giới hạn bở đồ thị hàm số $y=f(x)=\sqrt{x}.e^{x^2}$, trục hoành, đường  thẳng $x=1$. Tính thể tích $V$ của khối tròn xoay thu được khi quay $(H)$ quanh trục hoành.
	 \choice
		{\True $V=\dfrac{1}{4} \pi \left(e^2-1\right)$} 
		 {$V=\pi \left(e^2-1\right)$}
		{$V=\dfrac{1}{4}  \pi e^2-1$}
		 {$V=e^2-1$}
	\loigiai{ 
		Ta có: $\sqrt{x}\cdot e^{x^2}=0\Leftrightarrow x=0$.\\ 
		Thể tích khối tròn xoay thu được khi quay $(H)$ quanh trục hoành là: 
		\[V=\pi \displaystyle\int_0^1\left(\sqrt{x}.e^{x^2} \right) ^2\mathrm{d}x=\pi \displaystyle\int_0^1xe^{2x^2} \mathrm{d}x=\dfrac{\pi}{4} \displaystyle\int_0^1e^{2x^2} \mathrm{d}(2x^2)=\left.  \dfrac{\pi}{4} e^{2x^2} \right|_0^1=\dfrac{\pi}{4} \left(e^2-1\right).\]
	} 
\end{ex}

% nguồn: Phan1_Lop_12_Dai_so_Chuong_4
\begin{ex}%[2D4H3-3]%[ma:b9d7772e]
	Thể tích khối tròn xoay khi cho hình phẳng giới hạn bởi các  đường $y=x^2+1$, $x=0$, $x=1$, $y=0$ quay quanh trục $Ox$ là:
	 \choice
		{$\dfrac{4\pi}{3}$} 
		 {$\dfrac{28}{15}$}
		{\True $\dfrac{28\pi}{15}$} 
		 {$\dfrac{4}{3}$}
		\begin{center}
			\begin{tikzpicture}[line cap=butt,line join=miter,>=stealth]
					\tikzset{declare function={f (\x)=(\x)^2+1; xa=1;  ya=f (xa);}}
					\begin{scope}
						 \clip (-1,-3) rectangle (2,3);
						\fill[gray!35] plot[domain=0:1.05] (\x,{f  (\x)})--(1.05,0)--(0,0)--cycle;
						\draw[samples=100] plot[domain=0:3] (\x,{f (\x)}); 
						 \draw[samples=100,dashed] plot[domain=0:3] (\x,{-f (\x)});
						\draw  (1.05,0)--(1.05,ya);
					\end{scope}
					\draw[->] (-1,0)--(2.5,0) 
						 node[shift={(-100:7pt)},font=\normalsize]{$x$};
					 \draw[->] (0,-3)--(0,3.5)
						 node[shift={(180:7pt)},font=\normalsize]{$y$};
					 \draw (5pt,0) |- (0,5pt)
						 (0,0) node[shift={(240:9pt)},font=\normalsize]{$O$};
					\draw[densely dashed] (0,0) ellipse ({0.2} and 1) 
						({xa+0.05},0)  ellipse ({0.15*ya} and {ya});
					\path  (xa+0.1,ya) node[above right]{$y=x^2+1$};
				\end{tikzpicture}
	\end{center}
	 \loigiai{
		Thể tích khối tròn xoay là: $V=\pi \displaystyle\int_0^1\left(x^2+1\right)^2\mathrm{d}x=\dfrac{28\pi}{15}\cdot$ 
	} 
\end{ex}

% nguồn: Phan1_Lop_12_Dai_so_Chuong_4
\begin{ex}%[2D4H3-3]%[ma:3591f532]
	Cho hình phẳng $D$ giới hạn bởi các đường $y=\sqrt{x-1}$, trục hoành và $x=5$. Thể tích khối tròn xoay tạo thành khi quay $D$ quanh trục $Ox$  bằng
	\choice 
		{$\dfrac{15\pi}{2}$} 
		{$\dfrac{15}{2}$} 
		 {\True $8\pi$}
		 {$8$}
	\loigiai{ 
		Phương  trình hoành độ giao điểm của hàm số $y=\sqrt{x-1}$ và trục hoành là $\sqrt{x-1}=0\Rightarrow x=1$.\\
		Khi đó thể tích  khối tròn xoay tạo thành khi quay $D$ quanh trục $Ox$ là
		\[V=\pi \displaystyle\int_1^5\left(\sqrt{x-1} \right)^2\mathrm{d}x=\pi \displaystyle\int_1^5(x-1)  \; \mathrm{d}x=\pi \left. \left(\dfrac{x^2}{2}-x\right)\right|_1^5=8\pi.\]
	 }
\end{ex}

% nguồn: Phan1_Lop_12_Dai_so_Chuong_4
\begin{ex}%[2D4H3-3]%[ma:f3bcf69b]
	Tính thể tích khối tròn xoay được tạo thành khi quay hình phẳng giới hạn bởi đồ thị hàm số $y=3x-x^2$ và  trục hoành, quanh trục hoành.
	\choice 
		{\True $\dfrac{81\pi}{10}$ (đvtt)} 
		 {$\dfrac{85\pi}{10}$ (đvtt)}
		 {$\dfrac{41\pi}{7}$ (đvtt)}
		 {$\dfrac{8\pi}{7}$ (đvtt)}
	 \loigiai{
		 Ta có $3x-x^2=0\Leftrightarrow\hoac{&x=0\\&x=3.}$\\
		Thể tích khối  tròn xoay cần tìm là
		\begin{align*}
			V=\pi \displaystyle\int_0^3\left(3x-x^2\right)^2\mathrm{\,d}x=\pi  \displaystyle\int_0^3\left(9x^2-6x^3+x^4\right) \mathrm{\,d}x=\left.\pi\left(3x^3-\dfrac{3x^4}{2}+\dfrac{x^5}{5}\right)\right|_0^3=\dfrac{81\pi}{10}\text{ (đvtt)}.
		\end{align*}
	 }
\end{ex}

% nguồn: Phan1_Lop_12_Dai_so_Chuong_4
\begin{ex}%[2D4H3-3]%[ma:23868e5d]
	Cho hình phẳng  $(H)$ giới hạn bởi đồ thị hàm số $y=e^x$, các đường thẳng $x=0$, $x=\ln 3$ và trục hoành. Thể tích khối tròn xoay sinh bởi $(H)$ khi quay quanh trục hoành là
	 \choice
		 {$2\pi$}
		{\True $4\pi$} 
		{$4$} 
		{$\pi$} 
	\loigiai{ 
		\[V=\pi \displaystyle\int_0^{\ln 3}e^{2x} \; \mathrm{d}x=\pi  \dfrac{1}{2} e^{2x}|_0^{\ln 3}=4\pi.\]
	 }
\end{ex}

% nguồn: Phan1_Lop_12_Dai_so_Chuong_4
\begin{ex}%[2D4H3-3]%[ma:2f7d58a3]
	Cho $(H)$ là hình phẳng giới hạn bởi đồ  thị hàm số $y=x^2-x-2$ và trục hoành. Quay hình phẳng $(H)$ quay quanh trục hoành, ta được một khối nón tròn xoay có thể tích bằng
	\choice 
		{\True $\dfrac{81}{10}  \pi$}
		{$\dfrac{81}{10}$} 
		 {$\dfrac{9}{2}$}
		{$\dfrac{9\pi}{2}$} 
	\loigiai{ 
		Phương trình hoành độ giao điểm của đồ thị với trục hoành: $x^2-x-2=0\Leftrightarrow 
		 \hoac{&x= -1\\ &x= 2.}$\\
		Thể tích khối nón là  $V=\pi \displaystyle\int_{-1}^2\left(x^2-x-2\right)^2\mathrm{d}x=\dfrac{81\pi}{10}$.
	} 
\end{ex}

% nguồn: Phan1_Lop_12_Dai_so_Chuong_4
\begin{ex}%[2D4H3-3]%[ma:2a55eb99]
	Thể tích khối tròn xoay do hình phẳng $(H)$ giới hạn bởi các đường $y=\tan x$, $y=0$, $x=0$, $x=\dfrac{\pi}{3}$ quay quanh  trục $Ox$ tạo thành là
	 \choice
		{\True  $\dfrac{\pi}{3}(3\sqrt{3}-\pi)$}
		 {$\dfrac{\pi(\sqrt{3}-1)}{3}$}
		{$\dfrac{\pi}{3}(3\sqrt{3}- \pi)$} 
		 {$\pi\sqrt{3}$}
	 \loigiai{
		Thể tích khối tròn xoay cần  tìm là
		 $$\pi\displaystyle\int_0^{\tfrac{\pi}{3}} \tan^2x\mathrm{\,d}x=\pi\displaystyle\int_0^{\tfrac{\pi}{3}} \left(\dfrac{1}{\cos^2x}-1\right)\mathrm{\,d}x=\pi(\tan x-x)\Bigg|_0^{\tfrac{\pi}{3}}=\dfrac{\pi(3\sqrt{3}-\pi)}{3}.$$
	 }
\end{ex}

% nguồn: Phan1_Lop_12_Dai_so_Chuong_4
\begin{ex}%[2D4H3-3]%[ma:2f51e7ca]
	Cho hình phẳng giới hạn bởi đồ thị hàm số $y=x^2$ và đường thẳng $y=4$ quay quanh trục $Ox$. Thể tích khối tròn xoay sinh ra bằng: 
	 \choice
		 {$\dfrac{64\pi}{5}$}
		{$\dfrac{128\pi}{5}$} 
		 {\True $\dfrac{256\pi}{5}$}
		{$\dfrac{152\pi}{5}$} 
		\begin{center}
			\resizebox{!}{5cm}{\begin{tikzpicture}[line cap=butt,line join=miter,>=stealth]
		\tikzset{declare function={f (\x)=(\x)^2; xa=2;  ya=f (xa);}}
		\begin{scope}
			 \clip (-3,-4.5) rectangle (3,5);
			\fill[gray!35] plot[domain=-2:2] (\x,{f (\x)}); 
			\draw[samples=100] plot[domain=-3:3] (\x,{f  (\x)});
			\draw[samples=100,dashed] plot[domain=-3:3] (\x,{-f  (\x)});
			 \draw (-3,4)--+(6,0);
			 \draw[dashed] (2.1,0)--(2.1,4)
				(-2.1,0)--(-2.1,4); 
		\end{scope}
		\draw[->] (-3,0)--(3.5,0) 
			 node[shift={(-100:7pt)},font=\normalsize]{$x$};
		\draw[->]  (0,-4.5)--(0,5)
			node[shift={(180:7pt)},font=\normalsize]{$y$}; 
		 \draw (5pt,0) |- (0,5pt)
			(0,0)  node[shift={(255:9pt)},font=\normalsize]{$O$};
		\draw[dashed] ({xa+0.08},0) ellipse ({0.1*ya} and {ya+0.02}) 
			({-xa-0.08},0) ellipse ({0.1*ya} and  {ya+0.02});
		\path (-0.6,4)  node[above]{$y=4$};
	\end{tikzpicture}}
	\end{center}
	 \loigiai{
		Phương trình hoành độ giao điểm  là: $x^2=4\Leftrightarrow x^2-4=0\Leftrightarrow x=\pm 2$.\\
		Thể tích khối tròn xoay là: $V=\pi  \displaystyle\int_{-2}^2\left|(x^2)^2-4^2\right|\mathrm{d}x=\dfrac{256\pi}{5}$.
	} 
\end{ex}

% nguồn: Phan1_Lop_12_Dai_so_Chuong_4
\begin{ex}%[2D4H3-3]%[ma:f83b0843]
	Thể tích khối tròn xoay thu được khi quay hình phẳng giới hạn bởi hai đường $y=\sqrt{4x-x^2}$ và $y=0$ quanh trục $Ox$ bằng 
	 \choice
		{$\dfrac{34\pi}{3}$} 
		{$\dfrac{31\pi}{3}$} 
		{\True $\dfrac{32\pi}{3}$} 
		 {$\dfrac{35\pi}{3}$}
	 \loigiai{
		Hoành độ  giao điểm của đồ thị hàm số $y=\sqrt{4x-x^2}$ và đường thẳng $y=0$ là nghiệm phương trình $\sqrt{4x-x^2}=0\Leftrightarrow
		\hoac{&x= 0 \\ &x= 4.}$\\ 
		Do đó, thể tích khối tròn xoay khi cho hình phẳng giới hạn bởi đồ thị hàm số  $y=\sqrt{4x-x^2}$ và $y=0$ quay quanh trục $Ox$ là $V=\pi \displaystyle\int_0^4\left(\sqrt{4x-x^2} \right)^2\mathrm{d}x=\pi \displaystyle\int_0^4\left(4x-x^2\right)\mathrm{d}x=\pi \left. \left(2x^2-\dfrac{1}{3} x^3\right)\right|_0^4=\dfrac{32\pi}{3}$.
	 }
\end{ex}

% nguồn: Phan1_Lop_12_Dai_so_Chuong_4
\begin{ex}%[2D4H3-3]%[ma:a1c91380]
	Cho hình phẳng $D$ giới hạn bởi đường cong $y=\sqrt{x^2+1}$, trục hoành và các đường thẳng $x=0$, $x=1$. Khối tròn xoay tạo thành khi quay  $D$ quanh trục hoành có thể tích $V$ bằng bao nhiêu?
	\choice 
		 {$V=2$}
		{\True  $V=\dfrac{4 \pi}{3}$}
		{$V=2  \pi$}
		 {$V=\dfrac{4}{3}$}
	\loigiai{Thể tích $V$ của hình phẳng $D$ giới hạn bởi đường cong $y=\sqrt{x^2+1}$,  trục hoành và các đường thẳng $x=0$, $x=1$ là
	$$V= \pi  \displaystyle\int_{0}^{1} \left(\sqrt{x^2+1}\right)^2\mathrm{\,d}x=\dfrac{4 \pi}{3}.$$}
\end{ex}

% nguồn: Phan1_Lop_12_Dai_so_Chuong_4
\begin{ex}%[2D4H3-3]%[ma:572e6099]
	 Cho miền phẳng $(D)$ giới hạn bởi đồ thị hàm số $y=\sqrt{x}$, hai đường thẳng $x=1$, $x=2$ và trục hoành. Tính thể tích khối tròn xoay tạo thành khi quay $(D)$ quanh trục hoành.
	 \choice
		 {$\dfrac{3}{2}$}
		 {\True $\dfrac{3\pi}{2}$}
		 {$\dfrac{2\pi}{3}$}
		{$3\pi$} 
	\loigiai{ 
		Tính thể tích khối tròn xoay tạo thành  khi quay $(D)$ quanh trục hoành $V=\pi \displaystyle\int_1^2\left(\sqrt{x} \right)^2\mathrm{d}x=\dfrac{3\pi}{2}$.
	} 
\end{ex}

% nguồn: Phan1_Lop_12_Dai_so_Chuong_4
\begin{ex}%[2D4H3-3]%[ma:43c541c0]
	Tính thể tích $V$ của phần vật thể giới hạn bởi hai mặt phẳng $x=0$, $x=1$, có thiết diện bị cắt bởi mặt phẳng vuông góc với trục $Ox$ tại điểm có hoành độ $x$,  $(0\le x\le 1)$ là một tam giác đều có cạnh bằng $x$.
	\choice 
		 {$V=\dfrac{12\pi}{5}$}
		{$V=\dfrac{12}{5}$} 
		{$V=\dfrac{\sqrt{3} \pi}{12}$} 
		 {\True $V=\dfrac{\sqrt{3}}{12}$}
	 \loigiai{
		Thể tích vật thể là: $V=\displaystyle\int_0^1S(x)\mathrm{d}x=\displaystyle\int_0^1x^2\cdot\dfrac{\sqrt{3}}{4} \mathrm{d}x=\dfrac{\sqrt{3}}{12}$. 
	} 
\end{ex}

% nguồn: Phan1_Lop_12_Dai_so_Chuong_4
\begin{ex}%[2D4H3-3]%[ma:b69c2ace]
	Thể tích  vật thể tròn xoay khi cho hình phẳng giới hạn bởi đồ thị hàm số $y=x^2-4x+3$ và trục hoành quay quanh $Ox$ là
	\choice 
		 {$\dfrac{4}{3}$}
		{$\dfrac{16}{15}$} 
		 {$\dfrac{4\pi}{3}$}
		{\True  $\dfrac{16\pi}{15}$}
	\loigiai{ 
		Ta có  $x^2-4x+3=0\Leftrightarrow \hoac{&x=1 \\ &x=3.}$\\
		Vậy $V_{Ox}=\pi  \displaystyle\int_1^3\left(x^2-4x+3\right)^2\mathrm{d}x=\dfrac{16\pi}{15}$.
	} 
\end{ex}

% nguồn: Phan1_Lop_12_Dai_so_Chuong_4
\begin{ex}%[2D4H3-3]%[ma:6cf27476]
	Thể tích khối tròn xoay do hình phẳng giới hạn  bởi các đường $y=\sqrt{x}$, trục $Ox$ và hai đường thẳng $x=1$;
	$x=4$ khi quay quanh trục hoành được tính bởi công thức  nào?
	\choice 
		{\True $V=\pi  \displaystyle\int_1^4x\mathrm{d}x$}
		 {$V=\pi ^2\displaystyle\int_1^4x\mathrm{d}x$}
		{$V=\pi \displaystyle\int_1^4\sqrt{x} \mathrm{d}x$} 
		 {$V=\displaystyle\int_1^4\left|\sqrt{x} \right|\mathrm{d}x$}
	\loigiai{ 
		 Ta có: $V=\pi \displaystyle\int_1^4x\mathrm{d}x$
	} 
\end{ex}

% nguồn: Phan1_Lop_12_Dai_so_Chuong_4
\begin{ex}%[2D4H3-3]%[ma:e2bb8ba8]
	Cho hình phẳng giới hạn bởi đường cong $y=\sqrt{x+2}$ và hai trục toạ độ. Thể tích khối tròn xoay  thu được khi quay hình phẳng đó quanh trục $Ox$ bằng
	\choice 
		 {$\dfrac{4\sqrt{2}}{3}$}
		{$2$} 
		 {\True $2\pi$}
		{$\dfrac{4\sqrt{2}}{3} \pi$} 
	 \loigiai{
		\[V=\pi  \displaystyle\int_{-2}^0\left(\sqrt{x+2} \right)^2\mathrm{d}x=\left. \pi \left(x^2+2x\right)\right|_{-2}^0=2\pi.\]
	 }
\end{ex}

% nguồn: Phan1_Lop_12_Dai_so_Chuong_4
\begin{ex}%[2D4H3-3]%[ma:c5352220]
	Cho hình phẳng $(H)$ giới hạn bởi đồ thị $(P)\colon y=2x-x^2$ và trục $Ox$. Tính thể tích của khối tròn xoay tạo  thành khi cho $(H)$ quay quanh trục $Ox$.
	 \choice
		 {$V=\dfrac{19\pi}{15}$}
		{$V=\dfrac{13\pi}{15}$} 
		{$V=\dfrac{17\pi}{15}$} 
		 {\True $V=\dfrac{16\pi}{15}$}
	 \loigiai{
		Phương trình hoành độ giao điểm của đồ  thị $(P)$ và trục $Ox$ là: $2x-x^2=0\Leftrightarrow\hoac{&x= 0\\ &x= 2.}$\\
		Thể tích  khối tròn xoay cần tìm là $V=\pi \displaystyle\int_0^2\left(2x-x^2\right)^2\mathrm{d}x=\dfrac{16\pi}{5}$.
	} 
\end{ex}

% nguồn: Phan1_Lop_12_Dai_so_Chuong_4
\begin{ex}%[2D4H3-3]%[ma:3169caa1]
	Thể tích của khối tròn xoay sinh ra khi quay hình phẳng giới hạn bởi các đường $y=5x-6x^2$, $y=0$, $x=0$,  $x=1$ xung quanh trục $Ox$.
	 \choice
		{$\dfrac{1}{2}$} 
		{$\dfrac{8}{15}$} 
		{$\dfrac{\pi}{2}$} 
		{\True $\dfrac{8\pi}{15}$} 
	\loigiai{Thể tích khối  tròn xoay sinh ra là
		\begin{eqnarray*}
			V&=&\pi\displaystyle\int_0^1\left(5x-6x^2\right)^2\mathrm{\,d}x=\pi\displaystyle\int_0^1\left(25x^2-60x^3+36x^4\right)\mathrm{\,d}x\\ 
			&=&\left(\dfrac{25x^3}{3}-15x^4+\dfrac{36x^5}{5}\right)\Bigg|_0^1=\dfrac{8\pi}{15}. 
	\end{eqnarray*}}
\end{ex}

% nguồn: Phan1_Lop_12_Dai_so_Chuong_4
\begin{ex}%[2D4H3-3]%[ma:9dfebf58]
	Cho hình phẳng $(H)$ được giới hạn bởi đồ thị của hàm số $y=x^2$ và đồ thị của hàm số $y=\sqrt{x}$ (tham  khảo hình vẽ). Thể tích của khối tròn xoay thu được khi quay $(H)$ quanh trục $Ox$ bằng
	\begin{center}
		\begin{tikzpicture}[line cap=butt,line join=miter,>=stealth]
				\tikzset{declare function={f (\x)=(\x)^2; g (\x)=sqrt  (\x);}}
				\begin{scope}[overlay]
					\clip  plot[domain=0:2] (\x,{g (\x)})--(2,0)--cycle;
					\path[fill=gray!45] plot[domain=-2:2]  (\x,{f (\x)});
				\end{scope}
				 \draw[->] (-2,0)--(3.5,0)
					 node[shift={(-100:7pt)},font=\normalsize]{$x$};
				\draw[->]  (0,-0.5)--(0,3)
					node[shift={(180:7pt)},font=\normalsize]{$y$}; 
				\draw (5pt,0) |- (0,5pt) 
					 (0,0) node[shift={(225:9pt)},font=\normalsize]{$O$};
				\begin{scope}
					\clip  (-2,-0.5) rectangle (3,3);
					\draw[samples=100]  plot[domain=-3:3] (\x,{f (\x)});
					 \draw[samples=100] plot[domain=0:3] (\x,{g (\x)});
				\end{scope}
				\draw[dashed] (1,0) node[below]{$1$} |- (0,1)  node[left]{$1$};
				\path (1.5,{g  (1.4)+0.3}) node[below right]{$y=\sqrt{x}$};
				\path (1.5,{f  (1.5)}) node[left]{$y=x^2$};
			\end{tikzpicture}
	\end{center}
	\choice 
		{$V=\dfrac{9\pi}{10}$} 
		{\True  $V=\dfrac{3\pi}{10}$}
		 {$V=\dfrac{\pi}{10}$}
		 {$V=\dfrac{7\pi}{10}$}
	\loigiai{ 
		Thể tích khối tròn xoay thu được là $V=\pi \displaystyle\int_0^1\left|x^4-x\right|\mathrm{d}x=\dfrac{3\pi}{10}$. 
	 }
\end{ex}

% nguồn: Phan1_Lop_12_Dai_so_Chuong_4
\begin{ex}%[2D4H3-3]%[ma:036bc5f3]
	Thể tích khối tròn xoay thu được khi quay  hình phẳng giới hạn bởi hai đường $y=\left|x^2-3x+2\right|$ và $y=0$ quanh trục $Ox$ bằng
	\choice 
		{$\dfrac{\pi  ^2}{30}$}
		{$\dfrac{\pi}{6}$} 
		{\True $\dfrac{\pi}{30}$} 
		{$\dfrac{\pi  ^2}{6}$}
	\loigiai{ 
		Phương trình hoành độ giao điểm $\left|x^2-3x+2\right|=0\Leftrightarrow  \hoac{&x= 1 \\ &x= 2.}$\\
		Vậy $V=\pi \displaystyle\int_1^2\left|x^2-3x+2\right|^2\mathrm{d}x=\dfrac{\pi}{30}$. 
	} 
\end{ex}

% nguồn: Phan1_Lop_12_Dai_so_Chuong_4
\begin{ex}%[2D4H3-3]%[ma:cffb2b5c]
	Cho hình phẳng $(H)$ giới hạn bởi đồ thị hàm số $y=\sqrt{x}$, hai đường thẳng $x=1$, $x=4$ và trục hoành. Tính thể tích khối tròn xoay  tạo thành khi quay $(H)$ quanh trục hoành.
	 \choice
		 {\True $\dfrac{15\pi}{2}$}
		 {$\dfrac{15}{2}$}
		{$\dfrac{14\pi}{3}$} 
		 {$\dfrac{14}{3}$}
	 \loigiai{
		Thể tích  khối tròn xoay tạo thành khi quay $(H)$ quanh trục hoành
		\[V=\pi \displaystyle\int_1^4\left(\sqrt{x} \right)^2\mathrm{d}x=\pi  \displaystyle\int_1^4x\mathrm{d}x=\pi\cdot\left. \dfrac{x^2}{2} \right|_1^4=\dfrac{15\pi}{2}\cdot\]
	 }
\end{ex}

% nguồn: Phan1_Lop_12_Dai_so_Chuong_4
\begin{ex}%[2D4H3-3]%[ma:4f675292]
	Thể tích khối tròn xoay thu được khi quay hình phẳng giới hạn bởi các đường $y=4-x^2$ và $y=0$ quanh trục  $Ox$ bằng
	 \choice
		 {$\dfrac{32\pi}{3}$}
		{\True  $\dfrac{512\pi}{15}$}
		{$\dfrac{16\pi}{3}$} 
		{$\dfrac{256\pi}{15}$} 
	 \loigiai{
		Xét phương trình hoành độ giao  điểm: $4-x^2=0\Leftrightarrow x=\pm 2$.\\
		Thể tích khối tròn xoay thu được là: $V=\pi \displaystyle\int_{-2}^2\left(4-x^2\right)^2\mathrm{d}x=\dfrac{512\pi}{15}$. 
	 }
\end{ex}

% nguồn: Phan1_Lop_12_Dai_so_Chuong_4
\begin{ex}%[2D4H3-3]%[ma:1f34cf94]
	Thể tích khối  tròn xoay thu được khi quay hình phẳng giới hạn bởi hai đường $y=x^2-x$ và $y=0$ quanh trục $Ox$ bằng
	 \choice
		{$\dfrac{\pi}{3}$} 
		{$\dfrac{\pi}{15}$} 
		{\True $\dfrac{\pi}{30}$} 
		 {$\dfrac{\pi}{5}$}
	\loigiai{ 
		Xét phương  trình hoành độ giao điểm: $x^2-x=0\Leftrightarrow \hoac{&x= 0\\ &x= 1.}$\\
		Thể tích khối  tròn xoay là: $V=\pi \displaystyle\int_0^1\left(x^2-x\right)^2 \mathrm{d}x=\dfrac{\pi}{30}$.
	} 
\end{ex}

% nguồn: Phan1_Lop_12_Dai_so_Chuong_4
\begin{ex}%[2D4H3-3]%[ma:33e368a9]
	Cho $D$ là  hình phẳng giới hạn bởi đồ thị hàm số $y=\sqrt{x+1}$, trục tung, trục hoành và đường thẳng $x=2$. Thể tích của khối tròn xoay tạo thành khi quay $D$ quanh trục hoành bằng
	 \choice
		{$6\pi$} 
		 {$2\pi$}
		{$3\pi$} 
		{\True$4\pi$} 
	\loigiai{ 
		Thể tích cần tìm là  $V=\pi\displaystyle\int_{0}^{2} (x+1)\mathrm{\,d}x=\pi\left(\dfrac{x^2}{2}+x\right)\bigg|_{0}^{2}=4\pi$.
	 }
\end{ex}

% nguồn: Phan1_Lop_12_Dai_so_Chuong_4
\begin{ex}%[2D4H3-3]%[ma:088d36d9]
	Gọi $(H)$ là hình phẳng giới hạn bởi đồ thị hàm số $y=x^2-3x+2$ và trục hoành. Thể tích của khối tròn xoay sinh ra khi quay hình  $(H)$ quanh trục hoành bằng
	\choice 
		{$\dfrac{\pi}{6}$} 
		 {$\dfrac{1}{6}$}
		 {\True $\dfrac{\pi}{30}$}
		{$\dfrac{1}{30}$} 
	 \loigiai{
		Phương  trình hoành độ giao điểm: $y=0\Leftrightarrow x^2-3x+2=0\Leftrightarrow
		\hoac{&x=  1\\ &x= 2.}$\\
		Vậy $V=\pi \displaystyle\int_1^2\left(x^2-3x+2\right)^2 \mathrm{d}x=\dfrac{\pi}{30}$. 
	 }
\end{ex}

% nguồn: Phan1_Lop_12_Dai_so_Chuong_4
\begin{ex}%[2D4H3-3]%[ma:471866d4]
	Tính thể tích $V$ của phần vật thể giới hạn  bởi hai mặt phẳng $x=0$ và $x=3$, biết rằng khi cắt vật thể bởi mặt phẳng tùy ý vuông góc với trục $Ox$ tại điểm có hoành độ $x$ ($0\le x\le 3$) thì được thiết diện là một hình vuông có độ dài cạnh bằng $2\sqrt{9-x^2}$
	\choice 
		 {$90$}
		 {$72\pi$}
		 {$78\pi$}
		 {\True $72$}
	 \loigiai{
		Diện tích hình vuông  là $S=\left(2\sqrt{9-x^2} \right)^2=4\left(9-x^2\right)=36-4x^2$.\\
		Vậy thể tích vật thể là $V=\displaystyle\int_0^3S(x) \mathrm{d}x=\displaystyle\int_0^3\left(36-4x^2\right)  \mathrm{d}x=72$.
	} 
\end{ex}

% nguồn: Phan1_Lop_12_Dai_so_Chuong_4
\begin{ex}%[2D4H3-3]%[ma:400c35a5]
	Giả sử $D$ là hình phẳng giới hạn bởi đường parabol $y=x^2-3x+2$ và trục hoành. Quay $D$ quanh trục hoành ta thu được khối tròn xoay có thể tích  bằng
	 \choice
		 {\True $V=\dfrac{\pi}{30}$}
		{$V=\dfrac{1}{30}$} 
		 {$V=\dfrac{1}{6}$}
		{$V=\dfrac{\pi}{6}$} 
	 \loigiai{
		Ta có $y=0\Leftrightarrow x^2-3x+2=0\Leftrightarrow \hoac{&x=1  \\ &x= 2.}$\\
		Khi đó:  $V=\pi \displaystyle\int_1^2\left(x^2-3x+2\right)^2\mathrm{d}x=\dfrac{\pi}{30}$.
	 }
\end{ex}

% nguồn: Phan1_Lop_12_Dai_so_Chuong_4
\begin{ex}%[2D4H3-3]%[ma:4bd0f26a]
	Cho hình phẳng $(H)$ được giới hạn bởi hai đồ thị $y=x^2+x$ và $y=2x$. Quay hình $(H)$ quanh trục hoành, tính thể tích vật  thể thu được
	 \choice
		 {\True $\dfrac{3\pi}{10}$}
		{$\dfrac{5\pi}{6}$} 
		{$\dfrac{\pi}{6}$} 
		{$\dfrac{5}{6}$} 
	\loigiai{ 
		Xét phương trình:  $x^2+x=2x\Leftrightarrow x^2-x=0\Leftrightarrow
		\hoac{&x= 0\\ &x= 1.}$\\ 
		Ta có thể tích cần tìm là:  $V=\pi \displaystyle\int_0^1\left|\left(x^2+x\right)^2-(2x)^2\right|\mathrm{d}x=\dfrac{3\pi}{10}$
	} 
\end{ex}

% nguồn: Phan1_Lop_12_Dai_so_Chuong_4
\begin{ex}%[2D4H3-3]%[ma:b5aba08b]
	Thể tích khối  tròn xoay khi quay hình phẳng $(H)$ xác định bởi các đường $y=\dfrac{1}{3} x^3-x^2$ và $y=0$ quanh trục $Ox$ là
	\choice 
		{$\dfrac{71\pi}{35}$} 
		{$\dfrac{81}{35}$} 
		{$\dfrac{71}{35}$} 
		 {\True $\dfrac{81\pi}{35}$}
	\loigiai{ 
		Ta có $\dfrac{1}{3}  x^3-x^2=0\Leftrightarrow \hoac{&x= 0\\ &x= 3.}$\\
		Vậy thể tích khối tròn xoay cần tìm là $V=\pi \displaystyle\int_0^3\left(\dfrac{1}{3} x^3-x^2\right)^2\mathrm{d}x=\dfrac{81\pi}{35}\cdot$ 
	 }
\end{ex}

% nguồn: Phan1_Lop_12_Dai_so_Chuong_4
\begin{ex}%[2D4H3-3]%[ma:b7a17778]
	Cho hình phẳng $D$ giới hạn bởi đường cong  $y=\sqrt{2+\sin x}$, trục hoành và các đường thẳng $x=0$, $x=\pi$. Khối tròn xoay tạo thành khi xoay $D$ quanh trục hoành có thể tích $V$ bằng bao nhiêu?
	 \choice
		{\True $V=2\pi(\pi+1)$} 
		 {$V=2(\pi+1)$}
		{$V=2\pi$} 
		{$V=2\pi^2$} 
	\loigiai{Thể tích khối tròn xoay khi quay $D$  quanh trục hoành có thể tích $$V=\pi \displaystyle\int_{0}^{\pi} (2+\sin x)\mathrm{\,d}x =\pi \left(2x-\cos x\right)\Bigg|_{0}^{\pi}=\pi(2\pi+2)=2\pi (\pi+1).$$}
\end{ex}

% nguồn: Phan1_Lop_12_Dai_so_Chuong_4
\begin{ex}%[2D4H3-3]%[ma:e01de334]
	Cho  hình phẳng $(D)$ được giới hạn bởi các đường $x=0$, $x=\pi$, $y=0$ và $y=-\sin x$. Thể tích $V$ của khối tròn xoay tạo thành khi quay $(D)$ xung quanh trục $Ox$ được tính theo công thức
	 \choice
		 {$V=\pi \displaystyle\int_0^{\pi}\left|\sin x\right|\mathrm{d}x$}
		{\True $V=\pi \displaystyle\int_0^{\pi}\sin  ^2x\mathrm{d}x$}
		{$V=\pi \left|\displaystyle\int_0^{\pi}(-\sin x)\mathrm{d}x  \right|$}
		 {$V=\displaystyle\int_0^{\pi}\sin ^2x\mathrm{d}x$}
	 \loigiai{
		 Ta có $V=\pi \displaystyle\int_0^{\pi}(-\sin x)^2\mathrm{d}x=\pi \displaystyle\int_0^{\pi}\sin ^2x\mathrm{d}x$.
	 }
\end{ex}

% nguồn: Phan1_Lop_12_Dai_so_Chuong_4
\begin{ex}%[2D4H3-3]%[ma:43a7d32d]
	Tính thể tích của  khối tròn xoay khi cho hình phẳng giới hạn bởi đồ thị hàm số $y=x^2$, trục hoành và đường thẳng $x=2$, khi quay xung quanh trục $Ox$ bằng
	\choice 
		 {\True $\dfrac{32\pi}{5}$}
		 {$\dfrac{\pi}{6}$}
		{$\dfrac{5\pi}{6}$} 
		 {$\dfrac{4\pi}{5}$}
	 \loigiai{
		Thể tích khối tròn xoay  là: $V=\pi \displaystyle\int_0^2(x^2)^2\mathrm{d}x=\left. \pi \left(\dfrac{1}{5} x^5\right)\right|_0^2=\dfrac{32\pi}{5}$.
	} 
\end{ex}

% nguồn: Phan1_Lop_12_Dai_so_Chuong_4
\begin{ex}%[2D4H3-3]%[ma:a33abf62]
	\immini{Gọi $(H)$ là hình phẳng giới  hạn bởi các đường $y = 2^x$, $y = 2$, $x = 3$ (phần gạch chéo trong hình vẽ), khi $(H)$ quay quanh trục hoành ta được khối tròn xoay có thể tích được tính theo công thức nào dưới đây?
	\choice 
		{$V = \pi \displaystyle\int_1^3 4^x \, \mathrm{d}x$} 
		 {$V = \pi \displaystyle\int_1^3 (2^x - 2)^2 \, \mathrm{d}x$}
		{\True  $V = \pi \displaystyle\int_1^3 (4^x - 4) \, \mathrm{d}x$}
		{$V  = \displaystyle\int_1^3 (2^x - 2) \, \mathrm{d}x$}}
		{\begin{tikzpicture}[scale=1]
		\draw[->] (-0.2,0) -- (4,0) node[below] {$x$}; 
		\draw[->] (0,-0.2) --  (0,9) node[left] {$y$};
		\begin{scope}
			\clip (1,2)  -- plot[domain=1:3,samples=100] (\x,{2^(\x)}) -- (3,2) -- cycle;
			\fill[pattern=north east lines] (1,2) -- plot[domain=1:3,samples=100] (\x,{2^(\x)}) -- (3,2)  -- cycle;
		\end{scope}
		\draw[domain=0:3.2,samples=100] plot  (\x,{2^(\x)}) node[right] {$y=2^x$};
		\draw[dashed] (0,2)  -- (3.2,2);
		\draw[dashed] (1,0)  -- (1,2);
		 \draw[dashed] (3,0) -- (3,{2^3});
		\draw[dashed]  (0,8) -- (3.2,8);
		\node at (1,-0.3) {$1$}; 
		\node  at (3,-0.3) {$3$};
		\node at (-0.3,2) {$2$}; 
		\node at  (-0.3,8) {$8$};
		 \node at (-0.2,-0.2) {O};
	\end{tikzpicture}
		} 
	 \loigiai{
		Khi quay hình phẳng $(H)$ quanh trục hoành, thể tích  khối tròn xoay được tính bởi công thức
		 \[
		V = \pi \displaystyle\int_a^b \left[f(x)^2 - g(x)^2  \right] \mathrm{d}x.
		\] 
		 Trong đó
		\begin{itemize}
			\item $f(x) = 2^x$ là hàm nằm phía trên, 
			\item $g(x) = 2$ là hàm  nằm phía dưới,
			\item giới hạn tích phân từ $x  = 1$ đến $x = 3$.
		\end{itemize}
		 Khi đó $V = \pi \displaystyle\int_1^3 \left[(2^x)^2 - 2^2 \right] \mathrm{d}x = \pi \displaystyle\int_1^3 (4^x - 4) \, \mathrm{d}x$.
	} 
\end{ex}

% nguồn: Phan1_Lop_12_Dai_so_Chuong_4
\begin{ex}%[2D4H3-3]%[ma:98a7f2a5]
	Cho hình phẳng $D$ giới hạn bởi đường  cong $y = \sqrt{x^2 +1}$, trục hoành và các đường thẳng $x=0, x=1$. Khối tròn xoay tạo thành khi quay $D$ quanh trục hoành có thể tích $V$ bằng
	\choice 
		{$V  = 2\pi$}
		 {\True $V = \dfrac{4\pi}{3}$}
		{$V = \dfrac{4}{3} $} 
		 {$V=2$}
	 \loigiai{
		Khối tròn xoay tạo thành khi quay $D$ quanh trục hoành có thể tích  là\\
	$V = \pi \displaystyle\int_{0}^{1} (\sqrt{x^2 +1})^2 dx =  \pi \displaystyle\int_{0}^{1} (x^2+1)dx = \pi \left(\dfrac{x^3}{3}+x\right)\Big|_0^1 = \pi \left(\dfrac{1}{3}+1\right) = \dfrac{4\pi}{3}$.}
\end{ex}

% nguồn: Phan1_Lop_12_Dai_so_Chuong_4
\begin{ex}%[2D4H3-3]%[ma:d71bd434]
	Cho hình phẳng $(H)$ giới hạn bởi parapol $(P)\colon y=x^2$ và đường thẳng $d\colon y=2x$. Thể tích khối tròn xoay  sinh bởi $(H)$ khi quay quanh trục $Ox$ bằng
	 \choice
		 {\True $\dfrac{64\pi}{15}$}
		{$\dfrac{16\pi}{15}$} 
		{$\dfrac{256\pi}{15}$} 
		{$\dfrac{4\pi}{3}$} 
	 \loigiai{
		Phương trình hoành độ giao  điểm $x^2=2x\Leftrightarrow x^2-2x=0\Leftrightarrow \hoac{&x= 0 \\ &x= 2.}$.
		\[V=\pi \displaystyle\int_0^2\left[(2x)^2-(x^2)^2\right] \; \mathrm{d}x=\pi \displaystyle\int_0^2\left(4x^2-x^4\right) \; \mathrm{d}x=\left.  \pi \left(\dfrac{4}{3} x^3-\dfrac{x^5}{5} \right)\right|_0^2=\dfrac{64\pi}{15}.\]
	} 
\end{ex}

% nguồn: Phan1_Lop_12_Dai_so_Chuong_4
\begin{ex}%[2D4V3-3]%[ma:e18acca5]
	Cho hình phẳng giới hạn bởi  các đường $y=x-\dfrac{1}{2} x^3$, $y=0$, $x=0$, $x=\sqrt{2}$. Thể tích vật thể tròn xoay được sinh ra bởi hình phẳng đó khi quay quanh trục $Ox$ có giá trị thuộc khoảng nào sau đây?
	\choice 
		{$\left(\dfrac{3}{2};2\right)$} 
		{$\left(1;\dfrac{3}{2}\right)$} 
		{$\left(0;\dfrac{1}{2}\right)$} 
		{\True $\left(\dfrac{1}{2};1\right)$} 
		\begin{center}
			\resizebox{!}{4cm}{\begin{tikzpicture}[line cap=butt,line join=miter,>=stealth]
		\tikzset{declare function={f (\x)=\x-0.5*(\x)^3; xa=sqrt (2); xb=sqrt (2/3); yb=f (xb);}} 
		\begin{scope}
			\clip (-0.5,-1.5)  rectangle (3.5,1.5);
			\fill[gray!35] plot[domain=0:xa] (\x,{f (\x)}); 
			\draw[samples=100] plot[domain=-1:4] (\x,{f (\x)}); 
			 \draw[samples=100,dashed] plot[domain=-1:4] (\x,{-f (\x)});
		\end{scope}
		 \draw[->] (-0.5,0)--(2.5,0)
			 node[shift={(-100:7pt)},font=\normalsize]{$x$};
		 \draw[->] (0,-1.5)--(0,1.6)
			 node[shift={(180:7pt)},font=\normalsize]{$y$};
		\draw (5pt,0) |- (0,5pt) 
			(0,0) node[shift={(245:13pt)},font=\normalsize]{$O$}; 
		\draw (xb,yb) arc (90:270:{0.25*yb} and yb); 
		 \draw[dashed] (xb,yb) arc (90:-90:{0.25*yb} and yb);
		\path  (0.9,yb) node[above,font=\scriptsize]{$y=x-\tfrac{1}{2}x^3$};
	\end{tikzpicture}}
	\end{center}
	 \loigiai{
		Thể tích của vật thể là: $V=\pi \displaystyle\int_0^{\sqrt{2}}\left(x-\dfrac{1}{2} x^3\right)^2\mathrm{d}x=\dfrac{16\sqrt{2} \pi}{105} \approx 0{,}677\in \left(\dfrac{1}{2};  1\right)$.
	} 
\end{ex}

% nguồn: Phan1_Lop_12_Dai_so_Chuong_4
\begin{ex}%[2D4V3-3]%[ma:ab824c6b]
	Cho hình phẳng $(H)$ giới hạn bởi $y=2x-x^2$, $y=0$. Tính thể tích của khối tròn xoay thu  được khi quay $(H)$ xung quanh trục $Ox$ ta được $V=\pi \left(\dfrac{a}{b}+1\right)$ với $a,b$ là các số tự nhiên, $\dfrac{a}{b}$ là phân số tối giản. Khi đó
	\choice 
		 {$ab=16$}
		{$ab=12$} 
		 {\True $ab=15$}
		{$ab=18$} 
	\loigiai{ 
		Xét phương trình $2x-x^2=0\Leftrightarrow \hoac{&x=  0\\ &x= 2.}$\\
		Tính thể tích của khối tròn xoay thu được khi quay $(H)$ xung quanh trục  $Ox$ là
		\[V=\pi \displaystyle\int_0^2\left(2x-x^2\right)^2  \mathrm{d}x=\dfrac{16}{15} \pi=\left(1+\dfrac{1}{15}\right)\pi \Rightarrow a=1,b=15\Rightarrow ab=15.\]
	} 
\end{ex}


%%%%%%%%%%%%%%%% ĐÚNG/SAI %%%%%%%%%%%%%%%%

% nguồn: Phan1_Lop_12_Dai_so_Chuong_4
\begin{ex}%[2D4H3-3]%[ma:00fcc1c3]
	Cho đồ thị các hàm số $f(x)=x+1$  và $g(x)=(0{,}7)^x$ như hình vẽ
	\begin{center}
		\begin{tikzpicture}[scale=1,font=\footnotesize,line join=round,line cap=round,>=stealth]
				\draw[->] (-3.2,0)--(4.1,0) node[below left] {$x$}; 
				\draw[->] (0,-1.5)--(0,4.1)  node[below left] {$y$};
				\draw (0,0) node [below left] {$O$}  circle (1pt)
					 (-1,0) node [below] {$-1$} circle (1pt)
					 (1,0) node [below] {$1$} circle (1pt)
					(2,0)  node [below] {$2$} circle (1pt)
					(0,1)  node [left=0.1] {$1$} circle (1pt)
					(0,2) node [left]  {$2$} circle (1pt);
				\begin{scope}
					\clip (-3,-3) rectangle (4,4); 
					\draw[samples=200,domain=-2.5:3,smooth,variable=\x]  plot (\x,{(\x)+1})
						(3.8,3) node [below left=0.2] {$f  (x)=x+1$};
					 \draw[samples=200,domain=-4:4,smooth,variable=\x] plot (\x,{(0.7)^(\x)})
						(-0.5,3) node [below left=0.2] {$g  (x)=0{,}7^x$};
				\end{scope}
			\end{tikzpicture}
	\end{center}
	 \choiceTF
		 {\True $\displaystyle\int_{-1}^0 f(x)\mathrm{\,d}x=\dfrac{1}{2}$}
		{Thể tích vật thể tròn xoay được sinh ra khi quay hình phẳng giới hạn bởi các đường $y=f(x)$, trục hoành và hai đường  thẳng $x=-1$; $x=0$ quanh trục hoành bằng $9\pi$}
		{Thể tích vật thể tròn xoay được sinh ra khi quay hình phẳng giới hạn bởi các đường $y=g(x)$, trục hoành và hai đường thẳng $x=0$; $x=2$ quanh trục hoành có giá  trị xấp xỉ bằng $7{,}9$ (làm tròn kết quả đến hàng phần mười)}
		{\True Khối tròn xoay sinh ra khi quay hình phẳng giới hạn bởi các đường $y=f(x)$, $y=g(x)$, trục  hoành và hai đường thẳng $x=-1$, $x=2$ quanh trục hoành có thể tích xấp xỉ bằng $4{,}4$ (làm tròn kết quả đến hàng phần mười)}
	 \loigiai{
		\begin{itemchoice}
			\itemch 
			Ta có $\displaystyle\int_{-1}^0  f(x)\mathrm{\,d}x=\displaystyle\int_{-1}^0 (x+1)\mathrm{\,d}x=\left(\dfrac{x^2}{2}+x\right)\Bigg|_{-1}^0=0-\left(\dfrac{1}{2}-1\right)=\dfrac{1}{2}$.
			\itemch 
			Thể tích cần tìm  là
			 \[V_1=\pi \displaystyle\int_{-1}^0 f^2(x)\mathrm{\,d}x=\pi \displaystyle\int_{-1}^0 (x+1)^2\mathrm{\,d}x=\pi \displaystyle\int_{-1}^0(x^2+2x+1)\mathrm{\,d}x=\pi\left(\dfrac{x^3}{3}+x^2+x\right)\Bigg|_{-1}^0=\dfrac{\pi}{3}.\]
			 \itemch
			Thể tích cần tìm là \[V_2=\pi \displaystyle\int_0^2 g^2(x)\mathrm{\,d}x=\pi \displaystyle\int_0^2  (0{,}7)^{2x}\mathrm{\,d}x=\pi \displaystyle\int_0^2(0{,}49)^x\mathrm{\,d}x=\pi\cdot\dfrac{0{,}49^x}{\ln 0{,}49}\Bigg|_0^2=\dfrac{-0{,}7599\pi }{\ln0{,}49}\approx 3{,}3.\]
			 \itemch
			Từ đồ thị hàm số suy ra thể tích khối tròn xoay sinh ra khi quay hình phẳng giới hạn bởi các đường $y=f(x)$, $y=g(x)$, trục hoành và hai đường thẳng $x=-1$, $x=2$ quanh trục hoành là  $V=V_1+V_2=\dfrac{\pi}{3}+\dfrac{-0{,}7599\pi }{\ln0{,}49}\approx 4{,}4$.
		\end{itemchoice}
	 }
\end{ex}

% nguồn: Phan1_Lop_12_Dai_so_Chuong_4
\begin{ex}%[2D4H3-3]%[ma:45fc792f]
	 \immini{
	Đồ thị các đường $y=x$, $y=-\sqrt{x}$, $x=3$ cho bởi hình vẽ. Gọi $S_1$ là  diện tích hình phẳng giới hạn bởi các đường $y=x$, trục hoành và hai đường thẳng $x=0$, $x=3$. Gọi $S_2$ là diện tích hình phẳng giới hạn bởi các đường $y=-\sqrt{x}$, trục hoành và hai đường thẳng $x=0$, $x=3$.
	 \choiceTF
		 {\True $S_1=\displaystyle\int_0^3 x\,\mathrm{\,d}x$}
		{$S_2=\displaystyle\int_0^3 -\sqrt{x}\,\mathrm{\,d}x$} 
		{\True Thể tích khối tròn xoay sinh ra khi quay quanh trục $Ox$  hình phẳng giới hạn bởi $y=x$, $y=0$, $x=0$, $x=3$ bằng $9\pi$}
		{Diện tích hình phẳng giới hạn bởi $y=-\sqrt{x}$, $y=x$, $x=3$, $x=0$ bằng $12{,}5$ (kết quả chính xác đến hàng phần chục)} 
		}{\begin{tikzpicture}[scale=0.7,font=\footnotesize,line join=round,
			line cap=round,>=stealth] 
		 \def \xmin{-2}
		\def  \xmax{6}
		\def \ymin{-3} 
		 \def \ymax{4}
		\draw[->]  (\xmin,0)--(\xmax,0) node[shift=(-110:0.2)] {$x$};
		\draw[->] (0,\ymin)--(0,\ymax) node[shift=(-150:0.2)] {$y$}; 
		\fill (0,0) circle (1pt) node[shift=(135:0.25)]{$O$}; 
		\draw (5,4) node[below]  {$y=x$}
			(5.5,-1.5) node{$y=-\sqrt{x}$}; 
		\clip (\xmin,\ymin) rectangle  (\xmax,\ymax);
		 \draw[smooth,samples=100,domain=\xmin:\xmax] plot (\x,{(\x)});
		\draw[smooth,samples=100,domain=0:\xmax] plot (\x,{-sqrt  (\x)});
		 \foreach \x/\g in {3/-45}
		 \draw (\x,1pt)--(\x,-1pt) + (\g:0.4) node {$\x$};
		\draw  (3,\ymin)--(3,\ymax);
	\end{tikzpicture}}
	 \loigiai{
		\begin{itemchoice}
			\itemch Ta có 
			\[S_1=\displaystyle\int_0^3 |x|\,\mathrm{\,d}x=\int_0^3 x\,\mathrm{\,d}x.\] 
			\itemch Ta có 
			\[S_2=\displaystyle\int_0^3  |-\sqrt{x}|\,\mathrm{\,d}x=\int_0^3 \sqrt{x}\,\mathrm{\,d}x.\]
			\itemch Thể tích khối tròn xoay sinh ra khi quay quanh trục $Ox$ hình phẳng giới hạn bởi $y=x$,  $y=0$, $x=0$, $x=3$ bằng
			\[V_1 = \displaystyle \pi  \int_0^3 x^2 \, \mathrm{\,d}x=9 \pi.\]
			\itemch Diện tích hình phẳng giới hạn bởi $y=-\sqrt{x}$, $y=x$, $x=3$, $x=0$  bằng
			\[S = \displaystyle \int_0^3  (x+\sqrt{x})\,\mathrm{\,d}x \approx 8{,}0.\]
	\end{itemchoice}}
\end{ex}

% nguồn: Phan1_Lop_12_Dai_so_Chuong_4
\begin{ex}%[2D4H3-3]%[ma:dbb7a78f]
	 Cho hai hàm số $f(x) = x\sqrt{x} + 8$, $g(x) = 5^x - e^x$ và $F(x)$ là một nguyên hàm của $f(x)$.
	 \choiceTF
		{$\displaystyle \int  g\left( x \right)\,\mathrm{d}x = 5^x \ln 5 - e^x + C$}
		{\True $\displaystyle \int_{1}^{3} f\left( x \right)\,\mathrm{d}x = F\left( 3 \right) - F\left( 1  \right)$}
		{\True Diện tích hình phẳng giới  hạn bởi đồ thị hàm số $y = g\left( x \right)$, trục $Ox$ và các đường thẳng $x = 1$, $x = 2$ bằng $7{,}8$ (làm tròn đến hàng phần chục)}
		{Thể tích khối tròn xoay khi quay hình phẳng giới hạn bởi đồ thị hàm số $y = f(x)$, trục $Ox$, $x = 0$, $x = 3$ quanh trục $Ox$ bằng $312$ (làm tròn đến hàng đơn vị)} 
	 \loigiai{
		\begin{itemchoice}
			\itemch $\displaystyle \int g\left( x \right)\,\mathrm{d}x  = \dfrac{5^x}{\ln 5} - e^x + C$.
			\itemch $\displaystyle \int_{1}^{3} f\left(  x \right)\,\mathrm{d}x = F\left( 3 \right) - F\left( 1 \right)$.
			\itemch Diện tích hình phẳng cần tìm là $S=\displaystyle \int_{1}^{2} \left| 5^x -e^x \right|\,\mathrm{d}x = 7{,}8$. 
			\itemch Thể tích khối tròn xoay cần tìm là $V=\pi \displaystyle  \int_{0}^{3}\left(x\sqrt{x} + 8\right)^2\,\mathrm{d}x=980.$
		\end{itemchoice}
	} 
\end{ex}

% nguồn: Phan1_Lop_12_Dai_so_Chuong_4
\begin{ex}%[2D4H3-3]%[ma:c412709f]
	Cho hàm số $y=f(x)=2x+3$. Biết  $F(x)$ là một nguyên hàm của $f(x)$ trên $\mathbb{R}$. Khi đó
	\choiceTF 
		 {Biết $F(1)=2$ thì $F(x)=x^2+3x+2$}
		{\True Giá trị của $\displaystyle\int_0^2f(x) \mathrm{\,d}x-\displaystyle\int_5^2f(x) \mathrm{\,d}x+\displaystyle\int_{-1}^0f(x) \mathrm{\,d}x$  bằng $42$}
		{Diện tích hình phẳng giới hạn bởi các đường  $y=f(x)$, trục hoành và $x=-2$, $x=1$ bằng $6$}
		{\True Thể tích khối tròn xoay thu được khi quay hình phẳng giới hạn bởi hai đường $y=f(x)$ và $y=x^2-2x+6$ quanh  trục $Ox$ bằng $\dfrac{1\,556\pi}{15}$}
	\loigiai{ 
		\begin{itemchoice}
			\itemch Ta có $F(x)=\displaystyle\int  f(x) \mathrm{\,d}x=\displaystyle\int(2x+3) \mathrm{\,d}x=x^2+3x+C$.
			$$ 
			F(1)=2 \Leftrightarrow 1^2+3\cdot1+C=2  \Rightarrow C=-2.
			 $$
			 Vậy, $F(x)=x^2+3x-2$.
			\itemch Ta có 
			 $$
			\begin{aligned}
				\displaystyle\int_0^2 f(x) \mathrm{\,d}x-\displaystyle\int_5^2 f(x) \mathrm{\,d}x+\displaystyle\int_{-1}^0 f(x) \mathrm{\,d}x&=\displaystyle\int_{-1}^0 f(x) \mathrm{\,d}x+\displaystyle\int_0^2 f(x)  \mathrm{\,d}x-\displaystyle\int_5^2 f(x) \mathrm{\,d}x \\
				&=\displaystyle\int_{-1}^0 f(x) \mathrm{\,d}x+\displaystyle\int_0^2 f(x) \mathrm{\,d}x+\displaystyle\int_2^5 f(x) \mathrm{\,d}x 
				 \\
				&=\displaystyle\int_{-1}^5 f(x) \mathrm{\,d}x 
				\\ 
				&=\displaystyle\int_{-1}^5  (2x+3) \mathrm{\,d}x
				 \\
				 &=(x^2+3x)\bigg|_{-1}^5
				\\ 
				&=40-(-2)=42. 
			\end{aligned}
			 $$
			\itemch Diện tích hình phẳng cần tìm  là $S=\displaystyle\int_{-2}^1|2x+3|\mathrm{\,d}x=\dfrac{13}{2}$.
			\itemch Phương trình hoành độ giao  điểm $2x+3=x^2-2x+6\Leftrightarrow x^2-4x+3=0\Leftrightarrow\hoac{&x=1\\&x=3.}$\\
			Thể tích khối tròn xoay cần tìm là: $V=\pi \displaystyle\int_1^3\left|(2x+3)^2-\left(x^2-2x+6\right)^2\right|\mathrm{\,d}x=\dfrac{1\,556\pi}{15}$. 
		\end{itemchoice}
	 }
\end{ex}

% nguồn: Phan1_Lop_12_Dai_so_Chuong_4
\begin{ex}%[2D4H3-3]%[ma:e25763e7]
	Cho hàm số $f(x) = x^2 - 3x + 5$. Gọi $F(x)$  là nguyên hàm của $f(x)$ và $F(1) = 5$.
	\choiceTF 
		{$F(4) =  37$}
		{\True $\displaystyle \int_2^5 \left(x^2 - 3x  + 5\right) \mathrm{\,d}x = \dfrac{45}{2}$}
		{Diện tích  hình phẳng giới hạn bởi đồ thị hai hàm số $f(x) = x^2 - 3x + 5$, $g(x) = x+2$ và hai đường thẳng $x=1$, $x=3$ bằng $4$}
		{\True Thể tích khối tròn xoay tạo thành khi quay hình phẳng giới hạn bởi đồ thị hai hàm số $f(x) = x^2 - 3x + 5$, $g(x) = x+2$ và hai  đường thẳng $x=1$, $x=3$ quanh trục $Ox$ bằng $\dfrac{48}{5} \pi$}
	 \loigiai{
		\begin{itemchoice}
			\itemch Ta có $F(x) = \displaystyle \int \left(x^2 - 3x + 5\right) \mathrm{\,d}x = \dfrac{1}{3}x^3 - \dfrac{3}{2}x^2 + 5x +  C$. \\
			Ta có  $F(1) = \dfrac{1}{3} \cdot 1^3-\dfrac{3}{2}\cdot 1^2+5\cdot1 +C=5$ nên $C = \dfrac{7}{6}$. \\
			$\Rightarrow F(x)=\dfrac{1}{3}x^3-\dfrac{3}{2}x^2+5x+\dfrac{7}{6}$.\\ 
			Do đó, $F(4)  = \dfrac{37}{2} \neq 37$.
			\itemch Ta có $\displaystyle \int_2^5 \left(x^2 - 3x + 5\right) \mathrm{\,d}x = \left( \dfrac{x^3}{3} - \dfrac{3}{2}x^2  + 5x \right) \Bigg|_2^5 =\dfrac{45}{2}$.
			\itemch Phương trình hoành độ giao  điểm của hai đồ thị
			\[x^2 - 3x + 5 = x + 2 \Leftrightarrow x^2 - 4x + 3 = 0  \Leftrightarrow x = 1, x = 3.\]
			Diện tích hình phẳng  là
			 \[S = \displaystyle \int_1^3 \left|x^2 - 3x + 5 - (x+2)\right| \mathrm{\,d}x = \dfrac{4}{3}.\]
			\itemch  Thể tích khối tròn xoay quanh trục $Ox \colon
			V = \pi \displaystyle \int_{1}^{3} \left[\left(x^2 - 3x + 5\right)^2 - (x+2)^2 \right] \mathrm{\,d}x = \dfrac{48}{5}  \pi$.
		\end{itemchoice}
	} 
\end{ex}

% nguồn: Phan1_Lop_12_Dai_so_Chuong_4
\begin{ex}%[2D4H3-3]%[ma:a1464ff5]
	Khối chỏm cầu  có bán kính $R=5$ và chiều cao $h=1$ sinh ra khi quay hình phẳng giới hạn bởi cung tròn có phương trình $y=\sqrt{25-x^2}$ , trục hoành và hai đường thẳng $x=4$, $x=5$ xung quanh trục $Ox$.
	 \choiceTF
		 {Khoảng cách từ tâm của khối cầu đến khối chỏm cầu bằng $3$}
		{\True Thể tích  của khối chỏm cầu $V_1$ được tính theo công thức $V_1=\pi\displaystyle\int_4^5\left(25-x^2\right){\rm{d}}x$}
		{\True Thể  tích của khối chỏm cầu $V_1=\dfrac{14\pi}{3}$}
		{\True Gọi $V_2$ là thể tích của nửa khối cầu có bán kính bằng $5$. Tỉ số thể tích $\dfrac{V_1}{V_2}=\dfrac{7}{125}$} 
	\loigiai{ 
		\begin{itemchoice}
			\itemch Khoảng cách từ tâm của  khối cầu đến khối chỏm cầu bằng $5-1=4$.
			\itemch Thể tích của khối chỏm cầu được tính theo  công thức
			$$V_1=\pi\displaystyle\int_4^5\left(\sqrt{25-x^2}\right)^2\rm{d}x=\pi\displaystyle\int_4^5\left(25-x^2\right){\rm{d}}x.$$ 
			\itemch Ta có thể tích khối chỏm cầu là  $V=\pi\displaystyle\int_4^5\left(25-x^2\right){\rm{d}}x=\left.\pi\left(25x-\dfrac{x^3}{3}\right)\right|_4^5=\dfrac{14\pi}{3}$.
			\itemch Gọi $V_2$  là thể tích của nửa khối cầu có bán kính bằng $5$.\\
			Ta  có $V_2=\dfrac{1}{2}\cdot\dfrac{4}{3}\pi\cdot{5^3}=\dfrac{250}{3}\pi $ nên $\dfrac{V_1}{V_2}=\dfrac{\dfrac{14}{3}\pi}{\dfrac{250}{3}\pi}=\dfrac{7}{125}$.
		\end{itemchoice}
	 }
\end{ex}

% nguồn: Phan1_Lop_12_Dai_so_Chuong_4
\begin{ex}%[2D4H3-3]%[ma:a6096cb9]
	Cho hàm số $f(x) = x^3 - 3x^2 + 2x$ có đồ thị  như hình dưới đây\\
	\centerline{\begin{tikzpicture}[scale=0.7,line join=round,line cap=round,>=stealth,font=\scriptsize]
		\tikzset{every  node/.style={scale=1}}
		 \def\xmin{-1}\def\xmax{3}\def\ymin{-2}\def\ymax{2}
		\draw[->] (\xmin-0.2,0)--(\xmax+0.2,0) node[below]{$x$}; 
		\draw[->]  (0,\ymin-0.2)--(0,\ymax+0.2) node[right]{$y$};
		\draw  (0,0) node[below left]{$O$};
		 \foreach \x in {1,2}\draw (\x,0.1)--(\x,-0.1) node[below]{$\x$};
		\foreach \y in {}\draw  (0.1,\y)--(-0.1,\y) node[left]{$\y$};
		\clip (\xmin,\ymin) rectangle (\xmax,\ymax); 
		\draw[smooth,samples=200,domain=\xmin:\xmax] plot  (\x,{1*((\x)^3)+-3*((\x)^2)+2*(\x)+0});
	\end{tikzpicture} }
	\choiceTF 
		{\True $\displaystyle\int_0^1  f'(x) \mathrm{\,d}x = 0$}
		{Gọi $F(x)$ là một nguyên hàm của hàm số  $f(x)$, khi đó $F(x) = \displaystyle\int f'(x) \mathrm{\,d}x$}
		{Diện tích hình phẳng giới hạn bởi đồ thị hàm số $y=f(x)$, trục hoành và hai đường thẳng $x=0, x=2$ là $S=\displaystyle\int_0^2 f(x)  \mathrm{\,d}x$}
		{\True Gọi $(D)$ là hình phẳng  giới hạn bởi đồ thị hàm số $y=f(x)$, trục hoành, trục tung và đường thẳng $x=1$. Thể tích khối tròn xoay được tạo thành khi quay $(D)$ quanh trục hoành được tính theo công thức $V=\pi \displaystyle\int_0^1 f^2(x) \mathrm{\,d}x$}
	\loigiai{ 
		\begin{itemchoice}
			\itemch Ta  có $\displaystyle\int_0^1 f'(x) \mathrm{\,d}x=f(x)\bigg|_0^1=f(1)-f(0)=0$. \itemch
			Ta có $F(x)$ là một nguyên hàm của hàm số $f(x)$ nên $F(x)= \displaystyle\int f(x) \mathrm{\,d}x$ 
			\itemch Diện tích hình phẳng giới hạn bởi $y=f(x)$, trục hoành và $x=0,  x=2$ là $$S=\displaystyle\int_0^2 \left|f(x)\right| \mathrm{\,d}x
			=\displaystyle\int_0^1 \left|f(x)\right| \mathrm{\,d}x+\displaystyle\int_1^2 \left|f(x)\right|  \mathrm{\,d}x =\displaystyle\int_0^1 f(x) \mathrm{\,d}x-\displaystyle\int_1^2 f(x) \mathrm{\,d}x.$$
			\itemch Thể tích khối tròn  xoay được tạo thành khi quay $(D)$ quanh trục hoành được tính theo công thức $V=\pi \displaystyle\int_0^1 f^2(x) \mathrm{\,d}x$
		\end{itemchoice}
	} 
\end{ex}

% nguồn: Phan1_Lop_12_Dai_so_Chuong_4
\begin{ex}%[2D4H3-3]%[ma:4520df7d]
	Cho tam giác vuông $OAB$  có cạnh $OA=a$ nằm trên tục $Ox$ và $\widehat{AOB}=\alpha\left(0 <\alpha <\dfrac{\pi}{2}\right)$. Gọi $\beta $ là khối tròn xoay sinh ra khi quay miền tam giác $OAB$ xung quanh trục $Ox$.
	\begin{center}
		\begin{tikzpicture}[scale=1,font=\footnotesize,line join=round,line cap=round,>=stealth]
				\fill[yellow!30] (0,0)--(4,2)--(4,0)--cycle; 
				\draw[->] (4,0) -- (5.5,0) node[right]  {$x$};
				\draw[->] (0,-2.5)  -- (0,2.5) node[above] {$y$};
				\draw[->] (0,0) -- (-1.5,-1.5) node[below left] {$z$}; 
				\draw (0,0)--(4,2) 
					 (0,0)--(4,-2)
					 (4,2)--(4,0)
					(-0.5,0)--(0,0); 
				\draw[dashed]  (0,0)--(4,0);
				\draw  (4,0) ellipse (0.5 and 2);
				\fill (0,0) circle (1pt) node[above left]{$O$}; 
				\fill (4,0) circle (1pt)  node[below]{$A$};
				\fill (4,2) circle  (1pt) node[above]{$B$};
				\clip (4,0)  -- (0,0) -- (4,2);
				\draw (0,0)  circle (1cm);
				 \draw ($(0,0)+(1,0)$) node[above right]{$\alpha$};
				\node[below] at (2.5,-2.5) {Hình  $4.31$};
			\end{tikzpicture}
	\end{center}
	 \choiceTF
		{\True Khi $\alpha=\dfrac{\pi}{4}$ thì thể tích $V$  của khối $\beta $ là $\dfrac{\pi{a^3}}{3}$ (đvtt)}
		{\True Khi $\alpha=\dfrac{\pi}{6}$ thì thể tích $V$ của khối $\beta $ là $\dfrac{\pi{a^3}}{9}$ (đvtt)} 
		{\True Khi thể tích  $V$ của khối $\beta $ là $\dfrac{4\pi{a^3}}{3}$ thì giá trị $\cos\alpha <\dfrac{1}{2}$}
		{\True Khi  $\tan\alpha=\cot\alpha $ thì thể tích $V$ của khối $\beta $ là $\dfrac{\pi{a^3}}{3}$}
	\loigiai{ 
		\begin{itemchoice}
			\itemch Do $OB$ đi qua gốc tọa độ và tạo với  $Ox$ một góc $\dfrac{\pi}{4}$ nên $OB\colon y=\left(\tan\dfrac{\pi}{4}\right)x=x$.\\
			Khi đó, thể tích của khối $\beta $ theo  $V=\pi\displaystyle\int_0^a{x^2\rm{d}x}=\left.\dfrac{\pi{x^3}}{3}\right|_0^a=\dfrac{\pi{a^3}}{3}$ (đvtt).
			\itemch Do $OB$ đi qua gốc  tọa độ và tạo với $Ox$ một góc $\dfrac{\pi}{6}$ nên $OB\colon y=\left(\tan\dfrac{\pi}{6}\right)x=\dfrac{x}{\sqrt 3}$.\\
			Khi đó, thể tích của khối  $\beta $ theo $V=\pi\displaystyle\int_0^a{\left(\dfrac{x}{\sqrt 3}\right)^2\rm{d}x}=\pi\displaystyle\int_0^a{\dfrac{x^2}{3}{\rm{d}}x}=\left.\dfrac{\pi{x^3}}{9}\right|_0^a=\dfrac{\pi{a^3}}{9}$ (đvtt).
			\itemch  Do $OB$ đi qua gốc tọa độ và tạo với $Ox$ một góc $\alpha $ nên $OB\colon y=\tan\alpha\cdot x$.\\
			Khi đó, thể tích của khối $\beta $ theo $V=\pi\displaystyle\int_0^a{\left(\tan\alpha  \cdot x\right)^2\rm{d}x}=\left.\dfrac{\pi{\tan^2}\alpha \cdot x^3}{3}\right|_0^a=\dfrac{\pi{\tan^2}\alpha\cdot a^3}{3}$ (đvtt).\\
			Do $V=\dfrac{4\pi{a^3}}{3}\Leftrightarrow\dfrac{\pi{\tan^2}\alpha  \cdot a^3}{3}=\dfrac{4\pi{a^3}}{3}\Leftrightarrow{\tan ^2}\alpha=4\Rightarrow\dfrac{1}{\cos^2\alpha}=\tan ^2\alpha+1=5\Rightarrow\cos\alpha=\dfrac{\pm 1}{\sqrt 5}$.\\
			Mặt khác $0  <\alpha <\dfrac{\pi}{2}$ nên $\cos\alpha=\dfrac{1}{\sqrt 5}$.
			\itemch Do $OB$ đi qua gốc tọa độ và  tạo với $Ox$ một góc $\alpha $ nên $OB\colon y=\tan\alpha\cdot x$.\\
			Khi đó thể tích của khối $\beta $  theo $V=\pi\displaystyle\int_0^a{\left(\tan\alpha \cdot x\right)^2\rm{d}x}=\left.\dfrac{\pi{\tan^2}\alpha \cdot x^3}{3}\right|_0^a=\dfrac{\pi{\tan^2}\alpha\cdot a^3}{3}$ (đvtt).\\
			Do $\tan\alpha=\cot\alpha\Rightarrow{\tan ^2}\alpha=\cot\alpha \cdot\tan\alpha=1\Leftrightarrow\tan\alpha=\pm 1$.\\ 
			Mặt khác $0 <\alpha <\dfrac{\pi}{2}$ nên $\tan\alpha=1$ $\Rightarrow V=\dfrac{\pi{\tan^2}\alpha \cdot a^3}{3}=\dfrac{\pi  \cdot a^3}{3}$ (đvtt).
		\end{itemchoice}
	 }
\end{ex}

% nguồn: Phan1_Lop_12_Dai_so_Chuong_4
\begin{ex}%[2D4H3-3]%[ma:f2c62e59]
	Cho hình phẳng $(H)$ giới hạn bởi  đường cong $y=f(x)=\dfrac{1}{3}x^3-x^2$ và trục hoành.
	 \choiceTF
		{\True  $\displaystyle\int_1^3f(x)\,\mathrm{d}x=-2$}
		{\True  $\displaystyle\int f(x)\,\mathrm{d}x=\dfrac{x^4}{12}-\dfrac{x^3}{3}+C$}
		{Diện tích hình phẳng $(H)$ bằng $\dfrac{9}{2}$} 
		 {Thể tích khối tròn xoay khi quay hình phẳng $(H)$ quanh trục $Ox$ (làm tròn đến hàng phần chục) bằng $2{,}3$ (đơn vị thể tích)}
	\loigiai{Phương trình hoành độ giao điểm của hàm số $y=f(x)=\dfrac{1}{3}x^3-x^2$ và  trục hoành là
		\begin{eqnarray*}
			 &&\dfrac{1}{3}x^3-x^2=0\\
			 &\Rightarrow&\hoac{&x=0\\&x=3.}
		\end{eqnarray*}
		\begin{itemchoice}
			\itemch \textbf{Đúng}. $\displaystyle\int_1^3f(x)\,\mathrm{d}x=\int_1^3\left(\dfrac{1}{3}x^3-x^2\right)\,\mathrm{d}x=\dfrac{x^4}{12}-\dfrac{x^3}{3}\Bigg|_1^3=\dfrac{3^4}{12}-\dfrac{3^3}{3}-\left(\dfrac{1^4}{12}-\dfrac{1^3}{3}\right)=-2$. 
			 \itemch \textbf{Đúng}. $\displaystyle\int f(x)\,\mathrm{d}x=\int\left(\dfrac{1}{3}x^3-x^2\right)\,\mathrm{d}x=\dfrac{x^4}{12}-\dfrac{x^3}{3}+C$.
			\itemch \textbf{Sai}. 
			Diện tích hình  phẳng $(H)$ là
			\begin{eqnarray*}
				S=\displaystyle\int_0^3 \left|f(x) \right|\,\mathrm{d}x 
				=\int_0^3 \left|\dfrac{1}{3}x^3-x^2  \right|\,\mathrm{d}x=\dfrac{9}{4}\neq\dfrac{9}{2}.
			\end{eqnarray*}
			\itemch  \textbf{Sai}.
			Thể tích khối tròn xoay khi quay hình phẳng  $(H)$ quanh trục $Ox$ là
			\begin{eqnarray*}
				V=\pi\displaystyle\int_{0}^{3}\left(\dfrac{1}{3}x^3-x^2\right)^2\,\mathrm{d}x=\dfrac{81\pi}{35}\approx 7{,}3\, (\text{đơn vị  thể tích}).
			\end{eqnarray*}
		\end{itemchoice}
	 }
\end{ex}

% nguồn: Phan1_Lop_12_Dai_so_Chuong_4
\begin{ex}%[2D4V3-3]%[ma:576b3b4a]
	Cho hàm số $f(x)=x^2-2x$ có đồ thị là $(C)$ và đường  thẳng $d \colon y=x$.
	 \choiceTF
		{\True Tích phân $\displaystyle\int_0^1f(x)\mathrm{\,d}x$  bằng $-\dfrac{2}{3}$}
		{Hình phẳng giới hạn bởi $(C)$, trục hoành và  hai đường thẳng $x=1$, $x=2$ có diện tích bằng $\dfrac{4}{3}$}
		{\True Hình phẳng giới hạn bởi $(C)$ và $d$ có diện tích bằng  $\dfrac{9}{2}$}
		{\True Gọi $(H)$ là hình phẳng giới hạn bởi $(C)$, trục hoành và hai đường thẳng $x=0$, $x=1$. Khối tròn  xoay thu được khi cho $(H)$ quay quanh trục hoành có thể tích bằng $\dfrac{8\pi}{15}$}
	\loigiai{ 
		\begin{itemchoice}
			\itemch  \textbf{Đúng}.\\
			 $\displaystyle\int_0^1f(x)\mathrm{\,d}x=\displaystyle\int_0^1 (x^2-2x)\mathrm{\,d}x=-\dfrac{2}{3}$.
			\itemch  \textbf{Sai}.\\
			Diện tích hình phẳng giới  hạn bởi $(C)$, trục hoành và hai đường thẳng $x=1$, $x=2$ là $\displaystyle\int_1^2\left|x^2-2x\right|\mathrm{\,d}x=\dfrac{2}{3}$.
			\itemch \textbf{Đúng}.\\ 
			Xét $x^2-2x=x \Leftrightarrow x^2 -3x=0  \Leftrightarrow \hoac{&x=0 \\ & x=3.}$\\
			Diện tích hình phẳng  giới hạn bởi $(C)$ và $d$ là
			 $\displaystyle\int_0^3\left|(x^2-2x)-x\right|\mathrm{\,d}x=\dfrac{9}{2}$.
			\itemch \textbf{Đúng}.\\ 
			Thể tích khối  tròn xoay thu được khi cho $(H)$ quay quanh trục hoành là $$V=\pi \displaystyle\int_0^1 (x^2-2x)^2 \mathrm{\,d}x=\dfrac{8\pi}{15}.$$
		\end{itemchoice}
	} 
\end{ex}

% nguồn: Phan1_Lop_12_Dai_so_Chuong_4
\begin{ex}%[2D4V3-3]%[ma:6513ee71]
	Cho hàm số $y=-x^3+x^2+2 x$ có đồ thị $(C)$. Hình  phẳng $\left(H_1\right)$ được giới hạn bởi $(C)$, trục hoành, hai đường thẳng $x=-1, x=0$ và hình phẳng $\left(H_2\right)$ được giới hạn bởi $(C)$, trục hoành, hai đường thẳng $x=0, x=2$ (như hình vẽ).
	\begin{center}
		\begin{tikzpicture}[line join=round,line cap=round,>=stealth]
				 \tikzset{every node/.style={scale=0.9}}
				 \draw[->] (-1.6,0)--(2.6,0) node[below] {$x$};
				\draw[->]  (0,-1.1)--(0,3.1) node[left] {$y$};
				\draw (0,0) node  [below right] {$O$};
				 \foreach \x/\nx in {-1/-1,2/2}
				\draw[thin]  (\x,1pt)--(\x,-1pt) node [below left] {$\nx$};
				\begin{scope}
					\clip (-1.5,-1)  rectangle (2.5,3);
					 \draw[samples=200,domain=-1.5:2.5,smooth,variable=\x] plot (\x,{-1*((\x)^3)+1*((\x)^2)+2*(\x)+0});
					\draw[pattern  = north east lines,line width = 1.2pt] (-1,0) plot[domain=-1:2] (\x,{-1*(\x)^3+(\x)^2+2*(\x)})--(2,0)--(-1,0)--cycle;
					 \node[red] at (-0.6,-0.3) {$H_1$};
					 \node[red] at (1,1) {$H_2$};
				\end{scope}
			\end{tikzpicture}
	\end{center}
	\choiceTF 
		{Diện tích của hình $\left(H_1\right)$ được tính bởi  $\displaystyle\int_{-1}^0\left(-x^3+x^2+2 x\right)\mathrm{\,d}x$}
		{\True Thể tích của  khối tròn xoay được tạo thành khi quay hình $\left(H_2\right)$ quanh trục $O x$ bằng $\dfrac{464}{105} \pi$}
		{Tổng diện tích của hình $\left(H_1\right)$ và  $\left(H_2\right)$ bằng $\dfrac{9}{4}$}
		{\True Gọi $S_1, S_2$  lần lượt là diện tích của hình $\left(H_1\right)$ và $\left(H_2\right)$. Khi đó $S_2>6 S_1$}
	 \loigiai{
		\begin{itemchoice}
			\itemch Diện tích của hình $\left(H_1\right)$ được  tính bởi $-\displaystyle\int_{-1}^0\left(-x^3+x^2+2 x\right)\mathrm{\,d}x$ vì hình $\left(H_1\right)$ nằm dưới trục hoành.
			\itemch Thể tích của khối tròn xoay được tạo thành khi quay hình $\left(H_2\right)$ quanh  trục $Ox$ là
			$$V=\pi \displaystyle\int_0^2\left(-x^3+x^2+2  x\right)^2\mathrm{\,d}x=\dfrac{464}{105} \pi.$$
			\itemch Tổng diện tích của hình $\left(H_1\right)$  và $\left(H_2\right)$ là
			 $$S=S_1+S_2=-\displaystyle\int_{-1}^0\left(-x^3+x^2+2 x\right)\mathrm{\,d}x+\displaystyle\int_0^2\left(-x^3+x^2+2 x\right)\mathrm{\,d}x=\dfrac{5}{12}+\dfrac{\S}{3}=\dfrac{37}{12}.$$
			\itemch $S_1=\dfrac{5}{12}, S_2=\dfrac{8}{3}  \Rightarrow S_2>6 S_1$.
		\end{itemchoice}
	 }
\end{ex}

% nguồn: Phan1_Lop_12_Dai_so_Chuong_4
\begin{ex}%[2D4V3-3]%[ma:a2fcd7a6]
	Xét tính đúng - sai  của các khẳng định sau:
	\choiceTF 
		{Thể tích của vật thể tròn xoay được tạo thành khi quay hình phẳng giới hạn bởi các đường $y=f(x)$, $x=-3$, $x=1$, $y=0$ quanh trục $Ox$ có công  thức là $V=\displaystyle\int_{-3}^{1} |f(x)|\mathrm{\,d}x$}
		{Diện tích của hình phẳng giới hạn bởi đường thẳng $y=-3x-1$ và đồ thị hàm số  $y=-3x^2-15x-1$ bằng $34$}
		{\True Quay hình phẳng giới hạn bởi các đường $y=2$; $x=0$; $x=2$; $y=0$ quanh trục $Ox$ tạo  thành vật thể tròn xoay có thể tích bằng $8\pi$}
		 {Gọi tam giác cong là hình phẳng giới hạn bởi đồ thị các hàm số $y = -\dfrac{1}{2}x-\dfrac{1}{2}$; $y = x^2 + 8x + 16$ và $y=0$. Diện tích của tam giác cong đã cho bằng $\dfrac{10}{3}$
		\begin{tikzpicture}[scale=0.8,font=\footnotesize,line join=round,line cap=round,>=stealth]
				\draw[->] (-6.1,0)--(1.1,0) node[below left] {$x$}; 
				\draw[->] (0,-1.1)--(0,3.1) node[below  left] {$y$};
				\draw (0,0) node [above left]  {$O$} circle (1pt);
				\begin{scope}
					 \clip (-6,-1) rectangle (1,3);
					\draw[samples=200,domain=-6:1,smooth,variable=\x] plot  (\x,{(-0.5*(\x)-0.5)});
					\draw[samples=200,domain=-6:1,smooth,variable=\x] plot  (\x,{((\x)^2+8*(\x)+16)});
					\fill [pattern=north east lines] (-4,0) -- plot [domain=-4:-3,samples=100] (\x,{(\x+4)*(\x+4)}) -- (-1,0) --  cycle;
				\end{scope}
			\end{tikzpicture}
		 }
	 \loigiai{
		\begin{itemchoice}
			\itemch Thể tích của vật thể tròn xoay được tạo thành khi quay hình phẳng giới hạn bởi các đường $y=f(x)$, $x=-3$, $x=1$, $y=0$ quanh trục $Ox$ là  $V=\pi\displaystyle\int_{-3}^{1} f^2(x)\mathrm{\,d}x$
			\itemch Phương  trình hoành độ giao điểm
			$$-3x-1=-3x^2-15x-1\Leftrightarrow 3x^2+12x=0\Leftrightarrow \hoac{&x=-4\\&x=0.}$$ 
			Diện tích hình phẳng  là
			\begin{eqnarray*}
				S&=&\displaystyle\int_{-4}^0 \left|(-3x-1)-(-3x^2-15x-1)\right|\mathrm{\,d}x\\ 
				 &=&\displaystyle\int_{-4}^0 \left|3x^2+12x\right|\mathrm{\,d}x\\
				&=&32. 
			\end{eqnarray*}
			 \itemch Thể tích vật tròn xoay là $V=\pi \displaystyle\int_0^2 2^2\mathrm{\,d}x=4\pi x\Bigg|_0^2=8\pi$.
			 \itemch Ta có
			\begin{itemize}
				\item  $-\dfrac{x}{2}-\dfrac{1}{2}=0\Leftrightarrow x=-1$.
				\item $x^2+8x+16=0\Leftrightarrow  x=-4$.
				\item  $x^2+8x+16=-\dfrac{x}{2}-\dfrac{1}{2}\Leftrightarrow x^2+\dfrac{17}{2}x+\dfrac{33}{2}=0\Leftrightarrow \hoac{&x=-3\\&x=-\dfrac{11}{2}.}$
			\end{itemize}
			\begin{center}
				\begin{tikzpicture}[scale=0.8,font=\footnotesize,line join=round,line cap=round,>=stealth]
						\draw[->] (-6.1,0)--(1.1,0) node[below  left] {$x$};
						 \draw[->] (0,-1.1)--(0,3.1) node[below left] {$y$};
						\draw (0,0) node [above left] {$O$} circle (1pt) 
						(-3,0)  node [below] {$-3$} circle (1pt)
						(-1,0) node [below] {$-1$} circle  (1pt)
						(-4,0) node [below] {$-4$}  circle (1pt);
						\draw[dashed]  (-3,0)--(-3,1);
						\begin{scope}
						\clip (-6,-1)  rectangle (1,3);
						\draw[samples=200,domain=-6:1,smooth,variable=\x] plot (\x,{(-0.5*(\x)-0.5)}); 
						\draw[samples=200,domain=-6:1,smooth,variable=\x] plot  (\x,{((\x)^2+8*(\x)+16)});
						\fill [pattern=north  east lines] (-4,0) -- plot [domain=-4:-3,samples=100] (\x,{(\x+4)*(\x+4)}) -- (-1,0) -- cycle;
						\end{scope}
					\end{tikzpicture}
			\end{center}
			Diện  tích tam giác cong là
			\begin{eqnarray*}
				S&=&\displaystyle\int_{-4}^{-3}(x^2+8x+16) \mathrm{\,d}x+\displaystyle\int_{-3}^{-1}  \left(-\dfrac{x}{2}-\dfrac{1}{2}\right)\mathrm{\,d}x\\
				 S&=&\dfrac{(x+4)^3}{3}\Bigg|_{-4}^{-3}+\left(-\dfrac{x^2}{4}-\dfrac{x}{2}\right)\Bigg|_{-3}^{-1}\\
				&=&\dfrac{4}{3}. 
			\end{eqnarray*}
		\end{itemchoice}
	 }
\end{ex}

% nguồn: Phan1_Lop_12_Dai_so_Chuong_4
\begin{ex}%[2D4V3-3]%[ma:1ee28d00]
	Cho các đồ thị hàm số $(C_1) \colon y=\dfrac{\sqrt{3}}{9}x^3$, cung tròn có phương trình $(C_2) \colon y=\sqrt{4-x^2}$  và hình phẳng $(H)$ được tô màu như hình vẽ.
	\begin{center}
		\begin{tikzpicture}[x=0.9cm,y=0.9cm,>=stealth,font=\footnotesize]
				\begin{scope}
					 \clip (0,0) circle (2);
					\fill[blue!80,smooth]  (0,0) -- (2,0)--
						plot[domain=2:0](\x,{{0.192*(\x)^3}})--(0,0) 
						-- cycle; 
				\end{scope}
				\draw[smooth] plot[domain=0:1](\x,{{(4-(\x)^2)^0.5}}); 
				\draw[smooth] plot[domain=1:1.5](\x,{{(4-(\x)^2)^0.5}}); 
				 \draw[smooth] plot[domain=1.5:2](\x,{{(4-(\x)^2)^0.5}});
				\draw[smooth] plot[domain=-2.18:2.49](\x,{{3^0.5/9*(\x)^3}}); 
				 \draw (2,0) node[below]{$2$};
				\draw  (0,2) node[left]{$2$};
				\draw[->,color=black] (-2.18,0) -- (4,0) node[below]  {$x$};
				 \draw[->,color=black] (0,-2) -- (0,3) node[right]{$y$};
				 \draw (-0.1,0) node[above right] {$O$};
			\end{tikzpicture}
	\end{center}
	 \choiceTF
		{\True Hình phẳng đó được giới hạn bởi các đồ thị hàm số $y=\dfrac{\sqrt{3}}{9}x^3$, $y=\sqrt{4-x^2}$,  trục hoành và các đường thẳng $x=0$, $x=2$}
		{\True  Các đường $(C_1)$ và $(C_2)$ đều đi qua điểm $A(\sqrt{3};1)$}
		{Thể tích vật thể tròn xoay được tạo thành khi quay hình phẳng giới  hạn bởi đồ thị hàm số $(C_2) \colon y=\sqrt{4-x^2}$ và hai đường thẳng $x=0$, $x=2$ quanh trục hoành là $V=\displaystyle\int_{0}^{2} \left(\sqrt{4-x^2}\right)^2\mathrm{\,d}x$}
		{\True Biết thể tích của khối tròn xoay tạo thành  khi quay $(H)$ quanh trục hoành là \break $V=\left(-\dfrac{a}{b}\sqrt{3}+\dfrac{c}{d}\right){\pi}$, trong đó $a$, $b$, $c$, $d \in \mathbb{N}^*$ và $\dfrac{a}{b}$, $\dfrac{c}{d}$ là các phân số tối giản. Khi đó $a+b+c+d=46$}
	 \loigiai{
		\begin{itemchoice}
			\itemch Hình phẳng đó được giới  hạn bởi các đồ thị hàm số $y=\dfrac{\sqrt{3}}{9}x^3$, $y=\sqrt{4-x^2}$, trục hoành và các đường thẳng $x=0$, $x=2$.
			\itemch Giao điểm $A(x;y)$  của đồ thị hàm số $(C_1) \colon y=\dfrac{\sqrt{3}}{9}x^3$ và đồ thị hàm số $(C_2) \colon y=\sqrt{4-x^2}$ thỏa mãn $$\dfrac{\sqrt{3}}{9}x^3=\sqrt{4-x^2} \Rightarrow x = \sqrt{3} \Rightarrow y=\sqrt{4-(\sqrt{3})^2}=1.$$
			 Các đường $(C_1)$ và $(C_2)$ đều đi qua điểm $A(\sqrt{3};1)$.
			\itemch Thể tích vật thể tròn xoay được tạo thành khi quay hình phẳng giới hạn bởi đồ thị hàm  số $(C_2) \colon y=\sqrt{4-x^2}$ và hai đường thẳng $x=0$, $x=2$ quanh trục hoành là $$V=\pi\displaystyle\int_{0}^{2} \left(\sqrt{4-x^2}\right)^2\mathrm{\,d}x.$$
			\itemch Thể tích của khối tròn xoay tạo thành khi quay $(H)$ quanh trục  hoành là
			\begin{align*}
				V=&\pi\displaystyle\int_{0}^{\sqrt{3}} \left(\dfrac{\sqrt{3}}{9}x^3\right)^2\mathrm{\,d}x+\pi\displaystyle\int_{\sqrt{3}}^{2} \left(\sqrt{4-x^2}\right)^2\mathrm{\,d}x\\ 
				 &=\pi\displaystyle\int_{0}^{\sqrt{3}} \left(\dfrac{1}{27}x^6\right)\mathrm{\,d}x+\pi\displaystyle\int_{\sqrt{3}}^{2} \left(4-x^2\right)\mathrm{\,d}x\\
				&=\pi \cdot  \dfrac{1}{27} \cdot \dfrac{x^7}{7} \Biggr|_{0}^{\sqrt{3}}+ \pi \cdot \left(4x-\dfrac{x^3}{3}\right)\Biggr|_{\sqrt{3}}^{2}\\
				&=\left(-\dfrac{20 \sqrt{3}}{7}+\dfrac{16}{3}\right)  \pi .
			\end{align*}
			Theo đề bài,  $V=\left(-\dfrac{a}{b}\sqrt{3}+\dfrac{c}{d}\right){\pi}$ nên $a=20$, $b=7$, $c=16$, $d=3$ suy ra $a+b+c+d=46$.
	\end{itemchoice}}
\end{ex}


%%%%%%%%%%%%%%%% TRẢ LỜI NGẮN %%%%%%%%%%%%%%%%

% nguồn: Phan2_Lop_12_Dai_so_Chuong_4
\begin{bt}%[2D4N3-3]%[ma:e527d5dc]
	Thể tích vật thể tròn xoay khi quay hình phẳng  giới hạn bởi các đường $y=\dfrac{4}{x}$, $y=0$, $x=1$, $x=4$ quanh trục $Ox$ là bao nhiêu (làm tròn kết quả đến hàng phần mười) ?
	 \shortans{37,7}
	 \loigiai{
		Thể tích của khối  tròn xoay là: \[V=\pi\displaystyle\int_1^4\left(\dfrac{4}{x}\right)^2\,\mathrm{d}x=12\pi\approx 37{,}7.\]
	} 
\end{bt}

% nguồn: Phan2_Lop_12_Dai_so_Chuong_4
\begin{bt}%[2D4N3-3]%[ma:1a568698]
	Gọi $(D)$ là hình phẳng giới hạn bởi đồ thị hàm số $y=\mathrm{e}^x-1$, trục $Ox$ và  hai đường $x=-1$, $x=0$. Tính thể tích khối tròn xoay khi quay hình phẳng $(D)$ quanh trục $O x$ (kết quả được ghi dưới dạng số thập phân, làm tròn đến hàng phần trăm).
	\shortans{0,17} 
	 \loigiai{
		Tính thể tích  khối tròn xoay khi quay hình phẳng $(D)$ quanh trục $O x$ là
		 $$V=\pi \displaystyle\int_{-1}^{0}\left(\mathrm{e}^x-1\right)^2 \mathrm{\,d}x\approx 0{,}17.$$
	} 
\end{bt}

% nguồn: Phan2_Lop_12_Dai_so_Chuong_4
\begin{bt}%[2D4N3-3]%[ma:a47740ba]
	Tính thể  tích khối tròn xoay khi quay hình phẳng $(H)$ giới hạn bởi các đường $y=\dfrac{\mathrm{e}^x}{2}$, $y=0$ và $x=-1$, $x=1$ quanh trục hoành (làm tròn kết quả đến hàng phần trăm).
	 \shortans{2,85}
	\loigiai{ 
		 Thể tích khối tròn xoay là \[V=\pi\displaystyle\int_{-1}^1\dfrac{\mathrm{e}^{2x}}{4}\,\mathrm{d}x=\left.\dfrac{\pi}{8}{\mathrm{e}^{2x}}\right|_{-1}^1=\dfrac{\pi}{8}\left(\mathrm{e}^2-\dfrac{1}{\mathrm{e}^2}\right)\approx 2{,}85.\]
	 }
\end{bt}

% nguồn: Phan2_Lop_12_Dai_so_Chuong_4
\begin{bt}%[2D4H3-3]%[ma:d1c564a4]
	Thể tích vật thể tròn xoay khi quay hình phẳng giới hạn bởi  các đường $y=1-x^2$, $y=0$ quanh trục $Ox$ có kết quả dạng $\dfrac{a\pi}{b}$ với $\dfrac{a}{b}$ là phân số tối giản. Khi đó $S=a+b$ bằng bao nhiêu?
	\shortans{31} 
	 \loigiai{
		Phương  trình hoành độ giao điểm
		\[1-x^2=0\Leftrightarrow  x=\pm1.\]
		Thể tích khối  tròn xoay là \[V=\pi\displaystyle\int_{-1}^1\left(1-x^2\right)^2\,\mathrm{d}x=\dfrac{16\pi}{15}.\]
		 Do đó $a=16$, $b=15$.\\
		 Vậy $S=a+b=31$.
	} 
\end{bt}

% nguồn: Phan2_Lop_12_Dai_so_Chuong_4
\begin{bt}%[2D4H3-3]%[ma:54bccbed]
	Một hình cầu có bán kính $6$ dm, người ta cắt bỏ hai phần bằng hai mặt phẳng song song và cùng vuông góc với đường kính để làm mặt xung quanh của một chiếc lu chứa nước (như hình vẽ). Tính thể tích $V$ (lít) mà chiếc lu chứa được biết mặt phẳng cách tâm mặt cầu $4$ dm (làm tròn đến hàng đơn vị).\begin{center}
	\begin{tikzpicture}[scale=1,font=\footnotesize,line join=round,line cap=round,>=stealth]
			\def  \x{2}
			\def  \y{0.3}
			\def \z{2} 
			 \tkzDefPoints{0/0/O,-\x/0/G,\x/0/H,-\x/\z/O',\x/\z/K}
			 \coordinate (I) at (0,0.5*\z);
			\tkzDrawCircle  (I,O')
			\coordinate (D) at ($(I)+(I)-(O)$); 
			\draw[dashed]  (G) arc (180:0:{\x} and {\y});
			\draw (G)  arc (180:360:{\x} and {\y});
			\draw[dashed] (O') arc (180:0:{\x} and  {\y});
			\draw (O') arc (180:360:{\x}  and {\y});
			\tkzCalcLength[cm](I,G) \tkzGetLength{r} 
			\coordinate (E) at  ($(I)+(0:\r)$);
			\draw[dashed] (D)--(I)-- (O) node  [midway,right]{$4$ dm}
				(I)--(E)  node [midway,above]{$6$ dm};
			 \tkzDrawPoints(O,I,D,E)
		\end{tikzpicture} \begin{tikzpicture}[scale=1,font=\footnotesize,line join=round,line cap=round,>=stealth]
		\draw[->] (-2.5,0)  -- (3,0) node[below=0.2cm] {$x$};
		\draw (-1,-.05)  -- (-1,.05) node[below]{$-4$};
		\draw  (1,-.05) -- (1,.05) node[below]{$4$};
		\draw[->] (0,-2.5) --  (0,3) node[right] {$y$};
		\begin{scope}[xshift=-1cm,yshift=0cm,line cap=round,line join=round,>=stealth,rotate=-90]
			\def \x{2} 
			\def  \y{0.3}
			 \def \z{2}
			 \tkzDefPoints{0/0/O,-\x/0/G,\x/0/H,-\x/\z/O',\x/\z/K}
			 \coordinate (I) at (0,0.5*\z);
			\tkzDrawCircle (I,O') 
			\coordinate  (D) at ($(I)+(I)-(O)$);
			\draw[dashed] (G) arc (180:0:{\x} and {\y}); 
			\draw (G) arc  (180:360:{\x} and {\y});
			\draw[dashed] (O') arc  (180:0:{\x} and {\y});
			\draw (O')  arc (180:360:{\x} and {\y});
			\tkzCalcLength[cm](I,G)  \tkzGetLength{r}
			\coordinate (E)  at ($(I)+(0:\r)$);
			 \draw[dashed] (O)--(G)
				(D)--(O'); 
			\tkzDrawPoints(O,I,D,E) 
		\end{scope}
		 \tkzDefPoints{0/0/O}
		\tkzDrawPoints(O) 
		\tkzLabelPoints[below  left](O)
	\end{tikzpicture}
\end{center}
 \shortans{771}
 \loigiai{
	Thể tích cần tìm là thể tích của khối tròn xoay khi quay hình phẳng $(H)$ giới hạn bởi đồ thị hàm số $f(x)=\sqrt{36-x^2}$, trục hoành và các đường thẳng $x=-4$, $x=4$ quanh trục  hoành.\\
	Do đó $V=\pi\displaystyle\int_{-4}^4\left(36-x^2\right)dx=\dfrac{736\pi}{3}\approx 771$ lít. 
 }
\end{bt}

% nguồn: Phan2_Lop_12_Dai_so_Chuong_4
\begin{bt}%[2D4H3-3]%[ma:8b44da9a]
	\immini{Cho khối tròn xoay như hình bên. Tính thể tích của khối tròn xoay được tạo thành bởi hình phẳng cho ở hình đã cho khi quay quanh trục $Ox$ (viết  kết quả dưới dạng số thập phân và làm tròn đến hàng phần mười).}
	{\begin{tikzpicture}[line join=round,line cap=round,>=stealth]
		\draw[->] (-0.6,0)--(1.6,0) node[below left]  {$x$};
		\foreach \x/\nx  in {1/1}
		\draw[thin] (\x,1pt)--(\x,-1pt)  node [below] {$\nx$};
		 \draw[->] (0,-1.3)--(0,1.3) node[below left] {$y$};
		\fill (0,0)  node [below left] {$O$};
		\draw (1,0)--(1,1); 
		\begin{scope}
			\clip (-0.5,-1.5) rectangle  (1.5,1.5);
			\draw[samples=100,domain=0:1,smooth,variable=\x] plot (\x,{sqrt  ((\x))}) node [above]{$y=\sqrt{x}$};
			\draw[dashed,samples=100,domain=0:1,smooth,variable=\x] plot (\x,{-sqrt ((\x))}); 
		\end{scope}
		 \draw [dashed] (1,1) arc (90:-90:0.2 cm and 1 cm);
		\draw [dashed] (1,1) arc (90:270:0.2 cm and  1 cm);
	\end{tikzpicture}}
	 \shortans{1,6}
	\loigiai{ 
		Hình phẳng đã cho được giới hạn bởi đồ thị  hàm số $y=\sqrt{x}$, trục hoành và các đường thẳng $x=0$, $x=1$, khi quay hình phẳng đó quanh trục $O x$ ta được khối tròn xoay như hình đã cho. \\
		Thể tích khối tròn  xoay đó là
		 $$
		V=\pi \displaystyle\int_0^1\left(\sqrt{x}\right)  ^2 \mathrm{\,d} x=\pi\displaystyle \int_0^1 x \mathrm{\,d} x=\pi \cdot \left.\dfrac{x^2}{2}\right|_0 ^1=\dfrac{\pi}{2} \approx 1{,}6.
		 $$
	} 
\end{bt}

% nguồn: Phan2_Lop_12_Dai_so_Chuong_4
\begin{bt}%[2D4H3-3]%[ma:97180d4c]
	Người ta tạo ra mô hình một quả trứng ngỗng bằng cách quay hình phẳng giới hạn bởi đồ thị hàm số $y=\dfrac{1}{30} \sqrt{7569-400x^2}$ và trục hoành với $-4{,}35 \le x \le 4{,}35$ quanh trục $Ox$. Sử dụng mô hình này để tính thể tích quả trứng theo đơn vị cm$^3$ và làm tròn đến hàng đơn vị (với $x$, $y$ được đo theo  đơn vị cm).
	 \shortans{153}
	\loigiai{Thể tích quả trứng là thể tích của khối tròn xoay được tạo thành khi quay hình phẳng tạo bởi đồ thị của hàm số $y=\dfrac{1}{30} \sqrt{7569-400x^2}$, trục hoành và hai đường thẳng $x=-4{,}35$ và $x=-4{,}35$ quanh trục hoành. Do  đó
	$$V=\pi \displaystyle \int_{-4{,}35}^{4{,}35} \left(\dfrac{1}{30} \sqrt{7569-400x^2}\right)^2 \mathrm{\,d}x=\dfrac{24389 \pi}{500} \approx 153 \;  \text{ (cm$^3$)}.$$}
\end{bt}

% nguồn: Phan2_Lop_12_Dai_so_Chuong_4
\begin{bt}%[2D4H3-3]%[ma:6d56bd92]
	\immini{ 
	 Biết thể tích khối tròn xoay tạo nên khi quay xung quanh trục $Ox$ hình phẳng giới hạn bởi các đường $y=(1-x)^2, y=0, x=0, x=2$ có dạng $\dfrac{a \pi}{b}$ với $\dfrac{a}{b}$ là phân số tối giản. Tính giá trị biểu thức $S=24a+12b$.}
	 {
	\begin{tikzpicture}[line cap=butt,line join=miter,>=stealth]
			 \tikzset{declare function={xmin=-1; xmax=2.5; ymin=-1.5; ymax=1.5;
				 f (\x)=1*(1-\x)^2;
				 },
				 smooth,samples=450
			} 
			\draw[->] (xmin,0)--(xmax,0) node[shift={(-100:7pt)},font=\normalsize]{$ x  $};
			\draw[->] (0,ymin)--(0,ymax)  node[shift={(190:7pt)},font=\normalsize]{$ y $};
			\fill (0,0) node[shift={(225:7pt)},font=\normalsize]{$  O $};
			\fill [pattern=north west  lines] plot[domain=0:2] (\x,{f (\x)})--plot[domain=2:0] (\x,{0});
			 \draw[dashed] (2,0) node[below]{$2$}--(2,1);
			\clip  (xmin,ymin) rectangle (xmax,ymax);
			 \draw plot[domain=0:2] (\x,{f (\x)});
			\draw[dashed] plot[domain=0:2] (\x,{-f  (\x)});
			 \draw (0,1)
				..controls  +(180:.23) and +(180:.23)..(0,-1);
			 \draw[dashed] (0,1)
				..controls  +(0:.23) and +(0:.23)..(0,-1);
			\draw  (2,1)
				..controls +(180:.23)  and +(180:.23)..(2,-1);
			\draw  (2,1)
				..controls +(0:.23) and +(0:.23)..(2,-1); 
		\end{tikzpicture}
	 }
	 \shortans{108}
	\loigiai{ 
		Ta  có $V=\pi\displaystyle\int_0^2(1-x)^4\mathrm{\,d}x=\dfrac{2\pi}{5}$.\\
		Suy ra $a=2$, $b=5$.\\ 
		 Do đó $24a+12b=24\cdot 2+12\cdot 5=108$.
	 }
\end{bt}

% nguồn: Phan2_Lop_12_Dai_so_Chuong_4
\begin{bt}%[2D4H3-3]%[ma:ae301dc0]
	Biết thể tích khối tròn xoay tạo nên khi quay xung quanh trục $Ox$ hình phẳng giới hạn bởi các đường $y=\left(1-x\right)^2$, $y=0$, $x=0$, $x=2$ có dạng  $\dfrac{a\pi}{b}$ với $\dfrac{a}{b}$ là phân số tối giản. Tính giá trị biểu thức $S=24a+12b$.
	\shortans{108} 
	 \loigiai{
		Thể tích khối tròn xoay  là \[V=\pi\displaystyle\int_0^2\left(1-x\right)^4\,\mathrm{d}x=\dfrac{2\pi}{5}.\]
		Do đó $a=2$, $b=5$.  \\
		 Vậy $S=24a+12b=108$.
	 }
\end{bt}

% nguồn: Phan2_Lop_12_Dai_so_Chuong_4
\begin{bt}%[2D4V3-3]%[ma:0c73eaea]
	Từ một tấm tôn hình chữ nhật $ABCD$ với $AB=30$ cm, $AD=\dfrac{55\pi}{3}$ cm. Người ta cắt tấm tôn theo đường hình sin như hình vẽ bên để được hai miếng tôn. Biết $AM=20$ cm, $CN=15$ cm, $BE=5\pi $ cm. Tính thể tích $V$ (đơn vị: lít) của khối tròn xoay  được tạo thành khi xoay miếng tôn lớn quanh trục $AD$.
	\begin{center}
		\begin{tikzpicture}
				\begin{axis}[
						 axis lines=middle,
						 xlabel={$x$},ylabel={$y$},
						xtick={0,5.76}, 
						xticklabels={$0$,$\dfrac{55\pi}{30}$}, 
						ytick={0,3}, 
						 yticklabels={$0$,$3$},
						 domain=0:5.76,
						 samples=200,
						 ymin=0,ymax=4,
						 xmin=-0.3,xmax=6.7,
						axis  line style=->,
						 tick label style={font=\small},
						 clip=false]
					 \addplot [
						domain=0:5.76, 
						fill=gray!30] 
						 {sin (deg (x)) + 2} \closedcycle;
					\addplot {sin (deg (x)) +  2};
					\node[rotate=-40]  at (axis cs:2.2,2.3) {\small $y=\sin x + 2$};
					\draw (axis cs:0,0)  -- (axis cs:0,3) -- (axis cs:5.76,3) node[right] {$C$}
						-- (axis cs:5.76,0) node[above right]  {$D$}
						-- cycle; 
					\node[above left] at (axis  cs:0,0) {$A$};
					\node[left] at (axis cs:0,1.5)  {$M$};
					\node[above right] at (axis cs:0,3)  {$B$};
					\node[above] at  (axis cs:2.2,3) {$E$};
					\node[right] at  (axis cs:5.76,1.5) {$N$};
					 \node[below] at (axis cs:0,0) {$O$};
				\end{axis}
			\end{tikzpicture}
	\end{center}
	\shortans{83,8} 
	 \loigiai{
		Chọn  hệ trục $Oxy$ với đơn vị trên trục là $10$ cm như hình sau.\\
		\begin{center}
			\begin{tikzpicture}
					\begin{axis}[
						axis  lines=middle,
						xlabel={$x$},ylabel={$y$}, 
						xtick={0,5.76}, 
						xticklabels={$0$,$\dfrac{55\pi}{30}$}, 
						ytick={0,3}, 
						yticklabels={$0$,$3$}, 
						 domain=0:5.76,
						 samples=200,
						ymin=0,ymax=4, 
						xmin=-0.3,xmax=6.7, 
						 axis line style=->,
						tick label  style={font=\small},
						clip=false] 
						\addplot  [
						 domain=0:5.76,
						fill=gray!30] 
						{sin (deg (x)) + 2}  \closedcycle;
						\addplot {sin (deg (x)) + 2}; 
						\node[rotate=-40] at  (axis cs:2.2,2.3) {\small $y=\sin x + 2$};
						\draw (axis  cs:0,0) -- (axis cs:0,3) -- (axis cs:5.76,3) node[right] {$C$}
						-- (axis  cs:5.76,0) node[above right] {$D$}
						 -- cycle;
						\node[above left] at (axis cs:0,0)  {$A$};
						\node[left] at (axis cs:0,1.5)  {$M$};
						\node[above  right] at (axis cs:0,3) {$B$};
						\node[above] at (axis cs:2.2,3) {$E$}; 
						\node[right] at  (axis cs:5.76,1.5) {$N$};
						\node[below] at (axis cs:0,0)  {$O$};
					\end{axis}
				\end{tikzpicture}
		\end{center}
		Khi đó đường cắt miếng tôn hình sin trên là đồ thị của hàm số $y=\sin x+2$, $0\le  x\le\dfrac{55\pi}{30}$.\\
		 Thể tích cần tìm là $V=\pi\displaystyle\int_0^{\tfrac{55\pi}{30}}{\left(\sin x+2\right)^2\,\mathrm{d}x}\approx 83{,}8$ lít.
	 }
\end{bt}

% nguồn: Phan2_Lop_12_Dai_so_Chuong_4
\begin{bt}%[2D4V3-3]%[ma:cda65172]
	 \immini
	{ 
	Một vật trang trí có dạng khối tròn xoay tạo thành khi quay miền $(R)$ được giới hạn bởi đường gấp khúc $DABFE$ và cung tròn $ED$ (phần gạch chéo trong hình bên) xung  quanh trục $AB$. Biết $ABCD$ là hình chữ nhật cạnh $AB =3$ cm, $AD=2$ cm; $F$ là trung điểm của $BC$; điểm $E$ cách $AD$ một đoạn bằng $1$ cm. Tính thể tích của vật trang trí đó, làm tròn kết quả đến hàng phần mười, đơn vị cm$^3$.
	}{\resizebox{!}{5cm}{ 
	\begin{tikzpicture}
			\path (0,0)  coordinate (B)
				 (2,0) coordinate (C)
				(2,3) coordinate  (D)
				(0,3)  coordinate (A)
				($(B)!.5!(C)$) coordinate (F) 
				($(F)+(0,2)$)  coordinate (E)
				($(A)!.5!(D)$) coordinate (I); 
			\fill[cyan!25] (A)--(B)--(F)--(I)--cycle (E) arc  (-90:0:1cm)--(D)--(I)--cycle;
			\draw  (A)--(B)--(C)--(D)--cycle (E)--(F);
			\draw[dashed]  (E)--(I);
			\draw  (E) arc (-90:0:1cm);
			\foreach  \d/\g in {A/90,B/-90,C/-90,D/90,F/-90,E/-60}
			\draw[fill] (\d) circle  (0.5pt) node[shift={(\g:7pt)},font=\tiny]{$\d$};
		\end{tikzpicture}}
	 }
	 \shortans{16,5}
	\loigiai{ 
	\immini{ 
		Gắn hệ trục  tọa độ như hình vẽ bên.\\
		Quay hình chữ nhật $GEFB$ xung quanh trục $AB$ ta được khối trụ có bán kính đáy  bằng $1$ và chiều cao bằng $2$.\\
		Phương trình đường tròn tâm $I$ bán kính bằng $1$ là  $x^2+(y-1)^2=1$, suy ra phương trình của cung $DE$ là $y=1+\sqrt{1-x^2}$.\\
		Vây thể tích của  vật trang trí là
		 $$V=\pi \displaystyle\int_{0}^{1} \left(1+\sqrt{1-x^2}\right)^2 \mathrm{\,d}x+\pi \cdot 1^2\cdot 2\approx 16{,}5~(\mathrm{cm}^3).$$
		} 
		{\resizebox{0.35\linewidth}{!}{\begin{tikzpicture}
		\draw[->]  (-0.5,0)--(4,0) node[below]{$x$};
		\draw[->] (0,-0.5)--(0,3) node[left]{$y$}; 
		 \path (0,0) coordinate (A)
			(3,0) coordinate (B) 
			(3,2) coordinate  (C)
			(0,2) coordinate (D) 
			($(B)!.5!(C)$) coordinate (F) 
			 ($(F)-(2,0)$) coordinate (E)
			($(A)!.5!(D)$) coordinate  (I)
			(1,0) coordinate (G); 
		\fill[cyan!25] (A)--(B)--(F)--(I)--cycle  (E) arc (0:90:1cm)--(D)--(I)--cycle;
		\draw (A)--(B)--(C)--(D)--cycle (E)--(F); 
		\draw[dashed]  (E)--(I)
			 (E)--(G);
		 \draw (E) arc (0:90:1cm);
		\foreach \d/\g in  {A/-120,B/-90,C/45,D/180,F/0,E/45,I/180,G/-90}
		\draw[fill] (\d) circle  (0.5pt) node[shift={(\g:7pt)},font=\tiny]{$\d$};
	\end{tikzpicture}}
		 }
	 }
\end{bt}

% nguồn: Phan2_Lop_12_Dai_so_Chuong_4
\begin{bt}%[NBV2526]%[2D4V3-3]%[ma:c8fb7ce5]
	Để trang trí một bức tường hình chữ nhật có kích thước $3$\,m $\times 4$\, m trong phòng, bạn Hoa vẽ lên tường một hình như sau. Trên mỗi cạnh của  một hình lục giác đều có cạnh $2$\,m, vẽ một cánh hoa hình parabol. Đỉnh của parabol cách cạnh $3$\,dm và nằm phía ngoài hình lục giác đều. Đường parabol đi qua hai đầu mút của mỗi cạnh hình lục giác (tham khảo hình vẽ).
	\begin{center}
		\begin{tikzpicture}[line join=round,font=\footnotesize,scale=.2]
				\tikzset{declare function={f  (\x)=-3*(\x)^2+3;},smooth,samples=50}
				 \draw[cyan,fill=cyan!10] (15,-15) rectangle (55,15);
				\foreach \i  in {0,...,5}{
					\draw[scale=2,rotate=-90+60*\i,shift={(0,sqrt 3)}]  plot[domain=-1:1] (\x,{f (\x)})--cycle;
					\draw[xshift=26.5cm,scale=1,rotate=-90+60*\i,shift={(0,sqrt 3)}] plot[domain=-1:1] (\x,{f (\x)})--cycle; 
					 \draw[xshift=43.5cm,scale=1,rotate=-90+60*\i,shift={(0,sqrt 3)}] plot[domain=-1:1] (\x,{f (\x)})--cycle;
				} 
			\end{tikzpicture}
	\end{center}
	Hỏi bạn Hoa có thể vẽ tối đa bao nhiêu  hình có cùng kích thước như vậy lên bức tường cần trang trí?
	\shortans{34} 
	 \loigiai{
		Kích thước bức tường là $3\times 4$, diện tích là $S_{\text{ tường}} = 3\cdot  4 = 12\,\text{ m}^{2} = 1\,200\, \text{ dm}^{2}$.\\
		Diện tích của hình lục giác đều bằng $2$\,dm là  $6 \cdot \dfrac{2^2\sqrt{3}}{4} = 6\sqrt{3}$. \\
		Gắn trục tọa độ $Oxy$ cho một cánh hoa. Diện tích của một cánh hoa là diện tích hình phẳng giới hạn  bởi đồ thị hàm số $(P) \colon y = ax^2 + bx + c$ sao cho $(P)$ có đỉnh $I(0;3)$ và cắt trục hoành tại hai điểm $(-1;0)$ và $(1;0)$.
		\begin{center}
			\begin{tikzpicture}[>=stealth,line join=round,font=\footnotesize,scale=1]
					\tikzset{declare  function={f (\x)=-3*(\x)^2+3;},smooth,samples=50}
					\draw[->] (-1.6,0)--(2.3,0)  node[pos=.97,above]{$x$};
					\draw[->](0,-.7)--(0,3.7) node[pos=.95,left]{$y$}; 
					\draw plot[domain=-1.1:1.1] (\x,{f (\x)}); 
					\fill (-1,0) circle  (1pt) node[above left=-2pt]{$-1$}
						(1,0) circle (1pt) node[above right=-2pt]{$1$} 
						(0,0)  node[shift={(-135:6pt)}]{$O$}
						(0,3) circle (1pt) node[above  right=-2pt]{$3$};
				\end{tikzpicture}
		\end{center}
		Phương trình parabol  có dạng $y=ax^{2}+c$. Đi qua $I(0;3) \Rightarrow c=3$. Đi qua $(1;0) \Rightarrow 0 = a(1)^{2} + 3 \Rightarrow a=-3$.
		Phương trình parabol là $(P)\colon  y=-3x^{2}+3$.\\
		Diện tích một cánh hoa (phần giới hạn bởi parabol và trục hoành  $y=0$)
		$$S_{\text{ cánh  hoa}} = \displaystyle\int_{-1}^{1}|-3x^{2}+3|\mathrm{\,d}x = 4\,\text{ dm}^2.$$
		Diện tích một bông hoa được trang trí là 
		$$S_{\text{ hoa}} = S_{\text{ lục giác}} + 6 \cdot S_{\text{ cánh hoa}} = 6\sqrt{3} + 6 \cdot 4 = 24  + 6\sqrt{3}\,\text{ dm}^2.$$
		Diện tích bức  tường hình chữ nhật có kích thước $3$\,m$\times 4$\,m là $3 \cdot 4 = 12$\,m$^2$ = $1\,200$ \text{ dm}$^2$.\\
		Tỉ lệ diện tích của bức tường hình chữ nhật  và diện tích một bông hoa là
		\[\dfrac{1200}{24 + 6\sqrt{3}} \approx 34{,}89.  \]
		Vậy bạn  Hoa có thể trang trí tối đa $34$ hình có cùng kích thước như vậy lên bức tường cần trang
		 trí.
	} 
\end{bt}

% nguồn: Phan2_Lop_12_Dai_so_Chuong_4
\begin{bt}%[2D4V3-3]%[ma:2e1fa3b3]
	Một quả trứng khủng long đồ chơi bằng nhựa có thiết diện qua trục lớn là một đường elip. Biết độ dài mỗi trục là $12$ cm và  $8$ cm.
	Bên trong quả trứng  người ta cần thiết kế một chiếc hộp hình trụ để đựng các đồ chơi trẻ con như bóng đèn xanh đỏ, kẹo $\ldots$
	Hỏi khối trụ như  thế có thể tích tối đa bao nhiêu cm$^3$ (làm tròn đến hàng đơn~vị)?
	\begin{center}
		\begin{tikzpicture}[line join=round,line cap=round,>=stealth]
				\path  (0,0) coordinate (O);
				\draw (O) ellipse  (3 cm and 2 cm);
				 \path ($(O) + (40:3 and 2)$) coordinate (A)
					($(O) + (140:3 and 2)$)  coordinate (B)
					($(O) +  (220:3 and 2)$) coordinate (C)
					($(O) +  (-40:3 and 2)$) coordinate (D);
				 \draw[dashed] (A)--(B)
					 (C)--(D)
				(A) to [bend right=20](D)
			(B) to [bend right=20](C);
		\draw (A) to [bend left=20](D)
			(B) to [bend left=20](C);
		\foreach \t/\g  in {A/90,B/90,C/180,D/0}{
			\draw (\t) circle (1pt); 
		 }
		\end{tikzpicture}
	\end{center}
	\shortans{232} 
	\loigiai{ 
		\begin{center}
			\begin{tikzpicture}[line join=round,line cap=round,>=stealth]
					\path  (0,0) coordinate (O);
					\draw (O) ellipse (3 cm and 2 cm); 
					\path ($(O) +  (40:3 and 2)$) coordinate (A)
						($(O)  + (140:3 and 2)$) coordinate (B)
						($(O)  + (220:3 and 2)$) coordinate (C)
						($(O) + (-40:3  and 2)$) coordinate (D);
					 \draw[->] (-4,0)--(4,0) node[below]{$x$};
					\draw[->]  (0,-3)--(0,3) node[left]{$y$};
					\draw[dashed] (A)--(B) 
						 (C)--(D)
					(A) to [bend right=20](D)
				(B) to [bend right=20](C);
			\draw (A) to [bend left=20](D)
				(B) to [bend left=20](C);
			 \foreach \t/\g in {A/90,B/90,C/180,D/0}{
				\draw (\t) circle (1pt); 
			 }
			\foreach  \t/\g in {O/230,A/90}{
				\draw (\t) circle (1pt)  node[shift={(\g:7pt)},font=\scriptsize]{$ \t $};
			 }
			\end{tikzpicture}
		\end{center}
		Gắn  elip lên hệ trục tọa độ $Oxy$ như hình vẽ, elip có \[2a=12\Rightarrow a=6; 2b=8\Rightarrow b=4.\]
		Phương trình chính tắc  elip $(E)\colon\dfrac{x^2}{36}+\dfrac{y^2}{16}=1$.\\
		Đặt chiều cao và bán kính đáy hình trụ nội tiếp elip là $h$, $r$ thì điểm tiếp xúc $A\left(\dfrac{h}{2}; r\right)$ với $0< h < 12$; $0< r < 4$.\\ 
		 Điểm $A$ thuộc $(E)\colon\dfrac{x^2}{36}+\dfrac{y^2}{16}=1\Rightarrow \dfrac{\left(\dfrac{h}{2} \right)^2}{36}+\dfrac{r^2}{16}=1\\
		\Rightarrow \dfrac{h^2}{144}+\dfrac{r^2}{16}=1\Rightarrow  r^2=16\left(1-\dfrac{h^2}{144}\right)$.\\
		Thể tích khối trụ là $V=\pi  r^2h=\pi \cdot16\left(1-\dfrac{h^2}{144} \right)\cdot h=\pi \left(16h-\dfrac{h^3}{9} \right)$ \\
		hay $V=\pi \left(16h-\dfrac{h^3}{9} \right)$.\\ 
		Ta có $V'=\pi  \left(16-\dfrac{h^2}{3} \right)$; \\
		$V'=0\Rightarrow 16-\dfrac{h^2}{3}=0\Rightarrow  h^2=48\Rightarrow h=4\sqrt{3} \in (0; 12)$.\\
		Bảng biến thiên 
		\begin{center}
			\begin{tikzpicture}[font=\normalsize,t style/.style={style=solid}]
					 \tkzTabInit[nocadre=true,lgt=1.2,espcl=2.5,deltacl=0.5]
						{$h$/0.7,$V'$/0.7,$V$/2} 
						{$0$,$4\sqrt{3}$,$12$} 
					\tkzTabLine{,+,0,-} 
					 \path
						($(N13)!0.1!(N12)$)  node (A1){}
						($(N23)!0.8!(N22)$)  node (A2){}
						($(N33)!0.1!(N32)$) node  (A3){};
					\foreach \x/\y in {A1/A2,A2/A3}{ 
						\draw[-stealth]  (\x)--(\y);
					} 
				\end{tikzpicture}
		\end{center}
		Giá trị lớn  nhất của thể tích khối trụ là \[V_{\max}=V\left(4\sqrt{3} \right)=\pi \left(16\cdot 4\sqrt{3}-\dfrac{\left(4\sqrt{3} \right)^3}{9} \right)\approx 232~cm^3.\]
	} 
\end{bt}

% nguồn: Phan2_Lop_12_Dai_so_Chuong_4
\begin{bt}%[2D4V3-3]%[ma:71a48704]
	Cho hai nửa đường tròn như hình vẽ bên dưới, trong đó đường kính của nửa đường tròn lớn gấp đôi đường kính của nửa đường tròn nhỏ. Biết rằng nửa hình tròn đường kính $AB$  có diện tích $8\pi$ và $\widehat{ABC}=60^\circ$.
	Tính thể tích vật thể tròn xoay tạo thành khi quay hình phẳng $(H)$ (phần được tô đậm) quanh đường thẳng $AB$? kết quả được làm tròn đến hàng phần trăm. 
	\begin{center}
		\begin{tikzpicture}[declare function={r=3;}]
				\foreach  \x/\y/\ten in {0/0/O,8/0/B,4/0/I}{
					\path (\x,\y) coordinate  (\ten); }
				\path ($(O)!0.5!(I)$)  coordinate (T);
				 \pgfmathanglebetweenpoints{\pgfpointanchor{T}{center}}{\pgfpointanchor{O}{center}}
				\let\StartAngle\pgfmathresult 
				\draw[teal] (O) let \p1 = ($(T) - (I)$), 
				 \n1 = {veclen (\x1,\y1)} in arc (\StartAngle:{\StartAngle-180}:{\n1});
				\path ($(O)!0.5!(B)$)  coordinate (L);
				 \pgfmathanglebetweenpoints{\pgfpointanchor{L}{center}}{\pgfpointanchor{O}{center}}
				\let\StartAngle\pgfmathresult 
				\draw[teal] (O) let \p1 = ($(L)  - (B)$),
				\n1 = {veclen (\x1,\y1)} in arc  (\StartAngle:{\StartAngle-180}:{\n1});
				\tikzset{declare function={f  (\x)=(\x)/3^(1/2);}}
				\begin{scope}
					\clip (0,0) rectangle  (6,5);
					\draw[teal,samples=100]  plot[domain=0:8] (\x,{f (\x)});
				\end{scope}
				\path (6,{2*sqrt (3)}) coordinate (C); 
				\path[name  path=dthang] (O)--(C);
				\path[name path=dtron]  let \p1=($(T) - (O)$) in (T) circle ({veclen (\x1,\y1)});
				 \path[name intersections={of=dtron and dthang,by={D}}];
				\fill[cyan!30,opacity=0.7] 
					(D) --  (C)
				arc[start angle=60,end angle=0,radius=4]
				-- (4,0) 
			arc[start angle=0,end angle=60,radius=2]
			 -- cycle;
		 \draw[teal] (O)--(B)
			(6,1.3)  node{$(H)$};
		\foreach \t/\goc in {B/-90,I/-90,C/90,D/90}{ 
			 \draw[teal] (\t) circle (1pt) node[shift={(\goc:9pt)}] {$\t$};
		 }
		\draw[teal] (0,0) circle (1pt)  node[shift={(180:9pt)}] {$A$};
		\end{tikzpicture}
	\end{center}
	 \shortans{85,9}
	 \loigiai{
		\begin{center}
			\begin{tikzpicture}[declare function={r=3;}]
					 \foreach \x/\y/\ten in {0/0/O,8/0/B,4/0/I}{
						\path (\x,\y) coordinate (\ten);  }
					\path ($(O)!0.5!(I)$) coordinate  (T);
					\pgfmathanglebetweenpoints{\pgfpointanchor{T}{center}}{\pgfpointanchor{O}{center}} 
					 \let\StartAngle\pgfmathresult
					\draw[teal] (O) let \p1 = ($(T)  - (I)$),
					\n1  = {veclen (\x1,\y1)} in arc (\StartAngle:{\StartAngle-180}:{\n1});
					\path ($(O)!0.5!(B)$) coordinate  (L);
					 \pgfmathanglebetweenpoints{\pgfpointanchor{L}{center}}{\pgfpointanchor{O}{center}}
					\let\StartAngle\pgfmathresult 
					\draw[teal] (O) let \p1 = ($(L) - (B)$), 
					\n1 = {veclen (\x1,\y1)} in arc (\StartAngle:{\StartAngle-180}:{\n1}); 
					\tikzset{declare  function={f (\x)=(\x)/3^(1/2);}}
					\begin{scope}
						\clip (0,0)  rectangle (6,5);
						\draw[teal,samples=100] plot[domain=0:8] (\x,{f (\x)}); 
					\end{scope}
					\path  (6,{2*sqrt (3)}) coordinate (C);
					\path[name  path=dthang] (O)--(C);
					\path[name path=dtron] let \p1=($(T)  - (O)$) in (T) circle ({veclen (\x1,\y1)});
					\path[name  intersections={of=dtron and dthang,by={D}}];
					 \fill[cyan!30,opacity=0.7]
						(D)  -- (C)
					arc[start angle=60,end angle=0,radius=4]
					 -- (4,0)
				arc[start angle=0,end angle=60,radius=2]
				 -- cycle;
			 \draw[teal] (O)--(B)
				(D)--(C)  node[midway,above,rotate = 30]{$y=\frac{\sqrt{3}}{3}x$}
				(6,1.3) node{$(H)$}; 
			\foreach \t/\goc in  {B/-90,I/-90,C/90,D/90}{
				\draw[teal] (\t) circle (1pt) node[shift={(\goc:15pt)}]  {$\t$};
			 }
			 \draw[teal,->] (-1,0)--(9,0)
				node[shift={(-100:7pt)},font=\normalsize]{$x$}; 
			\draw[teal,->] (0,-1)--(0,5) 
				 node[shift={(180:7pt)},font=\normalsize]{$y$};
			\draw[teal]  (5pt,0) |- (0,5pt)
				(0,0)  node[shift={(225:9pt)},font=\normalsize]{$O$};
			\foreach \x in  {1,2,3,4,5,6,7,8}{
				 \path ($(\x,0)+(-90:5pt)$) node[font=\scriptsize]{$\x$};
			 }
			\end{tikzpicture}
		\end{center}
		Diện  tích nửa đường tròn đường kính $AB$ là $\dfrac{1}{2}\pi\left(\dfrac{AB}{2}\right)^2=8\pi \Rightarrow AB=8$.\\
		Xét hệ trục tọa độ $Oxy$ như hình vẽ với $O\equiv A$ và tia $AB$ trùng với tia $Ox$.\\ 
		Đường thẳng $AC$  có phương trình là $d\colon y=\dfrac{\sqrt{3}}{3}x$ (vì $AC$ đi qua $A(0;0)$ và có hệ số góc $k=\tan \widehat{BAC}=\tan 30^{\circ}=\dfrac{\sqrt{3}}{3}$).\\
		Gọi  $(C)$ là đường tròn tâm $(2;0)$, bán kính $R=2$. Khi đó phương trình $(C)\colon (x-2)^2+y^2=4$.\\
		Hoành độ giao điểm giữa $d$ và đường tròn $(C)$ thỏa mãn phương trình \[(x-2)^2+\left(\dfrac{\sqrt{3}}{3}x\right)^2=4\Leftrightarrow \dfrac{4}{3}x^2-4x=0\Leftrightarrow  \hoac{x&=0\\ x&=3.}\]
		Đường thẳng $BC$ qua $B(8;0)$ và vuông góc  với $d\colon y=\dfrac{\sqrt{3}}{3}x$ nên có phương trình $y=-\sqrt{3}(x-8)$.\\
		Hai đường thẳng $AC$, $BC$ cắt nhau tại $C$ thỏa hệ $  \heva{y&=\dfrac{\sqrt{3}}{3}x \\ y&=-\sqrt{3}(x-8) } \Rightarrow C(6; 2\sqrt{3}).$\\
		Thể tích khối tròn xoay khi quay tam  giác $ABC$ quanh $Ox$ là \[V=\pi\displaystyle\int_3^4\left[\dfrac{1}{3}x^2-(4-(x-2)^2)\right]\mathrm{d}x+\pi\displaystyle\int_4^6\dfrac{1}{3}x^2\mathrm{d}x+\pi\displaystyle\int_6^83(x-8)^2\mathrm{d}x=\dfrac{82}{3}\pi \approx 85{,}9.\]
	} 
\end{bt}

% nguồn: Phan2_Lop_12_Dai_so_Chuong_4
\begin{bt}%[2D4V3-3]%[ma:e4ba0740]
	 \immini[thm]{
	Một chậu cây có chiều cao là $30$ cm và đường kính miệng chậu là $30$ cm. Mặt cắt ngang của chậu cây là một đường parabol (tham khảo hình vẽ). Tính thể tích của chậu cây đó (đơn  vị: $\mathrm{dm}^3$, kết quả làm tròn đến hàng phần chục).
	 \shortans{10,6}
	}{ 
	\begin{tikzpicture}[line join=round,line cap=round,>=stealth,scale=1]
			\tikzset{label  style/.style={font=\footnotesize}}
			\draw[dashed,thin](-2.5,0)--(2.5,0); 
			 \draw[dashed,thin](-2.5,4)--(2.5,4);
			\draw[dashed,thin](0,-0.5)--(0,5); 
			 \draw[<->](2.2,0)--(2.2,4) node[midway,right]{$30$};
			\draw[|<->|](-2,5)--(2,5) node[midway,sloped,above]{$30$}; 
			\begin{scope}
				\clip  (-2.5,-0.5) rectangle (2.5,4.5);
				\draw (-2,4) plot[samples=200,domain=-2:2,smooth,variable=\x] (\x,{1*(\x)^2+0*(\x)+0}) arc (0:360:2  cm and 0.5 cm);
			\end{scope}
		\end{tikzpicture}
	} 
	\loigiai{ 
		Vì parabol đi qua gốc tộ độ $O$  nên $(P):y=ax^2$.\\
		 Vì parabol đi qua điểm $(-15,30)$ và $(15,30)$ nên $a=\dfrac{y}{x^2}=\dfrac{30}{15^2}=\dfrac{2}{15}$.\\
		Suy ra parabol  $(P):y=\dfrac{2}{15}x^2$.\\
		Vì vật  thể quay quanh trục $Oy$ nên $(P):x=\sqrt{\dfrac{15}{2}y}$.\\
		Thể tích chậu  cây là
		\[V=\pi \displaystyle\int_0^{30}\left(\dfrac{15}{2} y\right) \mathrm{\,d}y=3375\pi\ \left(\mathrm{cm}^3\right)\approx  10{,}6\ \left(\mathrm{dm}^3\right). \]
	 }
\end{bt}

% nguồn: Phan2_Lop_12_Dai_so_Chuong_4
\begin{bt}%[NBV2526]%[2D4V3-3]%[ma:72461eaf]
	Kiến trúc sư thiết kế một khu sinh hoạt cộng đồng có dạng hình vuông $ABCD$ có độ dài đường chéo $AC=45$\,m. Trong đó, phần được tô màu đậm là sân chơi, phần còn lại để trồng hoa. Mỗi phần trồng hoa có đường biên cong là một phần của parabol với các đỉnh thuộc một trục đối xứng của hình vuông, khoảng cách từ đỉnh đó đến đỉnh tương ứng của hình vuông bằng $15$\,m. Bản thiết kế khu sinh hoạt nhận tâm hình vuông $ABCD$ làm tâm đối xứng và trên hình bên ta có $AM=MN=NB$. Để trang trí cho phần sân chơi tốn $500$ nghìn đồng $/1$\,m$^2$, để trồng hoa tốn $200$ nghìn đồng $/1$\,m$^2$. Tổng số tiền cần dùng là bao nhiêu triệu đồng? (Kết quả làm tròn đến  hàng đơn vị).
	\begin{center}
		\begin{tikzpicture}[>=stealth,line join=round,scale=.2]
				 \tikzset{declare function={f (\x)=2/15*(\x)^2+15/2;},smooth,samples=50,dot/.style={circle,inner sep=1pt}}
				\fill[gray!20] (-22.5,0) coordinate (D)--(0,22.5) coordinate (A) 
					--(22.5,0) coordinate (B)--(0,-22.5) coordinate  (C)--cycle
					(7.5,15) coordinate  (M) node[shift={(40:8pt)}]{$M$}
					(15,7.5) coordinate  (N) node[shift={(50:8pt)}]{$N$};
				\foreach  \i in {0,...,3}{
					 \draw[fill=white,rotate=90*\i] plot[domain=-7.5:7.5] (\x,{f (\x)}) node[dot]{}
					--(0,22.5) node[dot]{}--(-7.5,15) node[dot]{}; 
					\draw[->] (90+90*\i:22.5)--node[sloped,above]{$15$~m} 
					 (90+90*\i:7.5);
				} 
				 \draw[red,rotate=-90] plot[domain=-7.5:7.5] (\x,{f (\x)});
				\draw  (-22.5,0)--(0,22.5)--(22.5,0)--(0,-22.5)--cycle;
				\foreach \i/\j in {A/90,B/0,C/-90,D/180}\fill  (\i) circle (1pt) node[shift={(\j:.25)}]{$\i$};
			\end{tikzpicture}
	\end{center}
	 \shortans{349}
	\loigiai{ 
		\begin{center}
			\begin{tikzpicture}[>=stealth,line join=round,scale=.2]
					 \tikzset{declare function={f (\x)=2/15*(\x)^2+15/2;},smooth,samples=50,dot/.style={circle,inner sep=1pt}}
					 \fill[gray!20] (-22.5,0) coordinate (D)--(0,22.5) coordinate (A)
						--(22.5,0) coordinate (B)--(0,-22.5) coordinate  (C)--cycle
						(7.5,15) coordinate (M) node[shift={(40:8pt)}]{$M$} 
						(15,7.5) coordinate  (N) node[shift={(50:8pt)}]{$N$};
					\foreach \i in {0,...,3}{ 
						\draw[fill=white,rotate=90*\i] plot[domain=-7.5:7.5] (\x,{f (\x)})  node[dot]{}
						--(0,22.5) node[dot]{}--(-7.5,15)  node[dot]{};
					 }
					 \draw[red,rotate=-90] plot[domain=-7.5:7.5] (\x,{f (\x)});
					 \draw (-22.5,0)--(0,22.5)--(22.5,0)--(0,-22.5)--cycle;
					\draw[dashed] (15/2,0) node[dot]{} node[above left]{$K$}  node[shift={(-120:12pt)},scale=.8]{$\dfrac{15}{2}$}
						--(M)--(0,15) node[dot]{}  node[above right]{$H$} node[left]{$15$};
					\draw pic["$45^\circ$",draw,angle radius = 15pt,angle  eccentricity=1.6]{angle = C--A--B};
					\draw[->] (0,0)  node[below left=-1pt]{$O$}
						 (-25,0)--(30,0) node[pos=.98,above]{$x$};
					\draw[->](0,-25)--(0,30) node[pos=.98,left]{$y$}; 
					\foreach \i/\j in {A/30,B/50,C/-35,D/-120}\fill (\i) circle (1pt) node[shift={(\j:.25)}]{$\i$}; 
				\end{tikzpicture}
		\end{center}
		Gán hệ tọa  độ $Oxy$ như hình vẽ: $A\left(0;\dfrac{45}{2}\right)$, $B\left(\dfrac{45}{2};0\right)$, $C\left(0;-\dfrac{45}{2}\right)$, $D\left(-\dfrac{45}{2};0\right)$.\\
		Do hình vuông $ABCD$ có $AC=45$\,m nên độ dài cạnh hình vuông là $AB  = \dfrac{AC}{\sqrt{2}} = \dfrac{45}{\sqrt{2}}$\,m.\\
		Diện tích hình vuông $ABCD$ là  $S_{ABCD} = AB^2 = \dfrac{45^2}{2} = 1012{,}5$\,m$^2$.\\
		Vì $AM=MN=NB$ nên $AM=MN=NB  = \dfrac{AB}{3} = \dfrac{15\sqrt{2}}{2}$\,m.\\
		Độ dài cạnh 
		$AH = AM \cdot \cos 45^\circ = \dfrac{15}{2}$; 
		$OH = OA - AH  = \dfrac{45}{2} - \dfrac{15}{2} = 15$;
		$HM = AM \cdot \sin 45^\circ = \dfrac{15}{2}$.\\ 
		Vậy tọa độ  của điểm $M$ là $\left(\dfrac{15}{2};15\right)$.\\
		Parabol  đi qua điểm $M$ và có đỉnh là điểm $E\left(0;\dfrac{15}{2}\right)$ có phương trình dạng $(P)\colon y = ax^2 + \dfrac{15}{2} \hfill{(1)}$.\\
		 Thay tọa độ điểm $M\left(\dfrac{15}{2};15\right)$ vào ($1$) ta được
		\[15 = a \cdot \left(\dfrac{15}{2}\right)^2 + \dfrac{15}{2} \Leftrightarrow a = \dfrac{2}{15}. \] 
		Vậy phương trình parabol $(P)$ là $y = \dfrac{2}{15}x^2  + \dfrac{15}{2}$.\\
		Diện tích hình phẳng giới hạn bởi $(P)$, đường thẳng $y=15$, trục $Oy$ và đường thẳng $x=\dfrac{15}{2}$ là 
		\begin{eqnarray*}
			S_1 &= &\displaystyle\int_{0}^{\tfrac{15}{2}}  \left[15 - \left(\dfrac{2}{15}x^2 + \dfrac{15}{2}\right)\right] \mathrm{d}x \\
			&=&\displaystyle\int_{0}^{\tfrac{15}{2}} \left(\dfrac{15}{2} -  \dfrac{2x^2}{15}\right) \mathrm{d}x \\
			& = & \left(\dfrac{15}{2}x  - \dfrac{2x^3}{45}\right)\Big|_{0}^{\tfrac{15}{2}}\\
			&=&  37{,}5\,\text{ m}^2.
		\end{eqnarray*}
		Diện tích tam giác $\triangle AHM$ là $S_2 = S_{AHM} = \dfrac{1}{2}AH \cdot HM = \dfrac{1}{2}  \cdot \dfrac{15}{2} \cdot \dfrac{15}{2} = 28{,}125$\,m$^2$.\\
		Diện tích trồng hoa là $S_{\text{ hoa}} = 8\cdot (S_1+S_2) = 8\cdot (37{,}5+28{,}125)  = 525$\,m$^2$.\\
		Diện tích sân chơi là  $S_{\text{ sân}} = S_{ABCD} - S_{\text{ hoa}} = 1012{,}5 - 525 = 487{,}5$\,m$^2$.\\
		Tổng số tiền cần dùng là $487{,}5 \cdot 500\,000 + 525 \cdot 200\,000 = 348\,750\,000$ đồng.\\ 
		Vậy số tiền cần dùng để trang trí  sân và trồng hoa là $349$ triệu đồng.
	 }
\end{bt}

% nguồn: Phan2_Lop_12_Dai_so_Chuong_4
\begin{bt}%[2D4V3-3]%[ma:aeeb220e]
	Gọi $V$ là thể tích khối tròn xoay tạo thành khi quay hình phẳng giới hạn bởi các đường $y=\sqrt x$, $y=0$ và $x=4$ quanh trục $Ox$. Đường thẳng $x=a$ $\left(0 < a < 4\right)$ cắt đồ thị hàm số $y=\sqrt x  $ tại $M$ (xem hình vẽ sau). Gọi $V_1$ là thể tích khối tròn xoay tạo thành khi quay tam giác $OMH$ quanh trục $Ox$. Biết rằng $V=2V_1$. Khi đó giá trị của $a$ là bao nhiêu?
	\begin{center}
		\begin{tikzpicture}[scale=1.2,font=\footnotesize,>=stealth]
				\def\xmin{-0.5} \def\xmax{4.5} 
				\def\ymin{-2.5}  \def\ymax{2.5}
				\def\a{2} 
				\fill[gray!50,domain=0:4,samples=100] 
					(0,0) 
					--  plot (\x,{sqrt (\x)})
					 -- (4,0)
					 -- cycle;
				\draw[->] (\xmin,0) -- (\xmax,0) node[below]  {$x$};
				\draw[->] (0,\ymin)  -- (0,\ymax) node[left] {$y$};
				\draw[domain=0:4,samples=100]  plot (\x,{sqrt (\x)});
				 \draw[dashed,domain=0:4,samples=100] plot (\x,{-sqrt (\x)});
				 \draw[dashed] (\a,0) -- (\a,{sqrt (\a)}) coordinate (M);
				\node[above right] at (M)  {$M$};
				\coordinate (K) at  (\a,0);
				\node[below] at (K)  {$K$};
				\coordinate  (H) at (4,0);
				\node[below] at  (H) {$H$};
				\coordinate (O) at (0,0); 
				\node[below left]  at (O) {$O$};
				\draw (O) -- (M) -- (H); 
				\draw[dashed]  (\a,0) ellipse (0.15 and {sqrt (\a)});
				\draw[dashed] (4,0) ellipse (0.25 and  2);
				\draw[dashed]  (M) -- (\a,{-sqrt (\a)});
				 \draw[dashed] (H) -- (4,-2);
				\node[right] at (\a,0.7)  {$a$};
				\node[above  right] at (4,0) {$4$};
			\end{tikzpicture}
	\end{center}
	 \shortans{3}
	 \loigiai{
		Ta  có $V=\pi\displaystyle\int_0^4x\,\mathrm{d}x=8\pi $ và tọa độ $M\left(a;\sqrt a\right)$.\\
		Khi quay tam giác $OMH$ quanh trục  $Ox$ tạo thành hai hình nón có chung đáy
		\begin{itemize}
			\item Hình nón $\left(N_1\right)$ có đỉnh là $O$, chiều cao $h_1=OK=a$, bán kính đáy $R=MK=\sqrt  a $;
			\item Hình nón $\left(N_2\right)$ có đỉnh là $H$, chiều cao  $h_2=HK=4-a$, bán kính đáy $R=MK=\sqrt a$.
		\end{itemize}
		Do đó $V_1=\dfrac{1}{3}\pi{R^2}h{\kern 1pt_1}+\dfrac{1}{3}\pi{R^2}h{\kern 1pt_2}=\dfrac{1}{3}\pi{(\sqrt{a})^2}h{\kern 1pt_1}+\dfrac{1}{3}\pi{(\sqrt{a})^2}h{\kern 1pt_2}=\dfrac{4}{3}\pi a$.\\ 
		Khi đó  $V=2V_1\Leftrightarrow 8\pi=2\cdot\dfrac{4}{3}\pi a\Rightarrow a=3$.
	 }
\end{bt}

% nguồn: Phan2_Lop_12_Dai_so_Chuong_4
\begin{bt}%[2D4V3-3]%[ma:582e7c63]
	Một thùng làm kem có mặt cắt dạng parabol như hình. Biết phương trình đường biên có  dạng $f(x)=a\sqrt{x}$, hỏi dung tích của thùng bằng bao nhiêu lít? (kết quả làm tròn đến hàng phần mười).
	\begin{center}
		\begin{tikzpicture}[scale=0.8,line join=round,line cap=round,>=stealth]
				\tikzset{declare  function={f (\x)=1/8*(\x)^2;
					xI=0; yI=2; xmin=xI-6;  xmax=xI+5;
					 ymin=yI-0.5; ymax=yI+1;},
					smooth,samples=100 
				} 
				\fill[pattern=dots] (-4,2)--(-5,2)--(-5,-1)--(5,-1)--(5,2)--(4,2)--plot[domain=4:-4] (\x,{f (\x)}); 
				\draw[<-] (xmin,0) node[below]  {$y$}--(xmax,0);
				\draw[->] (0,-1.5)--(0,ymax) node[below right] {$x$}; 
				 \fill (0,0) node[above left]{$O$} circle (1pt);
				\draw[cyan,domain=-4:4] plot  (\x,{f (\x)});
				\draw[|<->|](0,2)--(4,2) node[pos=0.5,above]{$30$\,cm}; 
				\draw[|<->|](5.2,-1)--(5.2,2) node[pos=0.5,right]{$20$\,cm}; 
				 \draw (-4,2)--(-5,2)--(-5,-1)--(5,-1)--(5,2)--(4,2);
			\end{tikzpicture}
	\end{center}
	 \shortans{28,3}
	\loigiai{ 
		Từ hình vẽ ta có $f(20) = 30 \Rightarrow a\sqrt{20} = 30 \Rightarrow a =  \dfrac{15}{\sqrt{5}} = 3\sqrt{5}$.\\
		Vậy $f(x) =  3\sqrt{5}\sqrt{x}$.\\
		Thể  tích khối tròn xoay là\\
		$V = \pi \int_{0}^{20} f^2(x)\mathrm{d}x = \pi \int_{0}^{20} \left(3\sqrt{5}\right)^2 x\mathrm{d}x = \pi  \int_{0}^{20} 45x\mathrm{d}x
		= \pi  \cdot \left[\dfrac{45}{2}x^2\right]_{0}^{20} = \pi \cdot 9\,000 \approx 28\,274{,}33 \text{ cm}^3 \approx 28{,}3 \text{ lít}$.
	} 
\end{bt}

% nguồn: Phan2_Lop_12_Dai_so_Chuong_4
\begin{bt}%[2D4V3-3]%[ma:d4936d04]
	Cho hình phẳng giới hạn bởi các đường $y=x\ln x$, $y=0$, $x=\mathrm{e}$ quay xung quanh trục $Ox$ tạo thành khối tròn xoay có thể tích bằng $\dfrac{\pi}{a}\left(b{\mathrm{e}^3}-2\right)$ với. Tính $P=a^2-b^2$, biết một nguyên  hàm của hàm số $f(x)=\left(x\ln x\right)^2$ là $F(x)=\dfrac{x^3}{27}\left(9\ln^2x-6\ln x+2\right)$.
	 \shortans{704}
	\loigiai{ 
		Phương trình  hoành độ giao điểm
		 \[x\ln x=0\Leftrightarrow x=\mathrm{1}.\]
		 Thể tích của khối tròn xoay là \[V=\pi\displaystyle\int_1^\mathrm{e}{y^2\,\mathrm{d}x}=\pi\displaystyle\int_1^\mathrm{e}{\left(x\ln x\right)^2\,\mathrm{d}x}=\pi \cdot F(x)\bigg|_1^\mathrm{e}=\dfrac{\pi}{27}\left(5\mathrm{e}^3-2\right).\]
		 Do đó $a=27$, $b=5$.\\
		Vậy  $P=a^2-b^2=704$.
	} 
\end{bt}

% nguồn: Phan2_Lop_12_Dai_so_Chuong_4
\begin{bt}%[2D4V3-3]%[ma:df956fff]
	Một khối cầu có bán kính là $5$ (dm), người ta cắt bỏ hai phần của khối cầu bằng hai mặt phẳng song song và cách tâm một khoảng $3$ (dm) để làm một chiếc lu đựng nước (như hình vẽ). Tính thể tích  (đơn vị dm$^3$) mà chiếc lu chứa được (kết quả làm tròn đến hàng đơn vị).
	 \shortans{415}
	\begin{center}
		\begin{tikzpicture}[scale=.7,>=stealth,line join=round,line cap=round,font=\footnotesize]
				 \def\r{5}
				\pgfmathsetmacro\g{atan  (3/4)}
				\pgfmathsetmacro\a{\r*cos (\g)} 
				 \pgfmathsetmacro\y{\r*sin (\g)}
				 \pgfmathsetmacro\b{\a/6}
				 \path (0,0) coordinate (O)
					 +(\g:\r) coordinate (A)
					+({180-\g}:\r) coordinate (A') 
					 +({-\g}:\r) coordinate (B);
				\draw (B) arc ({-\g}:{\g}:{\r}) 
					 (A') arc ({180-\g}:{180+\g}:{\r})
					(A) arc (0:-180:{\a} and {\b}) 
					(B) arc (0:-180:{\a} and  {\b});;
				\draw [dashed] (A) arc ({\g}:{180-\g}:{\r}) 
					(B) arc  ({-\g}:{-180+\g}:{\r});
					(A) arc (0:-180:{\a}  and {\b})
					(B) arc (0:-180:{\a}  and {\b});
				\draw [dash  pattern=on 1pt off 1pt]
					(A) arc (0:180:{\a} and  {\b})
					(B)  arc (0:180:{\a} and {\b})
					(O)--++(0:\r) node [above,midway,fill=white] {$5$  dm}
					(O)--(0,{-\y}) node  [midway,sloped,above,fill=white] {$3$ dm}
					 (0,{-\r})--(0,\r)
					(A)--(0,{\y}) 
					 (B)--(0,{-\y});
				\foreach \i/\j in  {0/0,\r/0,0/\y,0/{-\y},0/\r,\a/\y,\a/{-\y},0/{-\r}}{
					\fill (\i,\j)  circle (2pt);
				} 
			\end{tikzpicture}
	\end{center}
	\loigiai{ 
		\begin{center}
			\begin{tikzpicture}[scale=.7,>=stealth,line join=round,line cap=round,font=\footnotesize]
					 \def\r{5}
					\pgfmathsetmacro\g{atan (3/4)} 
					 \pgfmathsetmacro\a{\r*cos (\g)}
					 \pgfmathsetmacro\y{\r*sin (\g)}
					\pgfmathsetmacro\b{\a/6} 
					\path (0,0)  coordinate (O)
						 +(\g:\r) coordinate (A)
						 +({180-\g}:\r) coordinate (A')
						 +({-\g}:\r) coordinate (B);
					 \fill [pattern=north east lines] (B) arc ({-\g}:{\g}:{\r})--(0,{\y})--(0,{-\y})--cycle;
					\draw  (B) arc ({-\g}:{\g}:{\r})
						(A') arc ({180-\g}:{180+\g}:{\r}) 
						(A) arc (0:-180:{\a} and  {\b})
						(B)  arc (0:-180:{\a} and {\b});;
					\draw [dashed] (A) arc  ({\g}:{180-\g}:{\r})
						(B) arc  ({-\g}:{-180+\g}:{\r});
						(A) arc (0:-180:{\a} and  {\b})
						(B) arc (0:-180:{\a} and {\b}); 
					\draw [dash pattern=on  1pt off 1pt]
						(A) arc (0:180:{\a} and {\b}) 
						 (B) arc (0:180:{\a} and {\b})
						 (O)--++(0:\r) node [above,midway,fill=white,inner sep=.5pt] {$5$ dm}
						 (O)--(0,{-\y}) node [midway,sloped,above,fill=white,inner sep=.5pt] {$3$ dm}
						 (0,{-\r})--(0,\r)
						 (A)--(0,{\y})
						(B)--(0,{-\y}); 
					\foreach \i/\j in  {0/0,\r/0,\a/\y,\a/{-\y}}{
						\fill (\i,\j) circle (2pt); 
					} 
					\draw [->,dash pattern=on  1pt off 1pt] (O)--++(-135:{\r/2}) node [left] {$z$};
					\draw [->]  ({\r},0)--++(0:1) node [below,xshift=-3pt] {$y$};
					\draw  [->] (0,{\r})--++(90:1) node [left,yshift=-4pt] {$x$};
					 \draw (0,{-\r})--++(-90:1);
					 \fill (O) circle (2pt) node [left] {$O$}
						(0,{\y}) circle (2pt) node [left]  {$3$}
						(0,{-\y}) circle (2pt) node [left]  {$-3$};
					\draw [->]  (O)++({\g/2}:{\r})--++(0:1) node [right] {$y=\sqrt{25-x^2}$};
				\end{tikzpicture}
		\end{center}
		 Đường tròn có phương trình $x^2 + y^2 = 25 $ nên $y = \pm\sqrt{25 - x^2}$.\\
		Thể tích phần hình tròn quay quanh trục $Ox$ từ $x = -3$ đến $x = 3$ là 
		\[ 
		V = \pi  \displaystyle\displaystyle\int\limits_{-3}^{3} (25 - x^2)\mathrm{\,d}x
		 = \pi \left[25x - \dfrac{x^3}{3} \right]_{-3}^{3}
		=  132\pi\approx 415.
		 \]
		Vậy thể tích cần tìm là  $415$ dm$^3$.
	 }
\end{bt}

% nguồn: Phan2_Lop_12_Dai_so_Chuong_4
\begin{bt}%[NBV2526]%[2D4V3-3]%[ma:a8451f58]
	Cho  hình vuông $OABC$ cạnh bằng $8$, điểm $M$ nằm trong hình vuông sao cho khoảng cách từ $M$ đến các cạnh $OA, OC$ cùng bằng $3$. Parabol $(P_{1})$ đi qua các điểm $O$, $A$, $M$. Parabol $(P_{2})$ đi qua các điểm $O$, $C$, $M$. Tính diện tích phần tô đậm (hình vẽ bên). (kết quả làm tròn kết quả đến hàng đơn vị).
	\begin{center}
		\begin{tikzpicture}[line join=round,font=\footnotesize,scale=.7]
				\tikzset{declare function={f  (\x)=-1/5*(\x)^2+8/5*\x;},smooth,samples=150}
				\path (0,0)  coordinate (O)
					(8,0) coordinate (A) 
					(8,8)  coordinate (B)
					(0,8) coordinate  (C)
					(3,3) coordinate (M); 
				\fill[cyan!20] (O) rectangle (B); 
				\fill[white] plot[domain=0:8]  (\x,{f (\x)});
				 \fill[white] plot[domain=0:8] ({f (\x)},\x);
				\draw plot[domain=0:8] (\x,{f (\x)}); 
				\draw plot[domain=0:8] ({f (\x)},\x); 
				\draw (O) rectangle (B); 
				 \foreach \i/\j in {A/-45,B/45,C/135,O/-135,M/-50}\fill (\i) circle (1pt) node[shift={(\j:.2)}]{$\i$};
			\end{tikzpicture}
	\end{center}
	\shortans{32} 
	 \loigiai{
		Cho hệ  trục tọa độ $Oxy$ sao cho tia $Ox$ trùng tia $OA$, tia $Oy$ trùng tia $OC$.\\
		Gọi $S_1$ là diện tích hình phẳng giới hạn bởi các đường $(P_1)$, $y=0$.\\ 
		$S_2$ là diện tích hình phẳng giới hạn bởi các đường $(P_2)$, $x=0$.\\ 
		$S_3$ là diện tích hình phẳng giới hạn bởi các đường $(P_1)$  và $(P_2)$.\\
		Từ hình vẽ ta thấy diện tích hình phẳng cần tìm là $$S =  S_{OABC} - [S_1 + S_2 - S_3] \hfill{(*)}$$
		Gọi  phương trình $(P_1)$ là $y = ax^2 + bx + c$.\\
		Parabol $(P_{1})$ đi  qua $O(0;0)$, $A(8;0)$, $M(3;3)$ nên $\heva{&c=0\\&64a+8b+c=0\\&9a+3b+c=3} \Leftrightarrow \heva{&a=-\dfrac{1}{5}\\
		&b=\dfrac{8}{5}\\&c=0.}$\\ 
		Vậy $(P_{1})\colon  y=-\dfrac{1}{5}x^{2}+\dfrac{8}{5}x$.\\
		Ta thấy $(P_1)$ và $(P_2)$ đối  xứng qua $OM$ nên
		\[S_1 = S_2 = \displaystyle\int_{0}^{8}  \left(-\dfrac{1}{5}x^2+\dfrac{8}{5}x \right)\mathrm{d}x = \dfrac{256}{15}. \]
		Gọi $S_4$ là diện tích hình phẳng  giới hạn bởi các đường $(P_1)$, $OM\colon y=x$, do tính chất đối xứng nên
		\[S_3 = 2S_4 = 2\int_{0}^{3}\left[\left(-\dfrac{1}{5}x^2+\dfrac{8}{5}x\right)-x\right]\mathrm{d}x  = \dfrac{9}{5}. \]
	 Thay vào $(*)$ ta được $S = 64 - \left(\dfrac{256}{15} + \dfrac{256}{15} - \dfrac{9}{5}\right) = \dfrac{95}{3} \approx 32$. }
\end{bt}

% nguồn: Phan2_Lop_12_Dai_so_Chuong_4
\begin{bt}%[2D4V3-3]%[ma:d23c3e5b]
	Cho các nửa đường tròn như hình  vẽ sau, trong đó đường kính của nửa đường tròn lớn gấp đôi đường kính của nửa đường tròn nhỏ. Biết rằng nửa hình tròn đường kính $AB$ có diện tích là $32\pi $ và góc $\widehat{BAC}=30^\circ $. Tính thể tích vật thể tròn xoay được tạo thành khi quay hình phẳng $(H)$ (phần tô đậm) xung quanh đường thẳng $AB$ (làm tròn kết quả đến hàng đơn vị).
	\begin{center}
		\begin{tikzpicture}[scale=.5,font=\footnotesize,line join=round,line cap=round,>=stealth]
				 \def\xmin{-2.5}\def\xmax{5}\def\ymin{-2.5}\def\ymax{2.5}
				\def\f (#1){(sqrt (16*(#1)-(#1)^2)} 
				\def\g (#1){(sqrt  (8*(#1)-(#1)^2}
				\fill[color=gray](6,3.46)--(12,6.93)--plot[domain=12:16](\x,{\f  (\x)})--(8,0)--plot[domain=8:6](\x,{\g (\x)})--(6,3.46)--cycle;
				 \draw (0,0) arc (180:0:4);
				\draw  (0,0) arc (180:0:8);
				 \draw (12,6.92) node[above]{$C$}--(0,0)--(16,0) node[right]{$B$};
				\coordinate (A) at  (0,0) node[left]{$A$};
				\coordinate (B) at  (16,0);
				\coordinate  (C) at (12,6.92);
				\draw (12,3)  node[above]{$(H)$};
				 \draw pic["$30^\circ$",draw,angle eccentricity=1.7,angle radius=0.5cm]{angle=B--A--C};
			\end{tikzpicture}
	\end{center}
	 \shortans{821}
	 \loigiai{
		Chọn hệ trục  $Oxy$ thích hợp như hình sau\\
		\begin{center}
			\begin{tikzpicture}[scale=.5,font=\footnotesize,line join=round,line cap=round,>=stealth]
					\def\xmin{-2.5}\def\xmax{5}\def\ymin{-2.5}\def\ymax{2.5} 
					\def\f (#1){(sqrt  (16*(#1)-(#1)^2)}
					 \def\g (#1){(sqrt (8*(#1)-(#1)^2}
					\draw[->](-1,0)--(17.5,0)  node[above]{$x$};
					\draw[->](0,-1)--(0,7.5)  node[left]{$y$};
					\fill[color=gray](6,3.46)--(12,6.93)--plot[domain=12:16](\x,{\f (\x)})--(8,0)--plot[domain=8:6](\x,{\g  (\x)})--(6,3.46)--cycle;
					\draw (0,0) arc  (180:0:4);
					\draw (0,0) arc (180:0:8); 
					\draw (12,6.92) node[above]{$C$}--(0,0)--(16,0) node[below]{$B$}; 
					\draw[dashed](6,0)  node[below]{$6$}--(6,3.46);
					\draw[dashed](12,0) node[below]{$12$}--(12,6.92); 
					 \coordinate (A) at (0,0) node[above left]{$A$};
					\coordinate (O) at  (0,0) node[below left]{$O$};
					\coordinate (B)  at (16,0);
					\coordinate (C) at  (10,6.92);
					\draw (12,3)  node[above]{$(H)$};
					\draw  (8,0) node[below]{$8$};
					\draw pic["$30^\circ$",draw,angle eccentricity=1.7,angle  radius=0.5cm]{angle=B--A--C};
				\end{tikzpicture}
		\end{center}
		Phương trình nửa  đường tròn lớn $f(x)=\sqrt{16x-x^2}$.\\
		 Phương trình nửa đường tròn nhỏ $g(x)=\sqrt{8x-x^2}$.\\
		Phương trình đường thẳng  $AC\colon y=\dfrac{\sqrt 3}{3}x$.\\
		Khi đó 
		 \allowdisplaybreaks
		\begin{eqnarray*}
			 V&=&\pi\displaystyle\int_6^8\left|\left(\dfrac{\sqrt 3}{3}x\right)^2-\left(8x-x^2\right)\right|\,\mathrm{d}x+\pi\displaystyle\int_8^{12}{\left(\dfrac{\sqrt 3}{3}x\right)^2\,\mathrm{d}x}+\pi\displaystyle\int_{12}^{16}{\left(16x-x^2\right)\,\mathrm{d}x}\\
			&=&\pi\displaystyle\int_6^8\left(\dfrac{4}{3}x^2-8x\right)\,\mathrm{d}x+\pi\displaystyle\int_8^{12}{\dfrac{1}{3}x^2\,\mathrm{d}x}+\pi\displaystyle\int_{12}^{16}{\left(16x-x^2\right)\,\mathrm{d}x}\\ 
			 &=&\dfrac{784\pi}{3}\approx 821.
		\end{eqnarray*}
	} 
\end{bt}

% nguồn: Phan2_Lop_12_Dai_so_Chuong_4
\begin{bt}%[NBV2526]%[2D4V3-3]%[ma:1c133784]
	Một công ty đang thiết kế một bảng quảng cáo hình  chữ nhật $ABCD$ có kích thước $AB=12$\,m và $AD=8$\,m. Phần trung tâm của bảng sẽ được in nội dung quảng cáo, được mô tả là phần tô đậm (xem hình minh họa). Hai đường cong trong hình là một phần của đồ thị hàm số có dạng $y=\dfrac{ax+b}{cx+d}$. Đường tiệm cận đứng và tiệm cận ngang của đồ thị hàm số này đều cách điểm $A$ một khoảng bằng $5$\,m. Đồ thị giao với cạnh $AB$ tại điểm $E$ thỏa mãn $\dfrac{AE}{AB}=\dfrac{7}{15}$. Diện tích phần in nội dung quảng cáo là bao nhiêu mét vuông (làm tròn kết quả đến hàng phần mười)?
	\begin{center}
		\begin{tikzpicture}[>=stealth,line join=round,font=\footnotesize,scale=.4]
				 \tikzset{declare function={f (\x)=-3/\x;},smooth,samples=150}
				\path (-5,-5) coordinate  (A)
					(7,-5) coordinate  (B)
					 (7,3) coordinate (C)
					 (-5,3) coordinate (D)
					(.6,{f  (.6)}) coordinate (E)
					(7,{f (7)}) coordinate (F); 
				\fill[gray!30] plot[domain=-5:-1] (\x,{f  (\x)})--(C)--(F)--plot[domain=7:.6] (\x,{f (\x)})--(A)--cycle;
				\draw (A) rectangle (C) plot[domain=-5:-1] (\x,{f (\x)}) plot[domain=7:.6] (\x,{f (\x)}); 
				 \draw[dashed] (-5,0)--(7,0)
					 (0,-5)--(0,3);
				\foreach \i/\j in {A/-135,B/-45,C/45,D/135,E/-90,F/0}\fill (\i) circle (1.5pt) node[shift={(\j:.2)}]{$\i$}; 
			\end{tikzpicture}
	\end{center}
	\shortans{73,5} 
	 \loigiai{
		\begin{center}
			\begin{tikzpicture}[>=stealth,line join=round,font=\footnotesize,scale=.4]
					\tikzset{declare function={f (\x)=-3/\x;},smooth,samples=150} 
					\path (-5,-5) coordinate (A) 
						(7,-5)  coordinate (B)
						(7,3)  coordinate (C)
						(-5,3) coordinate  (D)
						 (.6,{f (.6)}) coordinate (E)
						(7,{f  (7)}) coordinate (F);
					\fill[gray!30] plot[domain=-5:-1] (\x,{f  (\x)})--(C)--(F)--plot[domain=7:.6] (\x,{f (\x)})--(A)--cycle;
					\draw (A) rectangle (C) plot[domain=-5:-1]  (\x,{f (\x)}) plot[domain=7:.6] (\x,{f (\x)});
					\draw[dashed] (-5,0)--(7,0) 
						 (0,-5)--(0,3);
					\foreach \i/\j  in {A/-135,B/-45,C/45,D/135,E/-90,F/0}\fill (\i) circle (1.5pt) node[shift={(\j:.2)}]{$\i$};
				\end{tikzpicture}
		\end{center}
		Gắn hệ trục tọa độ $Oxy$ như hình vẽ $A(-5;-5)$, $B(7;-5)$, $C(7;3)$,  $D(-5;3)$. Trục hoành $Ox$ trùng với cạnh $AD$ sau khi tịnh tiến lên $5$ đơn vị. Trục tung $Oy$ trùng với cạnh $AB$ sau khi tịnh tiến sang phải $5$ đơn vị.\\
		Các  đường tiệm cận cách điểm $A(-5;-5)$ một khoảng bằng $5$\,m là
		\begin{itemize}
			\item Tiệm cận đứng $x = -5 + 5 = 0$. 
			\item Tiệm cận ngang $y  = -5 + 5 = 0$.
		\end{itemize}
		Đồ  thị hàm số $y=\dfrac{ax+b}{cx+d}$ có tiệm cận đứng $x=-\dfrac{d}{c}=0$ và tiệm cận ngang $y=\dfrac{a}{c}=0$.
		Từ  đó suy ra $d=0$ và $a=0$.\\ Phương trình hàm số có dạng $y=\dfrac{b}{cx}$. \\
		Ta có  $\dfrac{AE}{AB}=\dfrac{7}{15} \Leftrightarrow AE=\dfrac{7}{15} \cdot AB = \dfrac{7}{15} \cdot 12 = 5{,}6$\,m. Vậy tọa độ điểm $E$ là $E(0{,}6;-5)$.\\
		Đồ thị hàm số đi qua điểm  $E(0{,}6;-5)$ nên
		$$-5 = \dfrac{b}{c\cdot 0{,}6} \Leftrightarrow \dfrac{b}{c} =-3 \text{ suy ra  } y=-\dfrac{3}{x}.$$
		Ta có 
		\begin{eqnarray*}
			S& = & \displaystyle\int_{-5}^{-1}\left(5+\dfrac{3}{x}\right)\mathrm{\,d}x\\ 
			& = &\left[5x  + 3\ln|x|\right]\Big|_{-5}^{-1}\\
			& = & 20  - 3\ln 5\,\left(\text{m}^2\right).
		\end{eqnarray*}
		Diện tích $S'$ (phần trắng dưới bên phải) giới hạn bởi đồ thị $y=-\dfrac{3}{x}$, đường thẳng $y=-5$ (cạnh $AB$) và đường thẳng $x=7$ (cạnh $BC$). Đồ thị cắt $y=-5$ tại $x=0{,}6$ (điểm $E$). Đồ thị cắt $x=7$  tại $y=-\dfrac{3}{7}$.
		\begin{eqnarray*}
			S'&=&\displaystyle\int_{0{,}6}^{7}\left(\dfrac{3}{x}\right)\mathrm{\,d}x  = 3\ln\left(\dfrac{35}{3}\right)\,\left(\text{m}^2\right).
		\end{eqnarray*}
		Diện tích phần in nội dung quảng cáo $S_{ABCD} - S - S' = 96 - (20 - 3\ln5) - 3\ln\left(\dfrac{35}{3}\right)\approx 73{,}5$\,(m$^2$). 
	 }
\end{bt}

% nguồn: Phan2_Lop_12_Dai_so_Chuong_4
\begin{bt}%[2D4V3-3]%[ma:9774f642]
	Một đèn pha dạng parabol có đường biên (mặt cắt) cho bởi $f(x)=\dfrac{3}{2}\sqrt{x}$. Chiều dài chóa là $4$\,cm. Hỏi thể tích không khí bên trong là bao nhiêu cm$^3$? (làm tròn  kết quả
	đến hàng phần  mười)
	\begin{center}
		\begin{tikzpicture}[scale=1,line join=round,line cap=round,>=stealth]
				\tikzset{declare function={f (\x)=1.5*sqrt (\x); 
					xI=0; yI=3;  xmin=xI; xmax=xI+4.5;
					 ymin=-3.5; ymax=yI+1;},
					 smooth,samples=100
				 }
				 \def\bongden{(-0,-0.21).. controls +(178:0.68) and +(-88:0.82)..
					(-1.24,1.1).. controls  +(92:0.32) and +(-160:0.5)..
					(-0.56,2.06).. controls +(53:0.07) and  +(-81:0.07)..
					(-0.48,2.23)..  controls +(96:0.07) and +(-90:0.04)..
					 (-0.56,2.38).. controls +(13:0.36) and +(168:0.36)..
					(0.52,2.38).. controls +(-90:0.04)  and +(84:0.07)..
					(0.44,2.23).. controls +(-96:0.07) and +(127:0.07).. 
					 (0.52,2.06).. controls +(-20:0.5) and +(89:0.32)..
					(1.2,1.1).. controls  +(-92:0.82) and +(2:0.68)..
					 (-0.04,-0.21);}
				\draw[->] (xmin-0.5,0)--(0,0) node[below left]{$O$} --(xmax+1,0) node[below  left] {$x$};
				\draw[->] (0,ymin)--(0,ymax) node[below  right] {$y$};
				\fill  (0,0) circle (1pt);
				\draw[yellow,line  width=2pt] (0.5,0.3)--(1.1,{f (1)})
					 (0.5,-0.3)--(1.1,{-f (1)})
					(0.4,0.3)--(0.5,{f (0.4)}) 
					(0.4,-0.3)--(0.5,{-f (0.4)}); 

				\draw[yellow,line width=2pt,->] (1.1,{f (1)})  -- (5.5,{f (1)});
				\draw[yellow,line width=2pt,->] (0.5,{f (0.4)}) --  (5.5,{f (0.4)});
				\draw[yellow,line width=2pt,->] (0.5,{-f  (0.4)}) -- (5.5,{-f (0.4)});
				\draw[yellow,line  width=2pt,->] (1.1,{-f (1)}) -- (5.5,{-f (1)});
				\draw[red,domain=0:4] plot (\x,{f (\x)}) plot (\x,{-f  (\x)})
					(4,0) ellipse (0.8 and 3); 
				 \fill
					(4,0)  circle (1pt) node[shift={(-90:7pt)}]{$4$};
				\path[inner color=orange,outer color=yellow,scale=0.25,rotate around  z=90,xshift=2.5cm]\bongden
			\end{tikzpicture}
	\end{center}
	 \shortans{56,6}
	 \loigiai{
		Thể tích khối tròn xoay quanh  trục $Ox$ là\\
		$V = \pi  \int_{0}^{4} f^2(x)\mathrm{d}x = \pi \int_{0}^{4} \left(\dfrac{3}{2}\sqrt{x}\right)^2 \mathrm{d}x = \pi \int_{0}^{4} \dfrac{9}{4}x\mathrm{d}x
		= \pi \left[\dfrac{9}{8}x^2\right]\Bigg|_{0}^{4}  = \pi \cdot \dfrac{9}{8} \cdot 16 = \pi \cdot 18 \approx 56{,}6$\,cm$^3$.
	} 
\end{bt}

% nguồn: Phan2_Lop_12_Dai_so_Chuong_4
\begin{bt}%[2D4V3-3]%[ma:1878863e]
	Tính thể tích của khối tròn xoay sinh ra khi quay quanh trục $Ox$ hình phẳng giới hạn bởi đồ thị hàm số $(C)\colon y=x^3-3x$ và đường thẳng $d:y=x$ (\textit{kết quả làm tròn  đến hàng phần mười}).
	\shortans{29,6} 
	\loigiai{ 
		\begin{center}
			\begin{tikzpicture}[line join=round,line cap=round,>=stealth,font=\footnotesize,scale=1]
					\draw[->]  (-2.6,0)--(2.6,0) node[below left] {$x$};
					\draw[->] (0,-3.1)--(0,3.1) node[below  left] {$y$};
					\begin{scope}
						\clip (-2.5,-2.8) rectangle (2.5,2.8); 
						\draw[samples=200,domain=-2.1:2.1,smooth,variable=\x]  plot (\x,{1*((\x)^3)-3*(\x)});
						\draw[color=blue,samples=200,domain=-2.3:2.3,smooth,variable=\x] plot  (\x,{1*(\x)+0});
					\end{scope}
					\fill 
						(0,0) circle  (1pt) node[shift=(150:0.4)]{$O$}
						(-{sqrt (3)},0)  circle (1pt) node[shift=(-45:0.4)]{$-\sqrt{3}$}
						({sqrt (3)},0) circle (1pt) node[shift=(135:0.4)]{$\sqrt{3}$} 
						(-2,0) circle (1pt) node[shift=(90:0.4)]{$-2$} 
						 (2,0) circle (1pt) node[shift=(-90:0.4)]{$2$}
						(2,2)  circle (1pt)
						(-2,-2)  circle (1pt);
					\draw[dashed]  (-2,-2)--(-2,0)
						(2,2)--(2,0); 
				\end{tikzpicture}
		\end{center}
		Phương trình hoành độ  giao điểm của $(C)$ và $d\colon x^3-3x=x\Leftrightarrow x^3-4x=0\Leftrightarrow\hoac{&x=0\\ &x=2\\ &x=-2.}$\\
		Gọi $(C')$ là  đồ thị đối xứng với đồ thị $(C)$ qua trục hoành.\\
		Suy ra $(C')\colon  y=-x^3+3x$.\\
		Phương trình hoành độ giao điểm của $(C')$ và $d$  là\\
		$$-x^3+3x=x\Leftrightarrow x^3-2x=0\Leftrightarrow \hoac{&x=0\\ &x=\sqrt{2}\\ &x=-\sqrt{2}.}$$\\ 
		\begin{center}
			\begin{tikzpicture}[line join=round,line cap=round,>=stealth,font=\footnotesize,scale=1]
					\draw[->] (-2.6,0)--(2.6,0) node[below left]  {$x$};
					 \draw[->] (0,-3.1)--(0,3.1) node[below left] {$y$};
					\begin{scope}
						\clip (-2.5,-2.8)  rectangle (2.5,2.8);
						\draw[samples=200,domain=-2.1:2.1,smooth,variable=\x]  plot (\x,{1*((\x)^3)-3*(\x)});
						\draw[color=red,samples=200,domain=-2.1:2.1,smooth,variable=\x]  plot (\x,{-1*((\x)^3)+3*(\x)});
						\draw[color=blue,samples=200,domain=-2.3:2.3,smooth,variable=\x] plot (\x,{1*(\x)+0}); 
					\end{scope}
					 \fill
						(0,0)  circle (1pt) node[shift=(150:0.4)]{$O$}
						(-{sqrt  (3)},0) circle (1pt) node[shift=(40:0.5)]{$-\sqrt{3}$}
						({sqrt (3)},0) circle  (1pt) node[shift=(-135:0.5)]{$\sqrt{3}$}
						(-2,0)  circle (1pt) node[shift=(100:0.25)]{$-2$}
						(2,0) circle  (1pt) node[shift=(-90:0.25)]{$2$}
						(2,2)  circle (1pt)
						(-2,-2)  circle (1pt)
						(-{sqrt (2)},0)  circle (1pt) node[shift=(-45:0.5)]{$-\sqrt{2}$}
						({sqrt (2)},0)  circle (1pt) node[shift=(135:0.5)]{$\sqrt{2}$};
					\draw[dashed]  (-2,-2)--(-2,0)
						 (2,2)--(2,0)
						({sqrt (2)},{sqrt (2)})  -- ({sqrt (2)},0)
						(-{sqrt (2)},-{sqrt (2)}) -- (-{sqrt  (2)},0);
				\end{tikzpicture}
		\end{center}
		Phương trình hoành độ giao  điểm của $(C)$ và trục hoành là
		 $$x^3-3x=0\Leftrightarrow\hoac{&x=0\\ &x=\sqrt{3}\\ &x=-\sqrt{3}.}$$\\
		\begin{center}
			\begin{tikzpicture}[line join=round,line cap=round,>=stealth,font=\footnotesize,scale=1]
					\draw[->] (-2.6,0)--(2.6,0)  node[below left] {$x$};
					\draw[->] (0,-3.1)--(0,3.1) node[below left] {$y$}; 
					\begin{scope}
						\clip (-2.5,-2.8) rectangle (2.5,2.8); 
						\draw[samples=200,domain=-2.1:2.1,smooth,variable=\x] plot (\x,{1*((\x)^3)-3*(\x)}); 
						\draw[color=red,samples=200,domain=-2.1:2.1,smooth,variable=\x]  plot (\x,{-1*((\x)^3)+3*(\x)});
						\draw[color=blue,samples=200,domain=-2.3:2.3,smooth,variable=\x]  plot (\x,{1*(\x)+0});
					\end{scope}
					\fill 
						(0,0)  circle (1pt) node[shift=(150:0.4)]{$O$};
					 \fill[gray,opacity=0.5]
						(0,0) 
						-- plot[samples=200,domain=0:{sqrt (2)},smooth]  (\x,{-\x^3 + 3*\x})
						-- ({sqrt (3)},0) 
						 plot[samples=200,domain={sqrt (2)}:2,smooth] (\x,{1*(\x)+0})
						 -- ({sqrt (3)},0);
					\draw[dashed]  (-2,-2)--(-2,0)
						(2,2)--(2,0) 
						({sqrt (2)},{sqrt (2)}) -- ({sqrt (2)},0) 
						(-{sqrt (2)},-{sqrt (2)}) -- (-{sqrt (2)},0); 
				\end{tikzpicture}
		\end{center}
		Từ đồ thị ta  có
		$$V=2\pi\left(\displaystyle\int_{0}^{\sqrt{2}}\left( x^3-3x\right) ^2 \mathrm{\,d}x+\displaystyle\int_{\sqrt{2}}^{\sqrt{3}}x^2 \mathrm{\,d}x+\displaystyle\int_{\sqrt{3}}^{2}\left[x^2-(x^3-3x)^2\right]\mathrm{\,d}x\right)\approx 29{,}6.$$ 
	 }
\end{bt}


%%%%%%%%%%%%%%%% TỰ LUẬN %%%%%%%%%%%%%%%%

% nguồn: Phan2_Lop_12_Dai_so_Chuong_4
\begin{bt}%[2D4H3-3]%[ma:67b9d893]
	Tính thể tích của khối tròn xoay tạo thành khi quay hình phẳng giới  hạn bởi các đường sau xung quanh trục $Ox$:
	\begin{enumerate}
		\item $y=1-x^2$, $ y=0$, $ x=-1$, $ x=1$; 
		\item $y=\sqrt{25-x^2}$, $  y=0$, $ x=2$, $ x=4$.
	\end{enumerate}
	 \loigiai{
		\begin{enumerate}
			 \item Ta có $V=\pi \displaystyle \int_{-1}^1 \left(1-x^2\right)^2 \mathrm{\,d}x=\dfrac{16 \pi}{15}$.
			\item Ta có $V=\pi \displaystyle \int_2^4 \left(\sqrt{25-x^2}\right)^2 \mathrm{\,d}x=\pi \displaystyle \int_2^4  \left(25-x^2\right) \mathrm{\,d}x=\dfrac{94\pi}{3}$.
		\end{enumerate}
	} 
\end{bt}

% nguồn: Phan2_Lop_12_Dai_so_Chuong_4
\begin{bt}%[2D4H3-3]%[ma:7b3b2b15]
	Tính thể  tích khối cầu có bán kính $R$.
	\loigiai{ 
	\immini{ 
		Khối cầu có bán kính $R$ là khối tròn xoay nhận được khi quay nửa hình tròn giới hạn bởi đồ  thị hàm số $y=\sqrt{R^2-x^2} (-R\le x\le R)$ và trục $Ox$ quanh trục $Ox $ .\\
		Từ đó, thể tích khối cầu là $$V=\pi\displaystyle\int_{-R}^{R} \left(R^2-x^2\right)\mathrm{\,d}x=\pi\left(R^2x-\dfrac{x^3}{3}\right)\bigg|_{-R}^{R}=\dfrac{4\pi  R^3}{3}.$$
		}{ 
		\begin{tikzpicture}[scale=0.7]
			\filldraw[yellow] (3,0) arc [start angle=0,end angle=180,radius=3];
			\draw (0,0) circle (1pt) node[below  left]{$O$};
			\draw[->] (-4,0) -- (4.6,0) node[below]{$x$}; 
			\draw[->] (0,-4) -- (0,4) node[left]{$y$}; 
		\draw (3,0) arc [start angle =0,end angle =180,radius=3];
		\draw[dashed] (-3,0) arc [start angle =180,end angle =360,radius=3];
		\draw[dashed] (0,0) ellipse [x  radius =0.7,y radius=3];
		\fill[pattern=north  east lines,pattern color=black] (2,0) ellipse [x radius=0.45,y radius=2.2];
		\draw[dashed] (2,0) ellipse [x radius=0.45,y radius=2.2]; 
		\draw (2,0) circle (1pt) node[below]{$x$}--(2,2.2) node[above right]{$y=\sqrt{R^2-x^2}$}; 
		\draw[->,>=Stealth] ($(3.9,0.3)+(0:0)$) node[below]{} arc  (90:370:{0.1} and {0.5});
		\path  (-3,0) node[below left]{$-R$}
			 (3,0) node[below right]{$R$};
		\end{tikzpicture}
		 }
	 }
\end{bt}

% nguồn: Phan2_Lop_12_Dai_so_Chuong_4
\begin{bt}%[2D4H3-3]%[ma:7e76d1d8]
	\immini{ 
	Sử dụng tích phân, tính thể tích khối  nón có
	bán kính đáy  $r$ và chiều cao $h$.}
	{ 
	\begin{tikzpicture}
			\fill[yellow] (0,0) ellipse [x radius=2,y  radius=0.5];
		\fill[yellow!20] (0,5)--(2,0) arc [x radius=2,y radius=0.5,start angle=0,end angle=180];
	\draw[dashed](2,0) arc [x radius=2,y radius=0.5,start angle=0,end angle=180];
	\draw(2,0) arc [x radius=2,y radius=0.5,start angle=0,end angle=-180];
	\draw  (-2,0)--(0,5)--(2,0);
	\draw[dashed](0,5)--(0,0)--(2,0); 
	\path (0.7,0) node[above]{$r$} 
		(0,2.5)  node[left]{$h$};
	\end{tikzpicture}
	} 
	\loigiai{ 
		Đặt hệ trục $Oxy$  sao cho đỉnh khối nón trùng với gốc toạ độ $O$.\\
		Đường thẳng đi qua  gốc toạ độ $O$ có dạng $y=ax$.
		Đồ thị  hàm số $y=ax$ đi qua điểm $(h;r)$ nên ta có:
		$r=ah\Rightarrow  a=\dfrac{r}{h}$.
		 Vậy $y=\dfrac{r}{h}x$.\\
		Gọi $D$ là hình phẳng giới hạn bởi đồ thị hàm số $y=\dfrac{r}{h}x$, trục hoành  và hai đường thẳng $x=0, x=h$.
		Thể tích khối nón là thể tích khối  tròn xoay tạo thành khi quay $D$ quanh trục $Ox$.\\
		Thể tích cần tìm là $V=\pi\displaystyle\int_{0}^{h}  \left(\dfrac{r}{h}x\right)^2\mathrm{\,d}x=\pi\displaystyle\int_{0}^{h} \dfrac{r^2}{h^2}x^2\mathrm{\,d}x=\pi\left(\dfrac{r^2x^3}{3h^2}\right)\bigg|_{0}^{h}=\dfrac{r^2h^3\pi}{3h^2}=\dfrac{1}{3}\pi r^2h$.
	 }
\end{bt}

% nguồn: Phan2_Lop_12_Dai_so_Chuong_4
\begin{bt}%[2D4H3-3]%[ma:9b3d4f6e]
	Hãy sử dụng tích phân để chứng minh công thức tính thể tích khối trụ  có chiều cao $h$ và bán kính đáy bằng $r$ (khối trụ là phần không gian được giới hạn bởi một hình trụ kể cả hình trụ đó).
	\loigiai{ 
		Chọn trục $Ox$ trùng với trục của khối trụ, hai đáy nằm trong hai mặt phẳng vuông  góc với $Ox$ tại $x=0$ và $x=h$.
		\begin{center}
			\begin{tikzpicture}[>=stealth,font=\footnotesize,line join=round,line cap=round,scale=0.8]
					 \path
						(0,0)  coordinate (O)
						(2,0)  coordinate (A)
						(-2,0) coordinate  (B)
						 (2,3) coordinate (A')
						 (-2,3) coordinate (B')
						(0,3) coordinate (h) 
						(0,1.5) coordinate (x) 
						 (0,4) coordinate (x');
					\fill[left  color=gray,right color=cyan]
						(A) arc (0:-180: 2  and 0.5)--(B)--(B') arc (180:0: 2 and 0.5)--(A')--cycle;
					\draw 
						(A) arc (0:-180: 2  and 0.5)
						(h) ellipse (2 and 0.5) 
						 (A)--(A')
						 (B)--(B')
						(A')--(h); 
					\draw[dashed] 
						(A) arc (0:180: 2 and  0.5)
						(A)--(B) 
						 (O)--(h)
						(x)--(2,1.5); 
					\draw[->](h)--(x'); 
					\node[below] at (O){$O$}; 
					\node[left] at  (x){$x$};
					\node[left] at (h){$h$}; 
					\node[left] at (-2,1.5){$h$}; 
					\node[above=-2pt] at  (-1,0){$r$};
					 \node[left] at (0,4){$x$};
				\end{tikzpicture}
		\end{center}
		Khi cắt khối trụ bởi mặt phẳng vuông góc với trục $O x$ tại điểm $x\,(0 \leq x  \leq h)$ thì phần chung giữa mặt phẳng và khối trụ là một hình tròn bán kính $r$ có diện tích $S(x)=\pi r^2$.\\
		Thể tích khối trụ  là
		 $$V=\displaystyle\int_0^hS(x)\mathrm{\,d}x=\displaystyle\int_0^h\pi r^2\mathrm{\,d} x=\left(\pi r^2x\right)\Big|_0^h=\pi r^2h\quad (\text{đvtt}).$$
	 }
\end{bt}
